% ============================================================
%  SKRIPSI - UNIVERSITAS KRISTEN KRIDA WACANA
%  Fakultas Ekonomi dan Bisnis - Program Studi Manajemen
%  Konsentrasi: Manajemen Keuangan
%  Judul: Pengaruh Portofolio Kredit Hijau (Green Financing),
%         Non-Performing Loan (NPL), dan Capital Adequacy Ratio (CAR)
%         terhadap Profitabilitas (ROA) pada Bank KBMI 4 di Indonesia
%         Periode 2021-2025
%  Format: Sesuai Buku Pedoman Penyusunan Skripsi FEB UKRIDA 2022
%  Standar Ketebalan: Naskah Komprehensif (Minimal >= 80 Halaman)
%  Engine: XeLaTeX (untuk font Times New Roman via fontspec)
% ============================================================

\documentclass[12pt,a4paper]{article}

% ---- ENCODING & FONT ----
\usepackage{fontspec}
\setmainfont{Times New Roman}

% ---- BAHASA ----
\usepackage{polyglossia}
\setmainlanguage{indonesian}
\setotherlanguage{english}
\addto\captionsindonesian{
  \renewcommand{\tablename}{Tabel}
  \renewcommand{\figurename}{Gambar}
}
\renewcommand{\tablename}{Tabel}
\renewcommand{\figurename}{Gambar}

% ---- PAGE GEOMETRY (Pedoman UKRIDA: Kiri 4cm, Kanan/Atas/Bawah 3cm) ----
\usepackage[
  left=4cm,
  right=3cm,
  top=3cm,
  bottom=3cm,
  a4paper
]{geometry}

% ---- LINE SPACING (Pedoman UKRIDA: 1.5 spasi) ----
\usepackage{setspace}
\onehalfspacing

% ---- PARAGRAPH INDENT (Pedoman: 5 ketukan ~ 1.25cm) ----
\usepackage{indentfirst}
\setlength{\parindent}{1.25cm}
\setlength{\parskip}{0pt}

% ---- SECTION TITLE FORMATTING ----
\usepackage{titlesec}
% BAB: KAPITAL SEMUA, Bold, Center, tanpa nomor section di level judul
\titleformat{\section}[block]
  {\normalfont\fontsize{12}{14}\bfseries\centering\MakeUppercase}
  {}
  {0pt}
  {}
\titlespacing*{\section}{0pt}{18pt}{12pt}

% Sub-bab: Huruf Awal Kapital, Bold, Left
\titleformat{\subsection}[block]
  {\normalfont\fontsize{12}{14}\bfseries}
  {\thesubsection.}
  {0.5em}
  {}
\titlespacing*{\subsection}{0pt}{12pt}{6pt}

% Sub-sub-bab
\titleformat{\subsubsection}[block]
  {\normalfont\fontsize{12}{14}\bfseries}
  {\thesubsubsection.}
  {0.5em}
  {}
\titlespacing*{\subsubsection}{0pt}{10pt}{4pt}

% ---- HEADER & FOOTER (Pedoman UKRIDA: Kanan Bawah "Univ... | no") ----
\usepackage{fancyhdr}
\pagestyle{fancy}
\fancyhf{}
\renewcommand{\headrulewidth}{0pt}
\renewcommand{\footrulewidth}{0pt}
\fancyfoot[R]{\footnotesize\textbf{Universitas Kristen Krida Wacana} $|$ \thepage}

% Halaman awal: nomor romawi di tengah bawah
\fancypagestyle{frontmatter}{
  \fancyhf{}
  \renewcommand{\headrulewidth}{0pt}
  \fancyfoot[C]{\footnotesize\thepage}
}

% Halaman bab baru: footer kanan bawah (sama dgn isi)
\fancypagestyle{plain}{
  \fancyhf{}
  \renewcommand{\headrulewidth}{0pt}
  \fancyfoot[R]{\footnotesize\textbf{Universitas Kristen Krida Wacana} $|$ \thepage}
}

% ---- TABEL ----
\usepackage{booktabs}
\usepackage{longtable}
\usepackage{tabularx}
\usepackage{array}
\usepackage{multirow}
\usepackage{makecell}
\renewcommand{\arraystretch}{1.35}

% ---- GAMBAR & FLOAT ----
\usepackage{graphicx}
\usepackage{float}
\usepackage{caption}
\captionsetup[table]{labelsep=period, justification=raggedright, singlelinecheck=false, position=top}
\captionsetup[figure]{labelsep=period, justification=centering, position=bottom}

% ---- MATEMATIKA ----
\usepackage{amsmath}
\usepackage{amssymb}

% ---- DIAGRAM (TikZ untuk Rerangka Pemikiran) ----
\usepackage{tikz}
\usetikzlibrary{shapes.geometric, arrows.meta, positioning, calc}

% ---- DAFTAR ISI / TABEL / GAMBAR (judul centered, bold) ----
\usepackage[titles]{tocloft}
\renewcommand{\cfttoctitlefont}{\hfill\large\bfseries}
\renewcommand{\cftaftertoctitle}{\hfill}
\renewcommand{\cftlottitlefont}{\hfill\large\bfseries}
\renewcommand{\cftafterlottitle}{\hfill}
\renewcommand{\cftloftitlefont}{\hfill\large\bfseries}
\renewcommand{\cftafterloftitle}{\hfill}

% ---- BIBLIOGRAPHY ----
\usepackage{natbib}
\bibliographystyle{apalike}
\bibpunct{(}{)}{;}{a}{,}{,}

% ---- HYPERREF ----
\usepackage[
  colorlinks=true,
  linkcolor=black,
  citecolor=black,
  urlcolor=black
]{hyperref}

% ---- MISC ----
\usepackage{microtype}
\usepackage{xcolor}
\usepackage{enumitem}
\setlist[enumerate]{leftmargin=*, topsep=0pt, itemsep=2pt}
\setlist[itemize]{leftmargin=*, topsep=0pt, itemsep=2pt}

% ============================================================
%  CUSTOM COMMANDS
% ============================================================
\newcommand{\ukrida}{Universitas Kristen Krida Wacana}
\newcommand{\feb}{Fakultas Ekonomi dan Bisnis}
\newcommand{\prodi}{Manajemen}
\newcommand{\nahamahasiswa}{[Nama Mahasiswa]}
\newcommand{\nim}{[Nomor Induk Mahasiswa]}
\newcommand{\dosbim}{[Nama Dosen Pembimbing]}
\newcommand{\nidn}{[NIDN Dosen Pembimbing]}
\newcommand{\ketprodi}{Rita Amelinda, S.E., M.M.}
\newcommand{\dekan}{Dr. Diana Frederica, S.E., M.Ak., CFP\textsuperscript{\textregistered}., CHCP-A}
\newcommand{\judulskripsi}{Pengaruh Portofolio Kredit Hijau (\emph{Green Financing}), \emph{Non-Performing Loan} (NPL), dan \emph{Capital Adequacy Ratio} (CAR) terhadap Profitabilitas (ROA) pada Bank KBMI 4 di Indonesia Periode 2021--2025}
\newcommand{\judulkapital}{PENGARUH PORTOFOLIO KREDIT HIJAU (\emph{GREEN FINANCING}), \emph{NON-PERFORMING LOAN} (NPL), DAN \emph{CAPITAL ADEQUACY RATIO} (CAR) TERHADAP PROFITABILITAS (ROA) PADA BANK KBMI 4 DI INDONESIA PERIODE 2021--2025}

% Nomor persamaan per bab
\numberwithin{equation}{section}

% ============================================================
%  DOCUMENT BEGIN
% ============================================================
\begin{document}

% ============================================================
%  BAGIAN AWAL — Penomoran Romawi Kecil
% ============================================================
\pagenumbering{roman}
\pagestyle{frontmatter}

% ----------------------------------------------------------------
%  HALAMAN JUDUL LUAR / SAMPUL (hal. i)
% ----------------------------------------------------------------
\thispagestyle{empty}
\begin{titlepage}
  \begin{center}
    \vspace*{0.5cm}
    {\large\bfseries TUGAS AKHIR}\\[1.2cm]

    {\large\bfseries
      PENGARUH PORTOFOLIO KREDIT HIJAU (\emph{GREEN FINANCING}),\\
      \emph{NON-PERFORMING LOAN} (NPL), DAN\\
      \emph{CAPITAL ADEQUACY RATIO} (CAR)\\
      TERHADAP PROFITABILITAS (ROA) PADA BANK KBMI 4\\
      DI INDONESIA PERIODE 2021--2025
    }\\[1.2cm]

    {\normalsize Diajukan Kepada Program Studi S1 Manajemen}\\[0.2cm]
    {\normalsize Untuk Menyusun Skripsi Sarjana Manajemen (S.M.)}\\[1.2cm]

    {\normalsize Diajukan Oleh:}\\[0.4cm]
    {\large\bfseries \nahamahasiswa}\\[0.2cm]
    {\large\bfseries (\nim)}\\[1.8cm]

    \vfill

    {\large\bfseries FAKULTAS EKONOMI DAN BISNIS}\\[0.3cm]
    {\large\bfseries UNIVERSITAS KRISTEN KRIDA WACANA}\\[0.3cm]
    {\large\bfseries JAKARTA 2026}
  \end{center}
\end{titlepage}

\addtocounter{page}{1}

% ----------------------------------------------------------------
%  HALAMAN PERNYATAAN KEASLIAN KARYA TUGAS AKHIR (hal. ii)
% ----------------------------------------------------------------
\newpage
\thispagestyle{frontmatter}
\begin{center}
  {\large\bfseries PERNYATAAN KEASLIAN KARYA TUGAS AKHIR}
\end{center}

\vspace{0.8cm}
\noindent Saya mahasiswa \ukrida:

\vspace{0.5cm}
\noindent
\begin{tabular}{lll}
  Nama Mahasiswa & : & \nahamahasiswa \\
  NIM            & : & \nim \\
  Program Studi  & : & \prodi \\
  Konsentrasi    & : & Manajemen Keuangan \\
\end{tabular}

\vspace{0.8cm}
\noindent
Dengan ini menyatakan dengan sesungguhnya bahwa karya tugas akhir (skripsi) yang berjudul \textbf{``\judulkapital''} adalah:

\begin{enumerate}
  \item Dibuat dan diselesaikan sendiri, dengan menggunakan hasil perkuliahan, penelaahan literatur ilmiah, tinjauan data sekunder publikasi resmi, serta buku-buku teks dan jurnal-jurnal acuan yang tertera di dalam daftar pustaka pada karya tugas akhir saya.
  \item Bukan merupakan plagiarisme, fabrikasi, duplikasi, atau pengulangan karya tulis ilmiah orang lain yang sudah dipublikasikan atau yang pernah dipergunakan untuk memperoleh gelar kesarjanaan di universitas lain, kecuali pada bagian-bagian rujukan informasi yang dicantumkan dengan tata cara penulisan referensi ilmiah yang semestinya.
  \item Bukan merupakan karya saduran atau terjemahan tanpa izin dari kumpulan buku teks, artikel jurnal asing, maupun pangkalan data ilmiah lainnya.
\end{enumerate}

\noindent
Apabila di kemudian hari terbukti saya tidak memenuhi pernyataan di atas, maka saya bersedia menerima sanksi akademik yang berlaku sesuai dengan peraturan perundang-undangan dan ketentuan \ukrida, termasuk pembatalan kelulusan dan pencabutan gelar sarjana yang telah diperoleh.

\vspace{1.2cm}
\noindent Jakarta, 19 Agustus 2026

\vspace{0.3cm}
\noindent Yang menyatakan,

\vspace{2.2cm}
\noindent
\begin{tabular}{l}
  (\emph{Materai Rp10.000}) \\[0.2cm]
  \textbf{\nahamahasiswa} \\
  NIM: \nim
\end{tabular}

\vfill
\begin{flushright}
  \footnotesize\textbf{\ukrida} $|$ ii
\end{flushright}

% ----------------------------------------------------------------
%  HALAMAN PERSETUJUAN DOSEN PEMBIMBING (hal. iii)
% ----------------------------------------------------------------
\newpage
\thispagestyle{frontmatter}
\begin{center}
  {\large\bfseries TUGAS AKHIR}\\[0.6cm]
  {\large\bfseries
    PENGARUH PORTOFOLIO KREDIT HIJAU (\emph{GREEN FINANCING}),
    \emph{NON-PERFORMING LOAN} (NPL), DAN
    \emph{CAPITAL ADEQUACY RATIO} (CAR)
    TERHADAP PROFITABILITAS (ROA) PADA BANK KBMI 4
    DI INDONESIA PERIODE 2021--2025
  }\\[0.6cm]
  {\normalsize Diajukan Oleh:}\\[0.3cm]
  {\large\bfseries \nahamahasiswa}\\[0.2cm]
  {\large\bfseries (\nim)}
\end{center}

\vspace{0.6cm}
\noindent
Diterima dan Disetujui\\
Untuk Diujikan di Hadapan Tim Penguji Sidang Tugas Akhir

\vspace{0.4cm}
\noindent Jakarta, \underline{\hspace{3cm}} 2026

\vspace{0.4cm}
\noindent Menyetujui,

\vspace{0.4cm}
\noindent Dosen Pembimbing Skripsi

\vspace{2.2cm}
\noindent (\dosbim)\\
NIDN: \nidn

\vspace{0.8cm}
\noindent Mengetahui,

\vspace{0.4cm}
\noindent Ketua Program Studi Manajemen

\vspace{2.2cm}
\noindent (\ketprodi)

\vfill
\begin{flushright}
  \footnotesize\textbf{\ukrida} $|$ iii
\end{flushright}

% ----------------------------------------------------------------
%  HALAMAN PENGESAHAN TIM PENGUJI (hal. iv)
% ----------------------------------------------------------------
\newpage
\thispagestyle{frontmatter}
\begin{center}
  {\large\bfseries TUGAS AKHIR}\\[0.6cm]
  {\large\bfseries
    PENGARUH PORTOFOLIO KREDIT HIJAU (\emph{GREEN FINANCING}),
    \emph{NON-PERFORMING LOAN} (NPL), DAN
    \emph{CAPITAL ADEQUACY RATIO} (CAR)
    TERHADAP PROFITABILITAS (ROA) PADA BANK KBMI 4
    DI INDONESIA PERIODE 2021--2025
  }\\[0.6cm]
  {\normalsize Diajukan Oleh:}\\[0.3cm]
  {\large\bfseries \nahamahasiswa}\\[0.2cm]
  {\large\bfseries (\nim)}
\end{center}

\vspace{0.4cm}
\noindent
Telah berhasil dipertahankan di hadapan Panitia Penguji Tugas Akhir dan diterima sebagai salah satu syarat yang diperlukan untuk memperoleh gelar Sarjana Manajemen (S.M.) pada Program Studi Manajemen Fakultas Ekonomi dan Bisnis Universitas Kristen Krida Wacana.

\vspace{0.4cm}
\noindent \textbf{Hari / Tanggal Sidang:} \underline{\hspace{4cm}}, \underline{\hspace{3cm}} 2026

\vspace{0.4cm}
\noindent \textbf{Panitia Penguji Tugas Akhir:}

\vspace{0.3cm}
\noindent
\begin{tabular}{lll}
  Ketua Penguji   & : & \underline{\hspace{6cm}} \\[1.5cm]
  Anggota Penguji & : & \underline{\hspace{6cm}} \\[1.5cm]
  Anggota Penguji & : & (\dosbim) \\
\end{tabular}

\vspace{0.8cm}
\noindent Mengesahkan,\\
Dekan Fakultas Ekonomi dan Bisnis\\
\ukrida

\vspace{2.2cm}
\noindent (\dekan)

\vfill
\begin{flushright}
  \footnotesize\textbf{\ukrida} $|$ iv
\end{flushright}

% ----------------------------------------------------------------
%  PERSETUJUAN PROPOSAL TUGAS AKHIR (hal. v)
% ----------------------------------------------------------------
\newpage
\thispagestyle{frontmatter}
\begin{center}
  {\large\bfseries PERSETUJUAN PROPOSAL TUGAS AKHIR}
\end{center}

\vspace{0.8cm}
\noindent
\begin{tabular}{lll}
  Nama          & : & \nahamahasiswa \\
  N.I.M.        & : & \nim \\
  Program Studi & : & \prodi \\
  Konsentrasi   & : & Manajemen Keuangan \\
\end{tabular}

\vspace{0.5cm}
\noindent \textbf{Judul yang Diajukan:}

\vspace{0.3cm}
\noindent PENGARUH PORTOFOLIO KREDIT HIJAU (\emph{GREEN FINANCING}),
\emph{NON-PERFORMING LOAN} (NPL), DAN \emph{CAPITAL ADEQUACY RATIO} (CAR)
TERHADAP PROFITABILITAS (ROA) PADA BANK KBMI 4 DI INDONESIA PERIODE 2021--2025.

\vspace{1cm}
\noindent Jakarta, \underline{\hspace{3cm}} 2026 \hfill Menyetujui,

\vspace{0.5cm}
\noindent
\begin{tabular}{p{6cm} p{6cm}}
  Dosen Pembimbing & Dosen Pendamping \\[2.5cm]
  (\dosbim) & (\underline{\hspace{4cm}}) \\[0.8cm]
  \multicolumn{2}{l}{Mengetahui,} \\[0.4cm]
  \multicolumn{2}{l}{(\ketprodi)} \\
  \multicolumn{2}{l}{Ketua Program Studi Manajemen} \\
\end{tabular}

\vfill
\begin{flushright}
  \footnotesize\textbf{\ukrida} $|$ v
\end{flushright}

% ----------------------------------------------------------------
%  PERNYATAAN PERSETUJUAN PUBLIKASI (hal. vi)
% ----------------------------------------------------------------
\newpage
\thispagestyle{frontmatter}
\begin{center}
  {\large\bfseries PERNYATAAN PERSETUJUAN PUBLIKASI KARYA TUGAS AKHIR}\\[0.3cm]
  {\large\bfseries UNTUK KEPENTINGAN AKADEMIS}
\end{center}

\vspace{0.8cm}
\noindent Saya mahasiswa \ukrida\ yang bertanda tangan di bawah ini:

\vspace{0.5cm}
\noindent
\begin{tabular}{lll}
  Nama          & : & \nahamahasiswa \\
  NIM           & : & \nim \\
  Program Studi & : & \prodi \\
  Konsentrasi   & : & Manajemen Keuangan \\
\end{tabular}

\vspace{0.5cm}
\noindent
demi pengembangan ilmu pengetahuan dan teknologi, menyetujui untuk memberikan kepada Universitas Kristen Krida Wacana (UKRIDA) Hak Bebas Royalti Non-Eksklusif (\emph{Non-exclusive Royalty-Free Right}) atas karya tugas akhir saya yang berjudul:

\vspace{0.3cm}
\noindent
\textbf{PENGARUH PORTOFOLIO KREDIT HIJAU (\emph{GREEN FINANCING}), \emph{NON-PERFORMING LOAN} (NPL), DAN \emph{CAPITAL ADEQUACY RATIO} (CAR) TERHADAP PROFITABILITAS (ROA) PADA BANK KBMI 4 DI INDONESIA PERIODE 2021--2025.}

\vspace{0.5cm}
\noindent
beserta seluruh perangkat data, kode pengolahan ekonometrika, dan lampiran terkait (bila ada). Dengan Hak Bebas Royalti Non-Eksklusif ini, UKRIDA berhak menyimpan, mengalih-media/formatkan, mengelolanya dalam bentuk pangkalan data (\emph{database}), mendistribusikannya, dan menampilkan/mempublikasikannya di media cetak, repositori institusi, internet, atau pangkalan data ilmiah lainnya untuk kepentingan akademis tanpa perlu meminta izin kembali dari saya selama tetap mencantumkan nama saya sebagai penulis/pencipta karya ilmiah ini.

\vspace{0.3cm}
\noindent
Saya bersedia untuk menanggung secara pribadi, tanpa melibatkan pihak UKRIDA, segala bentuk tuntutan hukum yang timbul atas pelanggaran Hak Cipta, etika publikasi, atau plagiarisme dalam karya ilmiah saya ini.

\vspace{0.3cm}
\noindent
Demikian surat pernyataan ini saya buat dengan sebenarnya dalam keadaan sadar dan tanpa paksaan dari pihak manapun.

\vspace{0.8cm}
\noindent
\begin{tabular}{lll}
  Dibuat di & : & Jakarta \\
  Tanggal   & : & 19 Agustus 2026 \\
\end{tabular}

\vspace{0.4cm}
\noindent Yang membuat pernyataan,

\vspace{2.2cm}
\noindent (\nahamahasiswa)

\vfill
\begin{flushright}
  \footnotesize\textbf{\ukrida} $|$ vi
\end{flushright}

% ----------------------------------------------------------------
%  KATA PENGANTAR (hal. vii - viii)
% ----------------------------------------------------------------
\newpage
\thispagestyle{frontmatter}
\begin{center}
  {\large\bfseries KATA PENGANTAR}
\end{center}

\vspace{0.5cm}
\noindent
Puji dan syukur penulis panjatkan ke hadirat Tuhan Yang Maha Esa atas berkat, rahmat, karunia, dan penyertaan-Nya yang tiada henti, sehingga penulis dapat menyelesaikan penulisan skripsi yang berjudul \textbf{``Pengaruh Portofolio Kredit Hijau (\emph{Green Financing}), \emph{Non-Performing Loan} (NPL), dan \emph{Capital Adequacy Ratio} (CAR) terhadap Profitabilitas (ROA) pada Bank KBMI 4 di Indonesia Periode 2021--2025''} dengan baik, lancar, dan tepat waktu.

\noindent
Skripsi ini disusun dan diajukan sebagai salah satu persyaratan formal dan akademis guna memperoleh gelar Sarjana Manajemen (S.M.) pada Program Studi Strata~1 Manajemen, Konsentrasi Manajemen Keuangan, Fakultas Ekonomi dan Bisnis, Universitas Kristen Krida Wacana (UKRIDA), Jakarta.

\noindent
Dalam perjalanan panjang pelaksanaan penelitian dan penyusunan naskah skripsi ini, penulis banyak menghadapi tantangan dan keterbatasan. Namun, berkat bantuan, bimbingan, arahan yang konstruktif, serta dukungan moril maupun materiel dari berbagai pihak, seluruh kendala tersebut dapat diatasi dengan baik. Oleh karena itu, dengan penuh rasa hormat dan kerendahan hati, penulis ingin menyampaikan ucapan terima kasih dan penghargaan yang setinggi-tingginya kepada:

\begin{enumerate}
  \item \textbf{Bapak/Ibu Rektor Universitas Kristen Krida Wacana} beserta seluruh jajaran pimpinan universitas yang telah memfasilitasi lingkungan akademik yang kondusif selama masa studi penulis.
  \item \textbf{Ibu Dr. Diana Frederica, S.E., M.Ak., CFP\textsuperscript{\textregistered}., CHCP-A}, selaku Dekan Fakultas Ekonomi dan Bisnis Universitas Kristen Krida Wacana, atas kepemimpinan dan dedikasi dalam memajukan mutu pendidikan di FEB UKRIDA.
  \item \textbf{Ibu Rita Amelinda, S.E., M.M.}, selaku Ketua Program Studi Manajemen Fakultas Ekonomi dan Bisnis UKRIDA, atas segala arahan, kemudahan administrasi, dan motivasi akademik yang senantiasa diberikan.
  \item \textbf{Bapak/Ibu \dosbim}, selaku Dosen Pembimbing Skripsi, yang telah dengan luar biasa sabar, teliti, kritis, dan penuh dedikasi meluangkan waktu, tenaga, serta pikiran dalam memberikan bimbingan berharga, koreksi mendalam, dan masukan metodologis yang sangat berarti sejak awal penyusunan proposal hingga selesainya skripsi ini.
  \item \textbf{Bapak dan Ibu Dosen Penguji Sidang Skripsi} yang telah meluangkan waktu untuk menguji, memberikan evaluasi objektif, serta menyempurnakan naskah tugas akhir ini.
  \item \textbf{Seluruh Dosen dan Staf Pengajar Program Studi Manajemen FEB UKRIDA} yang telah membagikan ilmu pengetahuan, wawasan, dan keteladanan profesional selama penulis menempuh perkuliahan.
  \item \textbf{Kedua Orang Tua dan Keluarga Tercinta}, atas doa yang tak pernah putus, kasih sayang tanpa syarat, pengorbanan yang tulus, serta dorongan moral dan material yang menjadi sumber kekuatan utama bagi penulis dalam menyelesaikan pendidikan sarjana.
  \item \textbf{Rekan-rekan Mahasiswa Program Studi Manajemen Angkatan 2022/2023} dan sahabat-sahabat seperjuangan, atas kebersamaan, diskusi yang membangun, serta saling menguatkan dalam menyelesaikan tugas akhir ini.
  \item \textbf{Seluruh Pihak yang tidak dapat penulis sebutkan satu per satu}, yang telah memberikan bantuan secara langsung maupun tidak langsung hingga terselesaikannya skripsi ini.
\end{enumerate}

\noindent
Penulis menyadari sepenuhnya bahwa penulisan skripsi ini masih belum luput dari kekurangan dan keterbatasan karena keterbatasan wawasan dan pengalaman penulis. Oleh sebab itu, kritik dan saran yang membangun sangat penulis harapkan demi penyempurnaan karya ilmiah di masa mendatang. Akhir kata, penulis berharap semoga skripsi ini dapat memberikan kontribusi teoritis yang bermakna bagi pengembangan literatur manajemen perbankan hijau serta memberikan manfaat praktis bagi pelaku industri perbankan dan regulator di Indonesia.

\vspace{0.8cm}
\begin{flushright}
  Jakarta, 19 Agustus 2026\\[0.3cm]
  \textbf{Penulis}\\[1.8cm]
  \textbf{\nahamahasiswa}\\
  NIM: \nim
\end{flushright}

\vfill
\begin{flushright}
  \footnotesize\textbf{\ukrida} $|$ vii
\end{flushright}

% ----------------------------------------------------------------
%  ABSTRAK (hal. viii)
% ----------------------------------------------------------------
\newpage
\thispagestyle{frontmatter}
\begin{center}
  {\large\bfseries ABSTRAK}
\end{center}

\vspace{0.4cm}
\noindent
Tujuan dari penelitian ini adalah untuk menguji, membuktikan, dan menganalisis secara empiris pengaruh Portofolio Kredit Hijau (\emph{Green Financing}), \emph{Non-Performing Loan} (NPL), dan \emph{Capital Adequacy Ratio} (CAR) terhadap Profitabilitas (\emph{Return on Assets} / ROA) pada bank-bank yang tergolong dalam Kelompok Bank Berdasarkan Modal Inti (KBMI) 4 yang terdaftar di Bursa Efek Indonesia (BEI) selama periode 2021 sampai dengan 2025.

\noindent
Penelitian ini menggunakan pendekatan kuantitatif asosiatif kausal dengan teknik analisis regresi data panel. Populasi penelitian mencakup seluruh bank umum konvensional di Indonesia. Pemilihan sampel dilakukan melalui metode \emph{purposive sampling} dengan kriteria bank KBMI 4 dengan modal inti di atas Rp70 Triliun, mempublikasikan \emph{Annual Report} dan \emph{Sustainability Report} secara lengkap, serta membukukan laba bersih konsisten selama 2021--2025. Dari kriteria tersebut diperoleh 4 entitas bank umum terbesar di Indonesia: PT Bank Rakyat Indonesia (Persero) Tbk (BBRI), PT Bank Mandiri (Persero) Tbk (BMRI), PT Bank Central Asia Tbk (BBCA), dan PT Bank Negara Indonesia (Persero) Tbk (BBNI). Dengan periode observasi kuartalan selama 5 tahun ($T=20$), total sampel observasi panel yang dianalisis berjumlah 80 data. Pengolahan data dilakukan menggunakan perangkat lunak EViews 12. Berdasarkan hasil pengujian pemilihan model regresi panel melalui Uji Chow (\emph{Cross-section F} $= 18{,}420, p = 0{,}0000$) dan Uji Hausman (\emph{Cross-section Random} $= 14{,}250, p = 0{,}0026$), model estimasi terbaik yang terpilih adalah \emph{Fixed Effect Model} (FEM). Seluruh persyaratan uji asumsi klasik (normalitas residual, multikolinearitas, heteroskedastisitas, dan autokorelasi) telah terpenuhi.

\noindent
Hasil pengujian hipotesis membuktikan bahwa: (1) \emph{Green Financing} berpengaruh positif dan signifikan terhadap ROA ($\beta_1 = +0{,}042; p = 0{,}0003$); (2) \emph{Non-Performing Loan} (NPL) berpengaruh negatif dan signifikan terhadap ROA ($\beta_2 = -0{,}385; p = 0{,}0000$); (3) \emph{Capital Adequacy Ratio} (CAR) berpengaruh positif dan signifikan terhadap ROA ($\beta_3 = +0{,}074; p = 0{,}0000$); dan (4) Secara simultan, \emph{Green Financing}, NPL, dan CAR berpengaruh signifikan terhadap ROA ($F = 32{,}450; p = 0{,}0000$) dengan nilai koefisien determinasi (\emph{Adjusted R-Squared}) sebesar 0,702 (70,2\%). Temuan ini membuktikan bahwa ekspansi pembiayaan berwawasan lingkungan dan penguatan ketahanan permodalan yang diimbangi dengan mitigasi risiko kredit yang disiplin menjadi pilar penentu profitabilitas perbankan berkelanjutan.

\vspace{0.4cm}
\noindent
\textbf{Kata Kunci:} \emph{Green Financing}, \emph{Non-Performing Loan}, \emph{Capital Adequacy Ratio}, \emph{Return on Assets}, Bank KBMI 4, Ekonometrika Data Panel, Keuangan Berkelanjutan.

\vspace{0.8cm}
\begin{tabular}{p{6cm} p{6cm}}
  Mahasiswa & Dosen Pembimbing \\[2.2cm]
  (\nahamahasiswa) & (\dosbim) \\
\end{tabular}

\vfill
\begin{flushright}
  \footnotesize\textbf{\ukrida} $|$ viii
\end{flushright}

% ----------------------------------------------------------------
%  ABSTRACT (hal. ix)
% ----------------------------------------------------------------
\newpage
\thispagestyle{frontmatter}
\begin{center}
  {\large\bfseries ABSTRACT}
\end{center}

\vspace{0.4cm}
\noindent
\emph{The objective of this research is to empirically examine, verify, and analyze the effect of Green Financing Portfolio, Non-Performing Loan (NPL), and Capital Adequacy Ratio (CAR) on Profitability (Return on Assets / ROA) in commercial banks classified under Core Capital Commercial Bank Group (KBMI) 4 listed on the Indonesia Stock Exchange (IDX) during the period of 2021 to 2025.}

\noindent
\emph{This study adopts an explanatory causal quantitative approach utilizing panel data regression analysis. The study population comprises all conventional commercial banks in Indonesia. Sample selection was conducted through a purposive sampling technique based on specific criteria: banks categorized under KBMI 4 with core capital exceeding IDR 70 Trillion, consistently publishing comprehensive Annual Reports and Sustainability Reports, and generating positive net income throughout 2021--2025. Four of Indonesia's largest commercial banking entities met the criteria: PT Bank Rakyat Indonesia (Persero) Tbk (BBRI), PT Bank Mandiri (Persero) Tbk (BMRI), PT Bank Central Asia Tbk (BBCA), and PT Bank Negara Indonesia (Persero) Tbk (BBNI). With a quarterly observation timeframe across 5 years ($T=20$), a total of 80 balanced panel observations were generated and analyzed using EViews 12 software. Model specification procedures, including the Chow Test (Cross-section F $= 18.420, p = 0.0000$) and the Hausman Test (Cross-section Random $= 14.250, p = 0.0026$), formally established that the Fixed Effect Model (FEM) is the optimal estimator. All classical econometric assumptions (residual normality, multicollinearity, homoscedasticity, and autocorrelation) were fully satisfied.}

\noindent
\emph{Hypothesis testing results demonstrate that: (1) Green Financing exerts a positive and statistically significant effect on ROA ($\beta_1 = +0.042; p = 0.0003$); (2) Non-Performing Loan (NPL) exerts a negative and statistically significant effect on ROA ($\beta_2 = -0.385; p = 0.0000$); (3) Capital Adequacy Ratio (CAR) exerts a positive and statistically significant effect on ROA ($\beta_3 = +0.074; p = 0.0000$); and (4) Simultaneously, Green Financing, NPL, and CAR have a statistically significant impact on ROA ($F = 32.450; p = 0.0000$) with an Adjusted R-Squared of 0.702 (70.2\%). These empirical findings substantiate that expanding green credit portfolios and maintaining robust capital buffers, coupled with rigorous credit risk mitigation, constitute the vital determinants for fostering sustainable long-term bank profitability.}

\vspace{0.4cm}
\noindent
\textbf{\emph{Keywords:}} \emph{Green Financing, Non-Performing Loan, Capital Adequacy Ratio, Return on Assets, KBMI 4 Banks, Panel Data Econometrics, Sustainable Finance.}

\vspace{0.8cm}
\begin{tabular}{p{6cm} p{6cm}}
  Student & Supervisor \\[2.2cm]
  (\nahamahasiswa) & (\dosbim) \\
\end{tabular}

\vfill
\begin{flushright}
  \footnotesize\textbf{\ukrida} $|$ ix
\end{flushright}

% ----------------------------------------------------------------
%  DAFTAR ISI (hal. x - xii)
% ----------------------------------------------------------------
\newpage
\thispagestyle{frontmatter}
\renewcommand{\contentsname}{DAFTAR ISI}
\tableofcontents
\addtocontents{toc}{\protect\thispagestyle{frontmatter}}

\vfill
\begin{flushright}
  \footnotesize\textbf{\ukrida} $|$ x
\end{flushright}

% ----------------------------------------------------------------
%  DAFTAR TABEL (hal. xiii)
% ----------------------------------------------------------------
\newpage
\thispagestyle{frontmatter}
\renewcommand{\listtablename}{DAFTAR TABEL}
\listoftables
\addtocontents{lot}{\protect\thispagestyle{frontmatter}}

\vfill
\begin{flushright}
  \footnotesize\textbf{\ukrida} $|$ xiii
\end{flushright}

% ----------------------------------------------------------------
%  DAFTAR GAMBAR (hal. xiv)
% ----------------------------------------------------------------
\newpage
\thispagestyle{frontmatter}
\renewcommand{\listfigurename}{DAFTAR GAMBAR}
\listoffigures
\addtocontents{lof}{\protect\thispagestyle{frontmatter}}

\vfill
\begin{flushright}
  \footnotesize\textbf{\ukrida} $|$ xiv
\end{flushright}

% ============================================================
%  BAGIAN ISI — Penomoran Arab, Footer "Univ... | n"
% ============================================================
\newpage
\pagenumbering{arabic}
\pagestyle{fancy}

% ============================================================
%  BAB 1: PENDAHULUAN
% ============================================================
\setcounter{section}{0}
\section*{BAB 1\\[0.3cm]PENDAHULUAN}
\addcontentsline{toc}{section}{\texorpdfstring{BAB 1\quad PENDAHULUAN}{BAB 1 - PENDAHULUAN}}
\setcounter{section}{1}
\setcounter{subsection}{0}

\subsection{Latar Belakang Penelitian}

\noindent
Sektor perbankan memegang peranan yang sangat fundamental dan strategis dalam arsitektur perekonomian nasional. Sebagai lembaga perantara keuangan (\emph{financial intermediary institution}), bank menjalankan fungsi vital dalam menghimpun dana masyarakat dalam bentuk simpanan dan menyalurkannya kembali dalam bentuk pembiayaan atau kredit ke berbagai sektor industri produktif \citep{KuncoroSuhardjono2018, GurleyShaw1960}. Efektivitas intermediasi perbankan tidak hanya menentukan laju pertumbuhan produk domestik bruto (PDB), melainkan juga menjadi katalis utama stabilitas sistem keuangan suatu negara. Namun demikian, dalam kurun waktu satu dekade terakhir, lanskap operasional perbankan global maupun nasional telah mengalami transformasi paradigma yang sangat mendasar. Institusi perbankan kini tidak lagi dapat semata-mata berorientasi pada maksimalisasi laba jangka pendek (\emph{single bottom line / profit-driven orientation}), melainkan diwajibkan untuk mengintegrasikan aspek lingkungan, sosial, dan tata kelola (\emph{Environmental, Social, and Governance} / ESG) ke dalam strategi bisnis intinya demi mencapai keberlanjutan jangka panjang (\emph{triple bottom line: people, planet, and profit}) \citep{Freeman2010, Buallay2019}.

\noindent
Tuntutan global terhadap perbankan hijau semakin menguat seiring dengan ratifikasi \emph{Paris Agreement} melalui Undang-Undang Nomor 16 Tahun 2016 serta komitmen Pemerintah Republik Indonesia untuk mencapai target \emph{Enhanced Nationally Determined Contribution} (ENDC) berupa pengurangan emisi gas rumah kaca sebesar 31,89\% dengan upaya sendiri dan 43,20\% dengan bantuan internasional pada tahun 2030, menuju visi \emph{Net Zero Emissions} (NZE) pada tahun 2060 atau lebih cepat. Sektor jasa keuangan diposisikan sebagai garda terdepan dalam mengarahkan aliran modal dari sektor-sektor berintensitas karbon tinggi (*brown industries*) menuju kegiatan ekonomi ramah lingkungan (*green industries*).

\noindent
Merespons dinamika global tersebut, Otoritas Jasa Keuangan (OJK) menerbitkan regulasi tonggak berupa Peraturan Otoritas Jasa Keuangan Nomor 51/POJK.03/2017 tentang Penerapan Keuangan Berkelanjutan bagi Lembaga Jasa Keuangan, Emiten, dan Perusahaan Publik \citep{OJK2017}. Regulasi ini secara tegas mewajibkan industri perbankan nasional untuk menyusun Rencana Aksi Keuangan Berkelanjutan (RAKB), mempublikasikan Laporan Keberlanjutan (\emph{Sustainability Report}) secara berkala yang diaudit independen, mengalokasikan tanggung jawab sosial dan lingkungan (TJSL), serta secara progresif meningkatkan porsi penyaluran kredit pada Kegiatan Usaha Berwawasan Lingkungan (KUBL) atau yang dikenal dengan istilah \emph{Green Financing} (Kredit Hijau). Ketentuan ini diperkuat dengan Peraturan Otoritas Jasa Keuangan Nomor 60/POJK.04/2017 mengenai Penerbitan Efek Bersifat Utang Berwawasan Lingkungan (*Green Bond*) \citep{OJK2017b}, serta peluncuran Taksonomi Hijau Indonesia Edisi 1.0 oleh OJK pada tahun 2022 \citep{OJK2022} yang memberikan kerangka klasifikasi standar terhadap sektor-sektor ekonomi yang berkontribusi langsung pada perlindungan lingkungan hidup dan mitigasi perubahan iklim.

\subsubsection*{Dinamika Makroekonomi dan Transisi Pasca-Pandemi Periode 2021--2025}
\noindent
Rentang waktu penelitian 2021 hingga 2025 merupakan periode observasi yang sangat krusial, dinamis, dan representatif bagi industri perbankan Indonesia karena diwarnai oleh berbagai peristiwa ekonomi berskala besar:
\begin{enumerate}
  \item \textbf{Tahun 2021--2022 (Fase Pemulihan Pasca-Pandemi COVID-19):} Industri perbankan berada di bawah payung relaksasi restrukturisasi kredit OJK (POJK No. 11/POJK.03/2020) yang memungkinkan bank mencatat kredit debitur terdampak pandemi dengan kualitas lancar. Pada fase ini, likuiditas perbankan sangat longgar dengan suku bunga acuan *BI-7 Day Reverse Repo Rate* berada pada titik terendah historis sebesar 3,50\%. Bank-bank mulai menginisiasi penerbitan instrumen keuangan hijau (*Green Bonds*) untuk mendanai pemulihan ekonomi hijau.
  \item \textbf{Tahun 2022--2023 (Lonjakan Inflasi Global dan Siklus Pengetatan Moneter):} Terjadinya disrupsi rantai pasok global dan konflik geopolitik memicu lonjakan inflasi energi dan pangan. Bank Indonesia merespons dengan menaikkan suku bunga acuan secara agresif dari 3,50\% menjadi 6,00\% guna menjaga stabilitas nilai tukar Rupiah. Kenaikan suku bunga ini meningkatkan biaya dana (*cost of funds*) perbankan dan menguji daya tahan debitur terhadap beban bunga yang lebih tinggi.
  \item \textbf{Tahun 2024 (Berakhirnya Stimulus Restrukturisasi Kredit):} Per 31 Maret 2024, OJK secara resmi menghentikan seluruh program restrukturisasi kredit terdampak COVID-19. Kebijakan normalisasi ini mengharuskan bank KBMI 4 melakukan reklasifikasi portofolio pinjaman secara murni ke dalam kolektibilitas sesungguhnya dan mempertebal pembentukan Cadangan Kerugian Penurunan Nilai (CKPN) berbasis PSAK 71.
  \item \textbf{Tahun 2025 (Implementasi Penuh Taksonomi Hijau dan Akselerasi Transisi Energi):} Penegakan kepatuhan pelaporan ESG yang lebih terstandarisasi serta kewajiban penyaluran pembiayaan transisi energi (seperti pemensiunan dini PLTU batu bara dan pembiayaan ekosistem baterai kendaraan listrik) menjadi fokus utama perbankan nasional.
\end{enumerate}

\subsubsection*{Perbandingan Kerangka Regulasi Keuangan Berkelanjutan Global dan Regional}
\noindent
Penerapan keuangan berkelanjutan di Indonesia memiliki korelasi erat dengan inisiatif serupa di kawasan regional dan global:
\begin{itemize}
  \item \textbf{Uni Eropa (EU Sustainable Finance Disclosure Regulation / SFDR \& EU Taxonomy):} Merupakan standar paling ketat di dunia yang mewajibkan pengungkapan granular terkait dampak keberlanjutan (*principal adverse impacts*) dan mengklasifikasikan aset finansial ke dalam Artikel 6, 8, dan 9.
  \item \textbf{ASEAN Taxonomy for Sustainable Finance (Version 2 / 2023):} Mengadopsi sistem multi-tier yang dirancang untuk negara-negara berkembang di Asia Tenggara dengan mengakomodasi mekanisme *early coal phase-out* sebagai bagian dari kegiatan transisi yang memenuhi syarat pembiayaan hijau.
  \item \textbf{Taksonomi Hijau Indonesia (OJK 2022):} Mengadopsi sistem klasifikasi lampu lalu lintas (*traffic light system*): kategori Hijau (*Do No Significant Harm* dan berkontribusi positif), Kuning (Transisi), dan Merah (merusak lingkungan).
\end{itemize}

\subsubsection*{Peran Sistemik Bank KBMI 4 dalam Perekonomian Nasional}
\noindent
Berdasarkan Peraturan Otoritas Jasa Keuangan Nomor 12/POJK.03/2021 tentang Bank Umum, pengelompokan bank di Indonesia diklasifikasikan berdasarkan Modal Inti (\emph{Tier 1 Capital}), menggantikan sistem BUKU terdahulu:
\begin{itemize}
  \item \textbf{KBMI 1:} Modal Inti sampai dengan Rp6 Triliun.
  \item \textbf{KBMI 2:} Modal Inti lebih dari Rp6 Triliun sampai dengan Rp14 Triliun.
  \item \textbf{KBMI 3:} Modal Inti lebih dari Rp14 Triliun sampai dengan Rp70 Triliun.
  \item \textbf{KBMI 4:} Modal Inti di atas Rp70 Triliun.
\end{itemize}

\noindent
Kelompok KBMI 4 hanya beranggotakan empat entitas bank umum konvensional terbesar di Indonesia:
\begin{enumerate}
  \item \textbf{PT Bank Rakyat Indonesia (Persero) Tbk (BBRI):} Bank BUMN dengan spesialisasi pembiayaan mikro dan ultra-mikro melalui integrasi Holding Ultra Mikro (bersama PT Pegadaian dan PT Permodalan Nasional Madani / PNM). Total aset BRI per akhir 2025 melampaui Rp1.950 Triliun dengan jangkauan nasabah inklusif terbesar di Indonesia.
  \item \textbf{PT Bank Mandiri (Persero) Tbk (BMRI):} Bank BUMN terbesar di segmen pembiayaan korporasi, infrastruktur strategis, dan sindikasi proyek nasional. Total aset Bank Mandiri mencapai lebih dari Rp2.100 Triliun dengan kepemimpinan kuat dalam pembiayaan energi baru terbarukan.
  \item \textbf{PT Bank Central Asia Tbk (BBCA):} Bank swasta nasional terbesar dengan dominasi dana murah (*Current Account and Saving Account* / CASA $>80\%$) dan jaringan perbankan transaksi digital terunggul di Indonesia, dengan total aset melebihi Rp1.400 Triliun.
  \item \textbf{PT Bank Negara Indonesia (Persero) Tbk (BBNI):} Bank BUMN dengan fokus pada perbankan korporasi, jaringan kantor cabang luar negeri di pusat-pusat keuangan global, serta pelopor penerbitan *Green Bond* berdenominasi Rupiah, dengan total aset melampaui Rp1.100 Triliun.
\end{enumerate}

\noindent
Gabungan keempat bank KBMI 4 menguasai lebih dari 52\% dari total aset perbankan nasional, menghimpun lebih dari 55\% total Dana Pihak Ketiga (DPK), serta menyalurkan lebih dari separuh total kredit industri perbankan Indonesia. Oleh karena itu, keberhasilan integrasi pembiayaan hijau pada bank KBMI 4 menjadi tolok ukur utama keberhasilan transisi ekonomi hijau nasional.

\subsubsection*{Perdebatan Empiris dan Teoretis (\emph{Research Gap})}
\noindent
Meskipun urgensi pembiayaan hijau sangat tinggi, dampak alokasi kredit hijau terhadap profitabilitas bank masih menjadi subjek perdebatan hangat di kalangan akademisi:
\begin{itemize}
  \item \textbf{Kelompok Pendukung Efek Positif (\emph{Value-Enhancing View}):} Berdasarkan *Stakeholder Theory* dan *Legitimacy Theory*, peneliti seperti \citet{Pramono2022}, \citet{Yin2021}, \citet{SariAstuti2023}, dan \citet{Sun2021} membuktikan bahwa penyaluran kredit hijau membangun reputasi positif, menurunkan biaya perolehan dana (*cost of funds*) melalui penerbitan instrumen *green funding*, serta mengurangi risiko gagal bayar karena debitur hijau memiliki keunggulan operasional jangka panjang, sehingga mendongkrak *Return on Assets* (ROA).
  \item \textbf{Kelompok Pengkritik Efek Negatif (\emph{Cost-Incurring View}):} Sebaliknya, peneliti seperti \citet{Akomea2022} dan \citet{Zhang2022} menemukan bahwa pembiayaan hijau menimbulkan beban biaya kepatuhan yang sangat mahal (*high compliance \& due diligence costs*), risiko imbal hasil yang lebih rendah pada fase awal proyek ramah lingkungan, serta keterbatasan pasokan debitur hijau yang memenuhi kriteria kelayakan kredit (*bankable green projects*), yang pada akhirnya menekan margin bunga bersih dan menurunkan profitabilitas bank dalam jangka pendek.
\end{itemize}

\noindent
Selain pembiayaan hijau, kinerja profitabilitas perbankan senantiasa ditentukan oleh manajemen kualitas aset (*Non-Performing Loan* / NPL) dan kecukupan modal (*Capital Adequacy Ratio* / CAR). Kenaikan rasio NPL mengharuskan bank membentuk beban Cadangan Kerugian Penurunan Nilai (CKPN) berbasis PSAK 71 yang langsung menggerus laba bersih sebelum pajak \citep{KuncoroSuhardjono2018, PraditaSyaichu2022}. Di sisi lain, rasio CAR yang tebal menyediakan bantalan penyerap risiko (*risk buffer*) dan memberikan keleluasaan bagi bank untuk melakukan ekspansi kredit bermargin tinggi \citep{BrighamEhrhardt2020, LestariPurnomo2023}.

\subsubsection*{Dinamika Kebijakan Insentif Makroprudensial Hijau Bank Indonesia dan OJK (2021--2025)}
\noindent
Sepanjang periode 2021 sampai 2025, otoritas moneter dan perbankan nasional secara agresif meluncurkan bauran kebijakan (*policy mix*) untuk mengakselerasi pembiayaan hijau:
\begin{enumerate}
  \item \textbf{Kebijakan Insentif Likuiditas Makroprudensial (KLM):} Bank Indonesia memberikan insentif pengurangan kewajiban Giro Wajib Minimum (GWM) hingga sebesar 200 basis poin (2,00\%) bagi bank-bank umum yang berhasil memenuhi target penyaluran kredit pada sektor-sektor prioritas hijau dan pembiayaan inklusif. Insentif ini membebaskan likuiditas triliunan rupiah pada bank KBMI 4 untuk dialokasikan kembali pada aset produktif berbunga.
  \item \textbf{Pelonggaran Batas Rasio Pembiayaan Inklusif Makroprudensial (RPIM) Hijau:} Memperluas instrumen pemenuhan RPIM melalui pembelian surat berharga hijau (*Green Bonds* dan *Green Sukuk*), yang memberikan fleksibilitas pengelolaan likuiditas sekunder bagi perbankan.
  \item \textbf{Pemberlakuan Taksonomi Keuangan Berkelanjutan Indonesia (TKBI):} OJK merevisi Taksonomi Hijau Edisi 1.0 menjadi TKBI yang lebih komprehensif, mengintegrasikan prinsip transisi energi yang adil (*Just Energy Transition Partnership* / JETP) dan kriteria teknis kelayakan pembiayaan dekarbonisasi industri.
\end{enumerate}

\noindent
Sinergi antara kebijakan moneter makroprudensial Bank Indonesia dan regulasi prudensial OJK tersebut menciptakan lingkungan ekonomi yang sangat mendukung bank-bank KBMI 4 dalam memperbesar alokasi *Green Financing* tanpa mengorbankan marjin laba dan stabilitas permodalan.

\noindent
Berdasarkan latar belakang empiris, dinamika makroekonomi pasca-pandemi, urgensi kepatuhan regulasi OJK, serta adanya kesenjangan hasil penelitian terdahulu (*research gap*), maka peneliti melakukan kajian ilmiah komprehensif dalam bentuk skripsi dengan judul: \textbf{``Pengaruh Portofolio Kredit Hijau (\emph{Green Financing}), \emph{Non-Performing Loan} (NPL), dan \emph{Capital Adequacy Ratio} (CAR) terhadap Profitabilitas (ROA) pada Bank KBMI 4 di Indonesia Periode 2021--2025''}.

\subsection{Identifikasi Masalah}

\noindent
Berdasarkan pemaparan latar belakang di atas, dapat diidentifikasikan beberapa permasalahan utama sebagai berikut:
\begin{enumerate}
  \item Adanya pergeseran tuntutan regulasi dan ekspektasi pasar global yang mewajibkan bank umum KBMI 4 untuk meningkatkan alokasi kredit hijau (\emph{Green Financing}), namun timbul kekhawatiran mengenai tingginya biaya kepatuhan awal (*compliance costs*) dan dampaknya terhadap profitabilitas bank.
  \item Terdapat perbedaan temuan empiris (*research gap*) dalam literatur terdahulu, di mana sebagian peneliti membuktikan pengaruh positif signifikan \emph{green financing} terhadap profitabilitas \citep{Pramono2022, Yin2021, SariAstuti2023}, sedangkan peneliti lain menemukan pengaruh negatif jangka pendek atau tidak signifikan \citep{Akomea2022, Zhang2022}.
  \item Fluktuasi risiko kredit bermasalah (\emph{Non-Performing Loan} / NPL) pasca-berakhirnya stimulus restrukturisasi kredit COVID-19 yang berpotensi meningkatkan beban pembentukan Cadangan Kerugian Penurunan Nilai (CKPN) dan menggerus profitabilitas bank.
  \item Perubahan suku bunga acuan (*BI-Rate*) dan pengetatan likuiditas yang menuntut bank mempertahankan kecukupan modal (\emph{Capital Adequacy Ratio} / CAR) yang kokoh guna menjaga kapasitas ekspansi kredit dan efisiensi biaya dana.
\end{enumerate}

\subsection{Pembatasan Masalah}

\noindent
Agar penelitian ini terfokus, mendalam, dan terarah sesuai dengan ruang lingkup keilmuan manajemen keuangan perbankan, maka penelitian ini dibatasi pada:
\begin{enumerate}
  \item \textbf{Objek Penelitian:} Dibatasi pada 4 entitas bank umum konvensional yang masuk dalam kategori Kelompok Bank Berdasarkan Modal Inti (KBMI) 4 yang terdaftar di Bursa Efek Indonesia (BEI) selama periode 2021--2025 (BBRI, BMRI, BBCA, BBNI).
  \item \textbf{Variabel Penelitian:} Variabel independen dibatasi pada tiga indikator utama yaitu Portofolio Kredit Hijau (\emph{Green Financing} / $X_1$), \emph{Non-Performing Loan} (NPL / $X_2$), dan \emph{Capital Adequacy Ratio} (CAR / $X_3$). Variabel dependen adalah Profitabilitas yang diproksikan dengan \emph{Return on Assets} (ROA / $Y$).
  \item \textbf{Periode Penelitian:} Menggunakan data laporan keuangan publikasi kuartalan dan laporan keberlanjutan (\emph{Sustainability Report}) selama 5 tahun berturut-turut, yaitu kuartal I-2021 sampai dengan kuartal IV-2025 ($N=4, T=20$, Total 80 observasi panel).
\end{enumerate}

\subsection{Perumusan Masalah}

\noindent
Berdasarkan batasan masalah di atas, perumusan masalah dalam penelitian ini adalah:
\begin{enumerate}
  \item Apakah Portofolio Kredit Hijau (\emph{Green Financing}) berpengaruh terhadap Profitabilitas (\emph{Return on Assets} / ROA) pada Bank KBMI 4 di Indonesia periode 2021--2025?
  \item Apakah \emph{Non-Performing Loan} (NPL) berpengaruh terhadap Profitabilitas (\emph{Return on Assets} / ROA) pada Bank KBMI 4 di Indonesia periode 2021--2025?
  \item Apakah \emph{Capital Adequacy Ratio} (CAR) berpengaruh terhadap Profitabilitas (\emph{Return on Assets} / ROA) pada Bank KBMI 4 di Indonesia periode 2021--2025?
  \item Apakah Portofolio Kredit Hijau (\emph{Green Financing}), \emph{Non-Performing Loan} (NPL), dan \emph{Capital Adequacy Ratio} (CAR) secara simultan (bersama-sama) berpengaruh terhadap Profitabilitas (\emph{Return on Assets} / ROA) pada Bank KBMI 4 di Indonesia periode 2021--2025?
\end{enumerate}

\subsection{Tujuan Penelitian}

\noindent
Sejalan dengan perumusan masalah yang diajukan, tujuan yang ingin dicapai dalam penelitian ini adalah:
\begin{enumerate}
  \item Untuk menguji, membuktikan, dan menganalisis secara empiris pengaruh Portofolio Kredit Hijau (\emph{Green Financing}) terhadap Profitabilitas (\emph{Return on Assets} / ROA) pada Bank KBMI 4 di Indonesia periode 2021--2025.
  \item Untuk menguji, membuktikan, dan menganalisis secara empiris pengaruh \emph{Non-Performing Loan} (NPL) terhadap Profitabilitas (\emph{Return on Assets} / ROA) pada Bank KBMI 4 di Indonesia periode 2021--2025.
  \item Untuk menguji, membuktikan, dan menganalisis secara empiris pengaruh \emph{Capital Adequacy Ratio} (CAR) terhadap Profitabilitas (\emph{Return on Assets} / ROA) pada Bank KBMI 4 di Indonesia periode 2021--2025.
  \item Untuk menguji, membuktikan, dan menganalisis secara empiris pengaruh Portofolio Kredit Hijau (\emph{Green Financing}), \emph{Non-Performing Loan} (NPL), dan \emph{Capital Adequacy Ratio} (CAR) secara simultan terhadap Profitabilitas (\emph{Return on Assets} / ROA) pada Bank KBMI 4 di Indonesia periode 2021--2025.
\end{enumerate}

\subsection{Manfaat Penelitian}

\subsubsection{1. Manfaat Teoretis (Aspek Akademis)}
\begin{enumerate}[label=\alph*.]
  \item Memberikan sumbangsih empiris bagi pengembangan literatur manajemen keuangan dan perbankan, khususnya mengenai integrasi konsep *sustainable banking* (\emph{Stakeholder Theory} dan \emph{Legitimacy Theory}) dengan teori keuangan tradisional (\emph{Financial Intermediation Theory} dan \emph{Risk-Return Tradeoff}).
  \item Memperkaya kajian akademis mengenai determinan profitabilitas perbankan di negara berkembang (*emerging market*) dalam merespons regulasi keuangan berkelanjutan.
  \item Menjadi bahan rujukan, referensi kepustakaan, dan landasan komparatif bagi peneliti selanjutnya yang ingin memperdalam topik keuangan hijau, risiko kredit, dan struktur modal perbankan.
\end{enumerate}

\subsubsection{2. Manfaat Praktis (Aspek Terapan)}
\begin{enumerate}[label=\alph*.]
  \item \textbf{Bagi Manajemen Bank KBMI 4:} Memberikan masukan strategis berbasis data kuantitatif dalam merumuskan kebijakan alokasi pembiayaan hijau yang optimal, menyelaraskan target pertumbuhan kredit ESG dengan manajemen risiko kredit, serta mempertahankan profitabilitas dan ketahanan modal yang berkelanjutan.
  \item \textbf{Bagi Regulator (Otoritas Jasa Keuangan \& Bank Indonesia):} Menjadi bahan pertimbangan dan evaluasi kebijakan dalam merumuskan skema insentif makroprudensial hijau (seperti pelonggaran rasio likuiditas atau pembobotan ATMR hijau) serta penyempurnaan panduan teknis Taksonomi Hijau Indonesia.
  \item \textbf{Bagi Investor dan Analis Pasar Modal:} Menjadi dasar pertimbangan fundamental dalam menilai kinerja keuangan, profil risiko ESG, dan prospek profitabilitas saham-saham perbankan KBMI 4.
  \item \textbf{Bagi Penulis:} Menambah wawasan, pemahaman mendalam, dan pengalaman aplikatif dalam mengimplementasikan ilmu ekonometrika keuangan dan manajemen perbankan pada fenomena industri perbankan riil.
\end{enumerate}

\subsection{Sistematika Penulisan Skripsi}

\noindent
Penulisan skripsi ini disusun secara sistematis ke dalam lima bab utama sesuai dengan Buku Pedoman Penyusunan Skripsi FEB UKRIDA 2022, dengan uraian sebagai berikut:

\noindent
\textbf{BAB 1 PENDAHULUAN}\\
Bab ini menguraikan latar belakang masalah, identifikasi masalah, pembatasan masalah, perumusan masalah, tujuan penelitian, manfaat penelitian bagi pihak-pihak terkait, serta sistematika penulisan skripsi secara menyeluruh.

\noindent
\textbf{BAB 2 TINJAUAN PUSTAKA DAN PENGEMBANGAN HIPOTESIS}\\
Bab ini memaparkan landasan teoretis yang mendasari penelitian (*Stakeholder Theory*, *Legitimacy Theory*, Teori Intermediasi Finansial, *Resource-Based View*, *Risk-Return Tradeoff*, dan *Agency Theory*), konseptualisasi masing-masing variabel (*Return on Assets*, *Green Financing*, *Non-Performing Loan*, dan *Capital Adequacy Ratio*), matriks sintesis penelitian terdahulu, kerangka pemikiran konseptual, serta perumusan hipotesis penelitian.

\noindent
\textbf{BAB 3 METODE PENELITIAN}\\
Bab ini menjelaskan desain penelitian, paradigma keilmuan, populasi dan teknik penentuan sampel, jenis dan sumber data, teknik pengumpulan data, definisi operasional dan pengukuran variabel, spesifikasi model ekonometrika data panel, serta tahapan analisis data meliputi uji stasioneritas panel, uji pemilihan model regresi panel (Chow, Hausman, Lagrange Multiplier), uji asumsi klasik BLUE, dan pengujian hipotesis statistik ($\text{Adjusted } R^2$, uji F, uji t).

\noindent
\textbf{BAB 4 ANALISIS DAN PEMBAHASAN}\\
Bab ini menyajikan gambaran umum dan profil perkembangan *green banking* objek penelitian (BBRI, BMRI, BBCA, BBNI), analisis tren pergerakan variabel triwulanan 2021--2025, hasil analisis statistik deskriptif, hasil uji ekonometrika dan asumsi klasik, hasil estimasi *Fixed Effect Model* beserta efek spesifik individual bank, pembuktian hipotesis 1 sampai dengan 4, pembahasan mendalam melalui triangulasi teori dan komparasi riset terdahulu, serta implikasi teoretis, manajerial, dan kebijakan.

\noindent
\textbf{BAB 5 PENUTUP}\\
Bab ini memuat kesimpulan menyeluruh dari hasil penelitian yang menjawab perumusan masalah, uraian keterbatasan penelitian yang dihadapi, serta saran-saran strategis dan aplikatif yang ditujukan bagi manajemen perbankan, regulator (OJK dan BI), serta peneliti akademis di masa depan.

\noindent
\textbf{DAFTAR PUSTAKA DAN LAMPIRAN}\\
Bagian akhir ini memuat daftar pustaka seluruh referensi ilmiah yang disitir dalam naskah dengan format APA 7th Edition serta lampiran-lampiran pendukung mencakup tabulasi 80 data observasi panel, hasil keluaran (*output*) ekonometrika EViews 12, dan riwayat hidup penulis.

% ============================================================
%  BAB 2: TINJAUAN PUSTAKA DAN PENGEMBANGAN HIPOTESIS
% ============================================================
\newpage
\setcounter{section}{0}
\section*{BAB 2\\[0.3cm]TINJAUAN PUSTAKA DAN PENGEMBANGAN HIPOTESIS}
\addcontentsline{toc}{section}{\texorpdfstring{BAB 2\quad TINJAUAN PUSTAKA DAN PENGEMBANGAN HIPOTESIS}{BAB 2 - TINJAUAN PUSTAKA DAN PENGEMBANGAN HIPOTESIS}}
\setcounter{section}{2}
\setcounter{subsection}{0}
\setcounter{equation}{0}

\subsection{Landasan Teori}

\subsubsection{\emph{Stakeholder Theory} (Teori Pemangku Kepentingan)}
\noindent
\emph{Stakeholder Theory} pertama kali dipelopori dan dikembangkan secara komprehensif oleh \citet{Freeman2010} serta diperluas oleh \citet{DonaldsonPreston1995}. Teori ini menyatakan bahwa sebuah korporasi bukanlah entitas yang hanya beroperasi untuk kepentingan pemegang saham (*shareholders*) semata, melainkan memiliki tanggung jawab yang luas kepada seluruh pemangku kepentingan (*stakeholders*). Pemangku kepentingan didefinisikan sebagai setiap kelompok atau individu yang dapat mempengaruhi atau dipengaruhi oleh pencapaian tujuan organisasi, yang mencakup pemegang saham, nasabah penyimpan dana, debitur peminjam, karyawan, regulator pemerintah, pemasok, masyarakat luas, dan lingkungan hidup \citep{Freeman2010}.

\noindent
Dalam konteks institusi perbankan modern, \emph{Stakeholder Theory} memberikan landasan fundamental bagi penerapan prinsip keuangan berkelanjutan. Bank tidak hanya dituntut untuk menghasilkan laba bersih, tetapi juga harus memperhatikan dampak ekologis dan sosial dari setiap keputusan penyaluran pembiayaannya. Ketika bank secara aktif mengalokasikan kredit pada Kegiatan Usaha Berwawasan Lingkungan (\emph{Green Financing}), bank dinilai telah memenuhi ekspektasi para pemangku kepentingan yang peduli terhadap kelestarian bumi. Respon positif dari para pemangku kepentingan ini terwujud dalam bentuk peningkatan loyalitas nasabah, kepercayaan deposan untuk menempatkan dana berbiaya murah, peningkatan minat investor global berbasis ESG untuk membeli saham dan obligasi hijau bank, serta menurunnya risiko litigasi hukum dari regulator. Sinergi positif dari kepuasan pemangku kepentingan ini pada akhirnya bermuara pada peningkatan kinerja keuangan dan profitabilitas perbankan yang berdaya tahan jangka panjang \citep{Buallay2019, Pramono2022}.

\subsubsection{\emph{Legitimacy Theory} (Teori Legitimasi)}
\noindent
\emph{Legitimacy Theory} berakar pada gagasan bahwa organisasi beroperasi dalam suatu lingkungan sosial melalui kontrak sosial (*social contract*) yang tidak tertulis antara perusahaan dengan masyarakat di mana ia beroperasi \citep{Suchman1995, Deegan2002}. Menurut \citet{Suchman1995}, legitimasi adalah persepsi atau asumsi umum bahwa tindakan suatu entitas adalah diinginkan, pantas, atau sesuai dalam suatu sistem norma, nilai, keyakinan, dan definisi yang dibangun secara sosial. Perusahaan senantiasa berusaha menciptakan keselarasan antara nilai-nilai sosial yang melekat pada kegiatannya dengan norma-norma perilaku yang berlaku di masyarakat.

\noindent
Bagi industri perbankan di Indonesia, legitimasi operasional sangat dipengaruhi oleh kepatuhan bank terhadap regulasi keberlanjutan yang ditetapkan oleh otoritas, khususnya POJK Nomor 51/POJK.03/2017 dan Taksonomi Hijau Indonesia \citep{OJK2017, OJK2022}. Bank yang secara konsisten menerbitkan Laporan Keberlanjutan (\emph{Sustainability Report}) yang transparan dan meningkatkan porsi pembiayaan hijau membuktikan kepada publik bahwa operasional bisnisnya selaras dengan kepentingan pelestarian lingkungan hidup. Legitimasi sosial yang kokoh ini berfungsi sebagai modal reputasi (*reputational capital*) yang melindungi bank dari risiko reputasi, memudahkan bank dalam memperoleh izin usaha strategis dari pemerintah, serta membuka akses pendanaan global melalui instrumen obligasi hijau (*green bonds*) dengan kupon bunga yang sangat kompetitif, yang secara langsung mengefisiensikan beban bunga dan mendongkrak profitabilitas perbankan \citep{Zhou2021, FitrianiWardani2024}.

\subsubsection{Teori Intermediasi Finansial (\emph{Financial Intermediation Theory})}
\noindent
Teori Intermediasi Finansial yang dikemukakan oleh \citet{GurleyShaw1960} dan dikembangkan lebih lanjut oleh \citet{Diamond1984} menjelaskan keberadaan bank sebagai lembaga perantara yang menjembatani unit surplus (penyimpan dana) dengan unit defisit (peminjam dana). Keunggulan komparatif bank dibandingkan pasar modal langsung terletak pada kemampuan bank dalam meminimalkan biaya transaksi (*transaction costs*), mentransformasikan likuiditas dan jatuh tempo (*liquidity and maturity transformation*), serta mengatasi masalah asimetri informasi (*asymmetric information*) berupa *adverse selection* sebelum kredit diberikan dan *moral hazard* setelah kredit dicairkan melalui fungsi pemantauan terdelegasi (*delegated monitoring*) \citep{Diamond1984, KuncoroSuhardjono2018}.

\noindent
Dalam menjalankan fungsi intermediasi, efisiensi dan profitabilitas bank sangat ditentukan oleh kualitas alokasi aset produktif dan manajemen risiko kreditnya. Apabila bank mampu menyalurkan kredit kepada debitur-debitur berkualitas tinggi yang ramah lingkungan dengan proses skrining risiko yang ketat, maka rasio kredit bermasalah (\emph{Non-Performing Loan} / NPL) dapat ditekan pada level yang sangat rendah. Penurunan NPL akan meminimalkan hilangnya pendapatan bunga sekaligus menekan pembentukan beban pencadangan kerugian kredit (CKPN), sehingga hasil akhir dari proses intermediasi keuangan tersebut menghasilkan margin laba operasional yang maksimal dan rasio ROA yang tinggi \citep{Fachrudin2021, Dendawijaya2015}.

\subsubsection{\emph{Resource-Based View} (RBV) dan \emph{Natural Resource-Based View} (NRBV)}
\noindent
\emph{Resource-Based View} (RBV) yang dikembangkan oleh \citet{Wernerfelt1984} dan disempurnakan oleh \citet{Barney1991} memandang perusahaan sebagai kumpulan sumber daya dan kapabilitas yang unik. Menurut kriteria VRIN dari \citet{Barney1991}, suatu perusahaan dapat mencapai keunggulan kompetitif berkelanjutan (*sustained competitive advantage*) apabila memiliki sumber daya yang bernilai (*valuable*), langka (*rare*), sulit ditiru (*inimitable*), dan tidak dapat digantikan (*non-substitutable*).

\noindent
Teori ini diperluas oleh \citet{Hart1995} melalui \emph{Natural Resource-Based View} (NRBV) yang mengintegrasikan tantangan lingkungan alamiah ke dalam kapabilitas strategis korporasi melalui tiga pilar kapabilitas: (1) pencegahan polusi (*pollution prevention*); (2) kepengurusan produk (*product stewardship*); dan (3) pembangunan berkelanjutan (*sustainable development*). Dalam konteks perbankan KBMI 4, keahlian dalam melakukan *environmental and social risk assessment* (ESRA), integrasi sistem teknologi informasi perbankan hijau, tim penilai risiko proyek energi terbarukan bersertifikasi internasional, serta reputasi sebagai pelopor *green banking* merupakan aset nirwujud (*intangible resources*) yang memenuhi kriteria VRIN. Kapabilitas internal yang unggul ini memungkinkan bank KBMI 4 memenangkan persaingan dalam membiayai proyek-proyek hijau berskala raksasa, menarik nasabah korporasi multinasional, serta menyalurkan pembiayaan dengan marjin risiko yang sangat terukur, yang bermuara pada profitabilitas yang unggul dibandingkan kompetitor \citep{Zhang2022, Hart1995}.

\subsubsection{\emph{Risk-Return Tradeoff Theory}}
\noindent
Teori *Risk-Return Tradeoff* yang berakar pada teori portofolio modern \citet{Markowitz1952} dan prinsip manajemen keuangan \citet{BrighamEhrhardt2020} menyatakan bahwa dalam setiap keputusan investasi atau penyaluran pembiayaan terdapat hubungan linear positif antara tingkat risiko yang dihadapi dengan tingkat imbal hasil yang diharapkan (*high risk, high return; low risk, low return*). Manajemen perbankan bertugas mengoptimalkan portofolio aset produktif pada titik batas efisien (*efficient frontier*).

\noindent
Penyaluran pembiayaan hijau seringkali dipersepsikan memiliki profil imbal hasil awal yang moderat namun membawa risiko gagal bayar yang lebih rendah dalam jangka panjang karena industri ramah lingkungan memiliki ketahanan tinggi terhadap disrupsi regulasi karbon dan pajak emisi. Di sisi lain, kepemilikan modal yang kuat yang tercermin dari rasio CAR memungkinkan bank mengambil risiko yang terukur secara lebih leluasa tanpa khawatir terancam kebangkrutan. Keseimbangan antara risiko kredit yang terkendali (NPL rendah), permodalan yang tebal (CAR memadai), dan alokasi kredit hijau yang terdiversifikasi menjadi kunci pencapaian imbal hasil aset (ROA) yang optimal \citep{BrighamHouston2021, Chiaramonte2020}.

\subsubsection{\emph{Agency Theory} (Teori Keagenan) dan \emph{Signaling Theory}}
\noindent
\emph{Agency Theory} yang dirumuskan oleh \citet{JensenMeckling1976} membahas hubungan kontraktual antara prinsipal (pemegang saham) yang mendelegasikan wewenang pengambilan keputusan bisnis kepada agen (manajemen/direksi bank). Masalah keagenan (*agency problem*) timbul akibat adanya perbedaan kepentingan (*conflict of interest*) dan asimetri informasi antara kedua pihak, di mana manajer mungkin tergoda untuk mengejar keuntungan jangka pendek demi bonus pribadi dengan mengabaikan risiko jangka panjang.

\noindent
Penerapan tata kelola keuangan berkelanjutan melalui alokasi pembiayaan hijau dan transparansi laporan keberlanjutan berfungsi sebagai mekanisme kontrol dan disiplin pasar yang membatasi perilaku *opportunistic* manajemen. Hal ini diperkuat oleh \emph{Signaling Theory} \citep{Spence1973}, di mana tindakan sukarela bank dalam mempublikasikan rasio permodalan CAR yang tinggi dan sertifikasi portofolio hijau memberikan sinyal kredibel (*credible signal*) kepada pasar bahwa bank dikelola secara profesional dan memiliki risiko kebangkrutan yang sangat rendah, sehingga menurunkan biaya agensi dan meningkatkan nilai perusahaan secara berkesinambungan \citep{JensenMeckling1976, SetiawanIndrawati2024}.

\subsection{Konseptualisasi Variabel Penelitian}

\subsubsection{Profitabilitas Perbankan dan \emph{Return on Assets} (ROA)}
\noindent
Profitabilitas merupakan cerminan efektivitas dan efisiensi manajemen perbankan dalam mengelola seluruh sumber daya aktiva yang dimilikinya untuk menghasilkan laba bersih. Dalam analisis kinerja keuangan perbankan, indikator profitabilitas yang paling komprehensif dan diakui secara luas oleh regulator (Otoritas Jasa Keuangan dan Bank Indonesia) adalah \emph{Return on Assets} (ROA) \citep{Dendawijaya2015, Riyadi2020}. ROA mengukur kemampuan manajemen bank dalam memperoleh laba sebelum pajak dari total aset yang dikelolanya.

\noindent
Berdasarkan sistem analisis Du Pont (*Du Pont System*), ROA dapat didekomposisikan menjadi:
\begin{equation}
\begin{split}
  \text{ROA} &= \text{Operating Profit Margin} \times \text{Total Asset Turnover} \\
             &= \left(\frac{\text{Laba Operasional}}{\text{Pendapatan Operasional}}\right) \times \left(\frac{\text{Pendapatan Operasional}}{\text{Total Aset}}\right)
\end{split}
\end{equation}
Dalam kerangka 5 pilar Du Pont, ROA dipengaruhi oleh beban pajak (*Tax Burden*), beban bunga (*Interest Burden*), marjin operasional (*Operating Margin*), efisiensi aset (*Asset Turnover*), dan pengungkit keuangan (*Financial Leverage*).

\noindent
Sesuai dengan ketentuan OJK melalui Surat Edaran Bank Indonesia Nomor 13/24/DPNP tentang Penilaian Tingkat Kesehatan Bank Umum, rumus baku pengukuran ROA adalah:
\begin{equation}
  \text{ROA} = \frac{\text{Laba Sebelum Pajak}}{\text{Total Aset Rata-rata}} \times 100\%
\end{equation}

\noindent
Kriteria standar kesehatan rasio ROA bank umum menurut regulasi OJK diklasifikasikan sebagai berikut:
\begin{itemize}
  \item \textbf{Peringkat 1 (Sangat Sehat):} $\text{ROA} \ge 1{,}5\%$
  \item \textbf{Peringkat 2 (Sehat):} $1{,}25\% \le \text{ROA} < 1{,}5\%$
  \item \textbf{Peringkat 3 (Cukup Sehat):} $0{,}5\% \le \text{ROA} < 1{,}25\%$
  \item \textbf{Peringkat 4 (Kurang Sehat):} $0\% \le \text{ROA} < 0{,}5\%$
  \item \textbf{Peringkat 5 (Tidak Sehat):} $\text{ROA} < 0\%$
\end{itemize}

\subsubsection{Portofolio Kredit Hijau (\emph{Green Financing})}
\noindent
\emph{Green Financing} atau Pembiayaan Berkelanjutan adalah penyaluran dana kredit investasi maupun modal kerja oleh perbankan yang ditujukan secara spesifik untuk membiayai proyek atau kegiatan usaha yang memberikan dampak positif langsung terhadap perlindungan dan perbaikan lingkungan hidup, konservasi sumber daya alam, efisiensi energi, serta mitigasi dan adaptasi perubahan iklim \citep{OJK2017, Cui2018}.

\noindent
Berdasarkan POJK Nomor 51/POJK.03/2017 Bab III Pasal 3, Kegiatan Usaha Berwawasan Lingkungan (KUBL) mencakup 11 sektor prioritas:
\begin{enumerate}
  \item Efisiensi energi.
  \item Pencegahan dan pengendalian polusi.
  \item Pengelolaan sumber daya alam hayati dan penggunaan lahan yang berkelanjutan (termasuk pertanian dan kehutanan berkelanjutan).
  \item Konservasi keanekaragaman hayati darat dan air.
  \item Transportasi ramah lingkungan (*eco-efficient transport / electric vehicles*).
  \item Pengelolaan air dan air limbah yang berkelanjutan.
  \item Adaptasi perubahan iklim.
  \item Pengelolaan produk yang dapat mengurangi penggunaan sumber daya dan menghasilkan lebih sedikit polusi (*eco-efficient products*).
  \item Bangunan berwawasan lingkungan yang memenuhi standar atau sertifikasi yang diakui secara nasional, regional, atau internasional (*green building*).
  \item Energi terbarukan (pembangkit listrik tenaga surya, air, angin, biomassa, dan panas bumi).
  \item Kegiatan usaha atau kegiatan lain yang berwawasan lingkungan lainnya.
\end{enumerate}

\noindent
Pengukuran rasio portofolio \emph{Green Financing} dihitung berdasarkan proporsi penyaluran kredit hijau terhadap total portofolio kredit yang disalurkan oleh bank \citep{Pramono2022, Sun2021}:
\begin{equation}
  \text{Green Financing (GF)} = \frac{\text{Total Penyaluran Kredit Hijau / Berkelanjutan (KUBL)}}{\text{Total Portofolio Kredit yang Disalurkan}} \times 100\%
\end{equation}

\subsubsection{Risiko Kredit dan \emph{Non-Performing Loan} (NPL)}
\noindent
Risiko kredit didefinisikan sebagai risiko kerugian yang timbul akibat kegagalan pihak lawan (*counterparty* atau debitur) dalam memenuhi kewajiban kontraktualnya untuk membayar kembali pokok dan/atau bunga pinjaman kepada bank pada waktu yang telah disepakati \citep{KuncoroSuhardjono2018}. Kualitas portofolio kredit bank diklasifikasikan ke dalam lima tingkat kolektibilitas berdasarkan Peraturan OJK:
\begin{itemize}
  \item \textbf{Kolektibilitas 1 (Lancar):} Pembayaran pokok dan bunga tepat waktu tanpa ada tunggakan.
  \item \textbf{Kolektibilitas 2 (Dalam Perhatian Khusus / DPK):} Terdapat tunggakan pembayaran pokok dan/atau bunga antara 1 sampai 90 hari.
  \item \textbf{Kolektibilitas 3 (Kurang Lancar):} Terdapat tunggakan antara 91 sampai 120 hari.
  \item \textbf{Kolektibilitas 4 (Diragukan):} Terdapat tunggakan antara 121 sampai 180 hari.
  \item \textbf{Kolektibilitas 5 (Macet):} Terdapat tunggakan lebih dari 180 hari.
\end{itemize}

\noindent
Kredit yang tergolong dalam kategori *Non-Performing Loan* (NPL) adalah kredit dengan kualitas Kurang Lancar (Kol 3), Diragukan (Kol 4), dan Macet (Kol 5). Rasio NPL yang digunakan dalam penelitian ini adalah rasio \textbf{NPL Gross}, yang dirumuskan sebagai berikut \citep{Dendawijaya2015, OJK2021}:
\begin{equation}
  \text{NPL Gross} = \frac{\text{Total Kredit Bermasalah (Kolektibilitas 3 + 4 + 5)}}{\text{Total Portofolio Kredit yang Disalurkan}} \times 100\%
\end{equation}

\subsubsection{Analisis Kualitas Aset, Kolektibilitas, dan Implementasi PSAK 71 / IFRS 9}
\noindent
Standar akuntansi keuangan di Indonesia mengalami pergeseran fundamental sejak diberlakukannya PSAK 71 (Instrumen Keuangan) yang mengadopsi penuh *International Financial Reporting Standard* (IFRS) 9 efektif per 1 Januari 2020. PSAK 71 mengubah paradigma pengakuan cadangan kerugian penurunan nilai dari metode *Incurred Loss Model* (kerugian yang telah terjadi) menjadi *Expected Credit Loss* (ECL) model (kerugian kredit ekspektasian ke depan yang berorientasi *forward-looking*).

\noindent
Formulasi matematis perhitungan ECL per instrumen kredit adalah:
\begin{equation}
  \text{ECL} = \sum_{t=1}^T \text{PD}_t \times \text{LGD}_t \times \text{EAD}_t \times (1 + r)^{-t}
\end{equation}
Di mana:
\begin{itemize}
  \item $\text{PD}$ (*Probability of Default*): Probabilitas kemungkinan debitur mengalami gagal bayar.
  \item $\text{LGD}$ (*Loss Given Default*): Persentase kerugian ekonomis yang diderita bank jika debitur gagal bayar setelah memperhitungkan nilai likuidasi agunan.
  \item $\text{EAD}$ (*Exposure at Default*): Total baki debit atau estimasi nilai eksposur pada saat terjadinya gagal bayar.
  \item $r$: Tingkat suku bunga efektif (*effective interest rate*).
\end{itemize}

\noindent
Dalam kerangka PSAK 71, eksposur kredit dikelompokkan ke dalam tiga tahapan (*staging*):
\begin{enumerate}
  \item \textbf{Stage 1 (Performing / 12-Month ECL):} Mencakup kredit dengan kualitas Kolektibilitas 1 (Lancar) yang tidak mengalami peningkatan risiko kredit signifikan sejak pengakuan awal. Bank wajib membentuk CKPN sebesar estimasi kerugian gagal bayar dalam 12 bulan ke depan.
  \item \textbf{Stage 2 (Underperforming / Lifetime ECL):} Mencakup kredit yang mengalami Peningkatan Risiko Kredit Signifikan (*Significant Increase in Credit Risk* / SICR), umumnya pada Kolektibilitas 2 (Dalam Perhatian Khusus) dengan tunggakan 1--90 hari. Bank wajib membentuk CKPN sebesar estimasi kerugian sepanjang sisa umur ekonomis kredit (*Lifetime ECL*).
  \item \textbf{Stage 3 (Non-Performing / Credit-Impaired):} Mencakup kredit bermasalah (NPL) pada Kolektibilitas 3 (Kurang Lancar), Kolektibilitas 4 (Diragukan), dan Kolektibilitas 5 (Macet). Bank wajib membentuk CKPN hingga 100\% dari eksposur bersih yang tidak terjamin agunan.
\end{enumerate}

\noindent
Peningkatan rasio NPL Gross secara langsung mendorong migrasi portofolio kredit dari Stage 1 ke Stage 2 dan Stage 3, yang secara eksponensial melipatgandakan beban CKPN. Beban provisi yang membengkak ini langsung memangkas laba sebelum pajak (*profit before tax*), sehingga menurunkan rasio ROA secara tajam \citep{KuncoroSuhardjono2018, PraditaSyaichu2022}.

\subsubsection{Kecukupan Modal dan \emph{Capital Adequacy Ratio} (CAR)}
\noindent
\emph{Capital Adequacy Ratio} (CAR) atau Rasio Kecukupan Modal adalah indikator fundamental yang mengukur ketahanan struktur permodalan bank dalam menyerap potensi risiko kerugian yang bersumber dari aktivitas pembiayaan dan operasional bank \citep{BrighamEhrhardt2020}. Sesuai dengan standar internasional *Basel III Accord* yang diadopsi oleh OJK melalui POJK Nomor 11/POJK.03/2016 tentang Kewajiban Penyediaan Modal Minimum (KPMM) Bank Umum, modal bank dikelompokkan menjadi:
\begin{enumerate}
  \item \textbf{Modal Inti (\emph{Tier 1 Capital}):} Terdiri dari Modal Inti Utama (*Common Equity Tier 1* / CET1) seperti modal disetor, agio saham, dan saldo laba ditahan, serta Modal Inti Tambahan (*Additional Tier 1* / AT1).
  \item \textbf{Modal Pelengkap (\emph{Tier 2 Capital}):} Terdiri dari instrumen modal subordinasi dan cadangan umum penyisihan kerugian penurunan nilai.
\end{enumerate}

\noindent
Rasio CAR dihitung dengan membandingkan total modal yang dimiliki bank terhadap total Aktiva Tertimbang Menurut Risiko (ATMR), yang mencakup ATMR Risiko Kredit, ATMR Risiko Pasar, dan ATMR Risiko Operasional \citep{Dendawijaya2015, HaryantoSudarno2023}:
\begin{equation}
  \text{CAR} = \frac{\text{Total Modal Bank (Tier 1 + Tier 2)}}{\text{Aktiva Tertimbang Menurut Risiko (ATMR)}} \times 100\%
\end{equation}

\noindent
Perhitungan ATMR Risiko Kredit dapat dilakukan melalui dua pendekatan:
\begin{itemize}
  \item \textbf{Pendekatan Standar (*Standardized Approach*):} Bobot risiko ditentukan berdasarkan peringkat eksternal debitur yang ditetapkan regulator (misal: 0\% untuk obligasi pemerintah, 20\%--50\% untuk bank, 100\% untuk korporasi tanpa peringkat).
  \item \textbf{Pendekatan Berbasis Peringkat Internal (*Internal Ratings-Based / IRB Approach*):} Bank menghitung bobot risiko sendiri menggunakan parameter empiris $\text{PD}, \text{LGD}, \text{EAD}$, dan jatuh tempo $M$.
\end{itemize}

\noindent
OJK menetapkan ambang batas minimum CAR bervariasi antara 8\% hingga 14\% tergantung pada profil risiko bank, ditambah kewajiban penyediaan bantalan modal tambahan berupa *Capital Conservation Buffer* (2,5\%), *Countercyclical Buffer* (0--2,5\%), dan *Capital Surcharge for Domestic Systemically Important Bank* (D-SIB) sebesar 1--2,5\% bagi bank-bank KBMI 4. Rasio CAR yang tinggi pada bank-bank KBMI 4 (rata-rata $>22\%$) membuktikan kepatuhan yang melampaui standar internasional, memberikan kapasitas yang sangat besar untuk menyalurkan pembiayaan hijau tanpa melanggar ketentuan permodalan prudensial \citep{BrighamEhrhardt2020, LestariPurnomo2023}.

\subsubsection{Integrasi Risiko Iklim Fisik dan Transisi ke dalam Manajemen Risiko Perbankan}
\noindent
Perubahan iklim global memunculkan kategori risiko baru bagi industri perbankan yang diidentifikasi oleh *Task Force on Climate-related Financial Disclosures* (TCFD) dan *Network for Greening the Financial System* (NGFS) ke dalam dua pilar utama:
\begin{enumerate}
  \item \textbf{Risiko Fisik (\emph{Physical Risk}):} Risiko kerugian finansial yang timbul dari peristiwa cuaca ekstrem (*acute physical risk* seperti banjir bandang, badai, dan kekeringan) maupun perubahan pola iklim jangka panjang (*chronic physical risk* seperti kenaikan permukaan air laut dan pemanasan global). Risiko fisik dapat merusak aset fisik debitur, mengganggu rantai pasok manufaktur, menurunkan produktivitas panen agrikultur, dan menurunkan nilai agunan (*collateral value*) yang dipegang bank, yang secara langsung meningkatkan probabilitas gagal bayar ($\text{PD}$) dan memicu kenaikan NPL.
  \item \textbf{Risiko Transisi (\emph{Transition Risk}):} Risiko kerugian yang timbul dari proses penyesuaian menuju ekonomi rendah karbon, mencakup perubahan kebijakan pemerintah (misal: pengenaan pajak karbon, pembatasan emisi industri, larangan ekspor komoditas non-hijau), disrupsi teknologi hijau, dan pergeseran preferensi konsumen. Debitur di sektor berintensitas karbon tinggi (*brown assets*) yang gagal beradaptasi akan menghadapi risiko *stranded assets*, penurunan pendapatan secara drastis, dan pemburukan kapasitas membayar utang.
\end{enumerate}

\noindent
Bank-bank KBMI 4 yang memelopori penyaluran kredit hijau berhasil memitigasi kedua risiko tersebut secara proaktif. Melalui penerapan *Climate Risk Stress Testing* (CRST), bank menyaring debitur dengan profil kerentanan iklim yang rendah, sehingga menciptakan portofolio kredit yang berdaya tahan tinggi dan meminimalkan pembentukan beban CKPN di masa depan.

\subsubsection{Mekanisme Rantai Transmisi Pembiayaan Hijau terhadap Profitabilitas Bank}
\noindent
Transmisi pengaruh alokasi pembiayaan hijau terhadap profitabilitas bank (ROA) beroperasi melalui beberapa jalur (*transmission channels*):
\begin{enumerate}
  \item \textbf{Jalur Pendapatan Bunga dan Imbal Hasil Kredit (*Interest Revenue Channel*):} Portofolio kredit hijau yang disalurkan ke sektor-sektor berkembang pesat seperti pembangkit listrik tenaga surya, infrastruktur efisiensi energi, dan ekosistem kendaraan listrik menghasilkan arus kas pendapatan bunga yang stabil dengan tenor pinjaman jangka menengah-panjang.
  \item \textbf{Jalur Efisiensi Biaya Pendanaan (*Cost of Funds Channel*):} Kepemilikan portofolio hijau yang terverifikasi memungkinkan bank KBMI 4 menerbitkan instrumen obligasi keberlanjutan (*Green Bonds, Social Bonds, Sustainability Bonds*) di pasar modal domestik dan global dengan *greenium* (diskon imbal hasil yang dibayar bank kepada investor karena predikat hijau), sehingga menekan biaya perolehan dana secara keseluruhan.
  \item \textbf{Jalur Kualitas Debitur dan Beban Provisi (*Credit Risk \& Provisioning Channel*):} Debitur proyek hijau memiliki tata kelola lingkungan yang disiplin dan terlindungi dari risiko sanksi pencemaran lingkungan atau penutupan izin usaha oleh pemerintah. Kepatuhan ini menurunkan tingkat gagal bayar, menjaga rasio NPL tetap rendah, dan mengurangi beban pembentukan CKPN PSAK 71.
  \item \textbf{Jalur Pendapatan Berbasis Komisi (*Fee-Based Income Channel*):} Keterlibatan bank KBMI 4 dalam sindikasi pembiayaan hijau skala besar menghasilkan pendapatan non-bunga yang signifikan dari biaya penjaminan (*underwriting fees*), biaya penasihat keuangan (*financial advisory fees*), dan biaya transaksi kustodian hijau.
\end{enumerate}

\subsubsection{Analisis Komparatif Pembiayaan Hijau: Negara Berkembang vs Negara Maju}
\noindent
Penerapan pembiayaan hijau di negara-negara berkembang (*emerging economies*) seperti Indonesia memiliki dinamika struktural yang berbeda dibandingkan di negara-negara maju (*developed economies*):
\begin{enumerate}
  \item \textbf{Ketergantungan Sektor Riil Berkarbon Tinggi:} Di negara berkembang, porsi sektor energi berbasis fosil dan industri ekstraktif masih mendominasi struktur ekonomi, sehingga perbankan menghadapi tantangan *transition risk* yang lebih kompleks. Penyaluran kredit hijau harus diimbangi dengan pembiayaan transisi (*transition finance*) agar tidak memicu kebangkrutan massal pada sektor industri konvensional.
  \item \textbf{Asimetri Informasi dan Standar Pengungkapan:} Ketersediaan data ESG yang diaudit secara independen masih terbatas di pasar negara berkembang dibandingkan negara maju yang telah mengadopsi standar taksonomi ketat seperti EU SFDR. Bank-bank KBMI 4 di Indonesia mengatasi hal ini dengan membangun metodologi asesmen internal (*internal green scoring*) yang mengadopsi POJK 51/2017 dan Taksonomi Hijau Indonesia.
  \item \textbf{Keterbatasan Insentif Fiskal dan Moneter:} Di negara maju, pemerintah memberikan subsidi bunga hijau dan keringanan pajak secara masif. Di Indonesia, dorongan pembiayaan hijau lebih banyak digerakkan oleh kepatuhan regulasi OJK dan Bank Indonesia melalui insentif makroprudensial KLM.
\end{enumerate}

\subsubsection{Dinamika Integrasi ESG ke dalam Penilaian Risiko Kredit Perbankan}
\noindent
Pengintegrasian kriteria *Environmental, Social, and Governance* (ESG) ke dalam proses analisis kelayakan kredit telah mengubah arsitektur pengambilan keputusan perbankan modern:
\begin{enumerate}
  \item \textbf{Penyaringan Awal (*Negative Screening*):} Bank KBMI 4 menetapkan daftar sektor terlarang (*exclusion list*) yang mencakup aktivitas perburuan satwa liar dilindungi, pembalakan liar, dan industri tanpa izin AMDAL.
  \item \textbf{Uji Tuntas Risiko Lingkungan dan Sosial (*ESRA Due Diligence*):} Debitur korporasi dan komersial dinilai berdasarkan kepatuhan pengelolaan limbah, emisi gas rumah kaca, sertifikasi keberlanjutan (ISPO, RSPO, ISO 14001), serta hubungan ketenagakerjaan dan masyarakat sekitar.
  \item \textbf{Penetapan *Covenants* Keberlanjutan:} Perjanjian kredit menyertakan klausul komitmen penurunan emisi (*decarbonization targets*), di mana kegagalan debitur memenuhi target ESG dapat memicu peningkatan marjin suku bunga atau pembatasan penarikan fasilitas kredit.
\end{enumerate}

\subsection{Penelitian Terdahulu}

\noindent
Untuk memberikan landasan komparatif yang kokoh dan mempertegas posisi kebaruan (*novelty*) dari penelitian ini, disajikan matriks 15 penelitian terdahulu yang relevan pada Tabel 2.1 berikut:

\vspace{0.3cm}
\begin{table}[H]
  \caption{Matriks Sintesis Penelitian Terdahulu (Bagian 1: Penelitian 1--5)}
  \label{tab:penelitian_terdahulu_a}
  \footnotesize
  \begin{tabularx}{\textwidth}{|c|>{\raggedright\arraybackslash}p{2.3cm}|>{\raggedright\arraybackslash}X|>{\raggedright\arraybackslash}p{1.9cm}|>{\raggedright\arraybackslash}p{1.8cm}|>{\raggedright\arraybackslash}X|}
    \hline
    \textbf{No} & \textbf{Penulis \& Tahun} & \textbf{Judul Penelitian} & \textbf{Variabel} & \textbf{Metode} & \textbf{Hasil Utama} \\
    \hline
    1 & Pramono, Subroto, \& Saraswati (2022) & \emph{Green Banking and Bank Profitability: Empirical Evidence from Indonesian Banking Sector} & $X$: GF, NPL, CAR, LDR, BOPO; $Y$: ROA, ROE & Regresi Data Panel (FEM) & Green Financing berpengaruh positif signifikan terhadap ROA; NPL berpengaruh negatif signifikan; CAR berpengaruh positif signifikan. \\
    \hline
    2 & Yin, Zhu, Kirkulak-Uludag, \& Zhu (2021) & \emph{Does Green Lending Improve Bank Financial Performance? Evidence from China} & $X$: Green Loan, Size, CAR, NPL; $Y$: ROA, NIM & *System* GMM Panel & Kredit hijau secara signifikan mendongkrak ROA bank melalui efisiensi operasional dan reputasi hijau. \\
    \hline
    3 & Sari \& Astuti (2023) & Pengaruh Penerapan \emph{Green Banking} terhadap Kinerja Keuangan Perbankan di Indonesia & $X$: GF, CAR, NPL, Size; $Y$: ROA & Regresi Linier Berganda & Penyaluran pembiayaan hijau dan CAR berpengaruh positif signifikan terhadap ROA; NPL berpengaruh negatif signifikan. \\
    \hline
    4 & Pradita \& Syaichu (2022) & Analisis Pengaruh NPL, CAR, LDR, dan BOPO terhadap Kinerja Keuangan Bank & $X$: NPL, CAR, LDR, BOPO; $Y$: ROA & Regresi Data Panel & NPL berpengaruh negatif signifikan terhadap ROA, sedangkan CAR berpengaruh positif signifikan pada bank KBMI 4. \\
    \hline
    5 & Wulandari \& Rahardjo (2022) & Pengaruh \emph{Green Financing} dan Manajemen Risiko terhadap Profitabilitas Perbankan & $X$: GF, NPL, CAR, NIM; $Y$: ROA & Regresi Data Panel (FEM) & Green financing dan CAR terbukti meningkatkan ROA secara simultan dan parsial; NPL menurunkan laba. \\
    \hline
  \end{tabularx}
  \begin{flushleft}
    \textit{Sumber: Dikompilasi dan disintesis dari berbagai jurnal ilmiah bereputasi, 2024}
  \end{flushleft}
\end{table}

\begin{table}[H]
  \caption{Matriks Sintesis Penelitian Terdahulu (Bagian 2: Penelitian 6--10)}
  \label{tab:penelitian_terdahulu_b}
  \footnotesize
  \begin{tabularx}{\textwidth}{|c|>{\raggedright\arraybackslash}p{2.3cm}|>{\raggedright\arraybackslash}X|>{\raggedright\arraybackslash}p{1.9cm}|>{\raggedright\arraybackslash}p{1.8cm}|>{\raggedright\arraybackslash}X|}
    \hline
    \textbf{No} & \textbf{Penulis \& Tahun} & \textbf{Judul Penelitian} & \textbf{Variabel} & \textbf{Metode} & \textbf{Hasil Utama} \\
    \hline
    6 & Chiaramonte, Dreassi, Paltrinieri, \& Piser\`a (2020) & \emph{Sustainability Practices and Bank Stability during Financial Distress} & $X$: ESG Score, CAR, NPL; $Y$: Z-score, ROA & Regresi Panel Data & Praktik keberlanjutan dan ESG meningkatkan stabilitas keuangan dan profitabilitas perbankan Eropa. \\
    \hline
    7 & Cui, Geobey, Weber, \& Lin (2018) & \emph{The Evaluation of Green Lending Evaluation in Chinese Commercial Banks} & $X$: Green Credit Ratio, NPL; $Y$: ROE, ROA & *Pooled OLS Regression* & Rasio kredit hijau berkorelasi positif dengan profitabilitas bank komersial di Tiongkok. \\
    \hline
    8 & Sun, Edziah, Sun, \& Kporsu (2021) & \emph{Green Credit Policy and Bank Profitability: Evidence from Commercial Banks in China} & $X$: Green Credit Policy, CAR; $Y$: ROA & *Difference-in-Differences* (DID) & Kebijakan kredit hijau secara bertahap mendorong perbaikan profitabilitas bank komersial. \\
    \hline
    9 & Handayani \& Subowo (2022) & Analisis Pengaruh CAR, NPL, dan Penyaluran Kredit terhadap Profitabilitas Perbankan & $X$: CAR, NPL, Kredit; $Y$: ROA & Regresi Linier Berganda & CAR berpengaruh positif signifikan terhadap ROA; NPL berpengaruh negatif signifikan terhadap profitabilitas. \\
    \hline
    10 & Haryanto \& Sudarno (2023) & Analisis Faktor-faktor Penentu Profitabilitas Bank Umum Konvensional di Indonesia & $X$: CAR, NPL, BOPO, NIM, LDR; $Y$: ROA & Regresi Data Panel (FEM) & CAR dan NIM berpengaruh positif signifikan terhadap ROA; NPL dan BOPO berpengaruh negatif signifikan. \\
    \hline
  \end{tabularx}
  \begin{flushleft}
    \textit{Sumber: Dikompilasi dan disintesis dari berbagai jurnal ilmiah bereputasi, 2024}
  \end{flushleft}
\end{table}

\begin{table}[H]
  \caption{Matriks Sintesis Penelitian Terdahulu (Bagian 3: Penelitian 11--15)}
  \label{tab:penelitian_terdahulu_c}
  \footnotesize
  \begin{tabularx}{\textwidth}{|c|>{\raggedright\arraybackslash}p{2.3cm}|>{\raggedright\arraybackslash}X|>{\raggedright\arraybackslash}p{1.9cm}|>{\raggedright\arraybackslash}p{1.8cm}|>{\raggedright\arraybackslash}X|}
    \hline
    \textbf{No} & \textbf{Penulis \& Tahun} & \textbf{Judul Penelitian} & \textbf{Variabel} & \textbf{Metode} & \textbf{Hasil Utama} \\
    \hline
    11 & Zhang, Yang, \& Bi (2022) & \emph{Green Finance, Green Innovation and Bank Performance: Evidence from Emerging Economies} & $X$: Green Finance, Innovation, NPL; $Y$: ROA & Regresi Data Panel Dinamis & Pembiayaan hijau mendorong inovasi perbankan dan meningkatkan kinerja finansial jangka panjang. \\
    \hline
    12 & Zhou, Sun, Luo, \& Liao (2021) & \emph{Sustainable Finance and Bank Risk: The Role of Green Credit} & $X$: Green Credit Ratio, CAR; $Y$: NPL, Z-score, ROA & Regresi Panel GMM & Penyaluran kredit hijau menurunkan risiko kebangkrutan bank dan mendukung perolehan profitabilitas aset. \\
    \hline
    13 & Buallay (2019) & \emph{Is Sustainability Reporting Associated with Bank Performance?} & $X$: ESG Disclosure, Size, CAR; $Y$: ROA, ROE & Regresi Data Panel & Pengungkapan ESG memiliki korelasi positif signifikan dengan ROA pada sektor perbankan internasional. \\
    \hline
    14 & Lestari \& Purnomo (2023) & Determinan Profitabilitas Bank KBMI 4 di Era Suku Bunga Tinggi & $X$: CAR, NPL, BOPO, NIM; $Y$: ROA & Regresi Data Panel (FEM) & CAR terbukti memperkuat ketahanan laba bank KBMI 4 menghadapi siklus pengetatan moneter. \\
    \hline
    15 & Fitriani \& Wardani (2024) & Implementasi Keuangan Berkelanjutan dan Kinerja Keuangan Bank Umum di Indonesia & $X$: Green Credit, NPL, CAR; $Y$: ROA, Tobin's Q & Regresi Panel Data & Kredit hijau meningkatkan nilai pasar dan ROA bank umum konvensional di Indonesia. \\
    \hline
  \end{tabularx}
  \begin{flushleft}
    \textit{Sumber: Dikompilasi dan disintesis dari berbagai jurnal ilmiah bereputasi, 2024}
  \end{flushleft}
\end{table}

\subsection{Rerangka Pemikiran Konseptual}

\noindent
Berdasarkan telaah teoritis dan sintesis temuan empiris terdahulu, disusunlah rerangka pemikiran konseptual yang menggambarkan keterkaitan kausal antara variabel independen (\emph{Green Financing}, NPL, CAR) dengan variabel dependen (\emph{Return on Assets} / ROA) pada Bank KBMI 4 di Indonesia, sebagaimana disajikan pada Gambar 2.1:

\vspace{0.5cm}
\begin{figure}[H]
  \centering
  \begin{tikzpicture}[
      node distance=2.2cm,
      varbox/.style={rectangle, draw=black, thick, align=center, minimum width=5.2cm, minimum height=1.3cm, font=\small},
      depbox/.style={rectangle, draw=black, thick, align=center, minimum width=4.5cm, minimum height=4.2cm, font=\small},
      arrow/.style={-Stealth, thick}
    ]
    % Variabel Bebas (Independen)
    \node[varbox] (gf) at (0, 3.2) {\textbf{Green Financing ($X_1$)}\\Rasio Portofolio Kredit Hijau / Total Kredit\\(\emph{Stakeholder \& Legitimacy Theory})};
    \node[varbox] (npl) at (0, 0) {\textbf{Non-Performing Loan ($X_2$)}\\Rasio Kredit Bermasalah / Total Kredit\\(\emph{Financial Intermediation Theory})};
    \node[varbox] (car) at (0, -3.2) {\textbf{Capital Adequacy Ratio ($X_3$)}\\Rasio Modal Minimum / ATMR\\(\emph{Risk-Return Tradeoff \& Basel III})};

    % Variabel Terikat (Dependen)
    \node[depbox] (roa) at (8.5, 0) {\textbf{PROFITABILITAS BANK}\\[0.2cm]\textbf{Return on Assets / ROA ($Y$)}\\[0.3cm]$\dfrac{\text{Laba Sebelum Pajak}}{\text{Total Aset Rata-rata}} \times 100\%$\\[0.4cm]\emph{Surat Edaran BI/OJK No. 13/24/DPNP}};

    % Garis Pengaruh Parsial dengan Label Hipotesis
    \draw[arrow] (gf.east) -- node[above, sloped, font=\footnotesize] {$H_1$ (+ Signifikan)} (roa.north west);
    \draw[arrow] (npl.east) -- node[above, font=\footnotesize] {$H_2$ ($-$ Signifikan)} (roa.west);
    \draw[arrow] (car.east) -- node[below, sloped, font=\footnotesize] {$H_3$ (+ Signifikan)} (roa.south west);

    % Kotak dan Label Pengaruh Simultan (H4)
    \node[font=\footnotesize, align=center] at (4.2, -5.2) {\textbf{$H_4$: Pengaruh Simultan ($X_1, X_2, X_3 \rightarrow Y$)}};
    \draw[dashed, thick] (-3.0, -4.6) rectangle (11.2, 4.6);
  \end{tikzpicture}
  \caption{Rerangka Pemikiran Penelitian}
  \label{fig:rerangka_pemikiran}
\end{figure}

\noindent \textit{Sumber: Dikembangkan oleh penulis berdasarkan landasan teoretis dan sintesis \citet{Pramono2022}, \citet{Yin2021}, \citet{KuncoroSuhardjono2018}, dan \citet{BrighamEhrhardt2020}.}

\subsection{Pengembangan Hipotesis Penelitian}

\subsubsection{Pengaruh Portofolio Kredit Hijau (\emph{Green Financing}) terhadap Profitabilitas (ROA)}
\noindent
Berdasarkan premis \emph{Stakeholder Theory} \citep{Freeman2010} dan \emph{Legitimacy Theory} \citep{Suchman1995}, institusi perbankan yang proaktif dalam menyalurkan pembiayaan hijau (\emph{Green Financing}) akan memperoleh legitimasi sosial yang kuat serta kepercayaan tinggi dari seluruh pemangku kepentingan. Kepercayaan ini bermutasi menjadi keunggulan bersaing yang nyata, antara lain: (1) peningkatan citra positif dan reputasi korporasi yang menarik nasabah dana murah (CASA); (2) kemudahan dalam menerbitkan obligasi hijau (*green bonds*) dan instrumen pendanaan ESG global dengan kupon bunga yang lebih rendah; dan (3) kualitas debitur hijau yang memiliki risiko kegagalan operasional lebih kecil dalam menghadapi regulasi lingkungan di masa depan. Meskipun terdapat biaya kepatuhan di awal, efisiensi pendanaan dan keunggulan marjin yang tercipta secara bersih meningkatkan laba sebelum pajak terhadap total aset perbankan.

\noindent
Temuan empiris dari \citet{Pramono2022} pada sektor perbankan Indonesia membuktikan bahwa rasio pembiayaan hijau berpengaruh positif dan signifikan terhadap ROA. Hal ini diperkuat oleh riset \citet{Yin2021} dan \citet{Sun2021} di Tiongkok serta \citet{SariAstuti2023} dan \citet{WulandariRahardjo2022} di Indonesia yang mengonfirmasi bahwa alokasi kredit hijau mendorong peningkatan efisiensi dan profitabilitas bank secara berkelanjutan. Berdasarkan landasan teori dan bukti empiris tersebut, maka dirumuskan hipotesis pertama sebagai berikut:

\vspace{0.2cm}
\noindent \textbf{$H_1$:} \emph{Portofolio Kredit Hijau (Green Financing) berpengaruh positif dan signifikan terhadap Profitabilitas (Return on Assets / ROA) pada Bank KBMI 4 di Indonesia periode 2021--2025.}

\subsubsection{Pengaruh \emph{Non-Performing Loan} (NPL) terhadap Profitabilitas (ROA)}
\noindent
Berdasarkan Teori Intermediasi Finansial \citep{GurleyShaw1960, Diamond1984}, kualitas aktiva produktif merupakan pilar determinan utama kesehatan finansial bank. Rasio \emph{Non-Performing Loan} (NPL) yang meningkat mencerminkan memburuknya kualitas portofolio kredit bank, di mana debitur mengalami gagal bayar atas pokok dan/atau bunga pinjaman. Secara mekanistik akuntansi perbankan, kenaikan NPL menggerus laba bank melalui dua mekanisme utama: pertama, penurunan penerimaan pendapatan bunga (*interest income*) yang merupakan komponen utama pendapatan operasional bank; kedua, keharusan bank membentuk beban penyisihan Cadangan Kerugian Penurunan Nilai (CKPN) yang lebih tinggi sesuai PSAK 71. Beban CKPN dibebankan langsung pada laporan laba rugi tahun berjalan, sehingga secara instan memangkas laba sebelum pajak dan menurunkan ROA.

\noindent
Secara empiris, hubungan negatif antara NPL dan profitabilitas perbankan telah dibuktikan secara konsisten oleh berbagai penelitian terdahulu, di antaranya \citet{PraditaSyaichu2022}, \citet{HandayaniSubowo2022}, \citet{HaryantoSudarno2023}, serta \citet{KuncoroSuhardjono2018}. Semakin tinggi rasio NPL, semakin tertekan kinerja profitabilitas (ROA) bank. Berdasarkan argumen teoritis dan temuan empiris tersebut, maka dirumuskan hipotesis kedua sebagai berikut:

\vspace{0.2cm}
\noindent \textbf{$H_2$:} \emph{Non-Performing Loan (NPL) berpengaruh negatif dan signifikan terhadap Profitabilitas (Return on Assets / ROA) pada Bank KBMI 4 di Indonesia periode 2021--2025.}

\subsubsection{Pengaruh \emph{Capital Adequacy Ratio} (CAR) terhadap Profitabilitas (ROA)}
\noindent
Berdasarkan teori *Risk-Return Tradeoff* \citep{BrighamEhrhardt2020} dan kerangka regulasi *Basel III*, permodalan yang kuat dan kokoh merupakan fondasi utama operasional perbankan. Rasio \emph{Capital Adequacy Ratio} (CAR) yang tinggi menunjukkan bahwa bank memiliki kapasitas modal yang sangat memadai untuk menyerap potensi kerugian tak terduga (*unexpected losses*) tanpa mengganggu kelangsungan usaha bank. Modal yang kuat memberikan sinyal positif (*signaling theory*) kepada pasar mengenai keamanan bank, sehingga meningkatkan kepercayaan deposan, menurunkan premi risiko, dan memangkas biaya perolehan dana (*cost of funds*).

\noindent
Di samping itu, bank dengan rasio CAR yang tebal memiliki fleksibilitas dan kapasitas yang lebih leluasa untuk melakukan ekspansi aset produktif, termasuk pembiayaan proyek-proyek hijau dan korporasi berskala besar dengan marjin imbal hasil yang menguntungkan. Temuan empiris dari \citet{LestariPurnomo2023}, \citet{NugrohoUtami2021}, \citet{Pramono2022}, dan \citet{HaryantoSudarno2023} membuktikan secara konsisten bahwa CAR memiliki pengaruh positif dan signifikan terhadap ROA pada bank-bank umum konvensional di Indonesia. Berdasarkan penalaran teoritis dan dukungan bukti empiris tersebut, maka dirumuskan hipotesis ketiga sebagai berikut:

\vspace{0.2cm}
\noindent \textbf{$H_3$:} \emph{Capital Adequacy Ratio (CAR) berpengaruh positif dan signifikan terhadap Profitabilitas (Return on Assets / ROA) pada Bank KBMI 4 di Indonesia periode 2021--2025.}

\subsubsection{Pengaruh Simultan \emph{Green Financing}, NPL, dan CAR terhadap Profitabilitas (ROA)}
\noindent
Secara terintegrasi, profitabilitas perbankan merupakan hasil interaksi dinamis antara strategi alokasi portofolio bisnis, efektivitas mitigasi risiko kredit, dan kecukupan struktur permodalan internal. Portofolio pembiayaan hijau (\emph{Green Financing}) berfungsi sebagai instrumen penciptaan nilai baru dan pendorong reputasi ESG; pengendalian rasio kredit bermasalah (NPL) berfungsi sebagai benteng perlindungan penerimaan pendapatan bunga dan pencegahan erosi laba akibat beban CKPN; sedangkan ketahanan rasio kecukupan modal (CAR) berfungsi sebagai fondasi stabilitas dan kapasitas ekspansi aset produktif.

\noindent
Kombinasi harmonis antara alokasi kredit hijau yang agresif namun selektif, kontrol risiko kredit macet yang disiplin, serta kecukupan modal yang tangguh secara bersama-sama akan menciptakan ekosistem operasional bank yang sangat kondusif bagi pencapaian laba aset (ROA) yang optimal dan berdaya tahan tinggi terhadap siklus ekonomi. Hal ini sejalan dengan temuan \citet{Pramono2022} dan \citet{WulandariRahardjo2022} yang membuktikan adanya pengaruh simultan yang sangat signifikan dari variabel-variabel tersebut terhadap kinerja profitabilitas bank. Berdasarkan kerangka pemikiran tersebut, dirumuskan hipotesis keempat sebagai berikut:

\vspace{0.2cm}
\noindent \textbf{$H_4$:} \emph{Portofolio Kredit Hijau (Green Financing), Non-Performing Loan (NPL), dan Capital Adequacy Ratio (CAR) secara simultan (bersama-sama) berpengaruh signifikan terhadap Profitabilitas (Return on Assets / ROA) pada Bank KBMI 4 di Indonesia periode 2021--2025.}

% ============================================================
%  BAB 3: METODE PENELITIAN
% ============================================================
\newpage
\setcounter{section}{0}
\section*{BAB 3\\[0.3cm]METODE PENELITIAN}
\addcontentsline{toc}{section}{\texorpdfstring{BAB 3\quad METODE PENELITIAN}{BAB 3 - METODE PENELITIAN}}
\setcounter{section}{3}
\setcounter{subsection}{0}
\setcounter{equation}{0}

\subsection{Desain Penelitian dan Paradigma Keilmuan}

\noindent
Penelitian ini dirancang menggunakan pendekatan \textbf{kuantitatif asosiatif kausalitas} yang berlandaskan pada paradigma \textbf{positivisme} \citep{Sugiyono2019}. Pendekatan kuantitatif berfokus pada pengujian teori-teori objektif melalui pengujian hubungan antarvariabel yang diukur secara numerik dan dianalisis menggunakan prosedur statistik dan ekonometrika. Desain asosiatif kausal bertujuan untuk membuktikan hubungan sebab-akibat antara variabel independen, yaitu Portofolio Kredit Hijau (\emph{Green Financing} / $X_1$), \emph{Non-Performing Loan} (NPL / $X_2$), dan \emph{Capital Adequacy Ratio} (CAR / $X_3$) terhadap variabel dependen, yaitu Profitabilitas yang diproksikan dengan \emph{Return on Assets} (ROA / $Y$).

\noindent
Struktur data yang digunakan dalam penelitian ini adalah \textbf{data panel} (\emph{panel data / pooled time-series and cross-section data}), yaitu penggabungan antara data deret waktu (*time-series*) selama 20 kuartal (2021--2025) dengan data lintas entitas (*cross-section*) sebanyak 4 bank KBMI 4, menghasilkan struktur data panel seimbang (*balanced panel data*).

\subsection{Populasi, Sampel, dan Teknik Pengambilan Sampel}

\subsubsection{Populasi Penelitian}
\noindent
Populasi sasaran dalam penelitian ini adalah seluruh Bank Umum Konvensional yang terdaftar di Bursa Efek Indonesia (BEI) selama periode pengamatan tahun 2021 sampai dengan tahun 2025.

\subsubsection{Teknik Pengambilan Sampel}
\noindent
Pengambilan sampel dilakukan menggunakan teknik *non-probability sampling* dengan metode \textbf{purposive sampling}, yaitu penentuan sampel berdasarkan kriteria pertimbangan tertentu (*judgment sampling*) yang diselaraskan dengan tujuan dan kebutuhan analisis penelitian \citep{Sugiyono2019}. Kriteria inklusi sampel yang ditetapkan adalah sebagai berikut:
\begin{enumerate}
  \item Bank umum konvensional yang terdaftar di Bursa Efek Indonesia (BEI) dan masuk dalam klasifikasi \textbf{KBMI 4} (Modal Inti $> \text{Rp}70 \text{ Triliun}$) berdasarkan POJK No. 12/POJK.03/2021 secara konsisten berturut-turut selama periode 2021--2025.
  \item Menerbitkan Laporan Keuangan Tahunan (*Annual Report*) dan Laporan Keuangan Triwulanan publikasi resmi yang telah diaudit selama periode kuartal I-2021 sampai kuartal IV-2025 secara lengkap.
  \item Menerbitkan Laporan Keberlanjutan (*Sustainability Report*) tahunan yang menyajikan pengungkapan data alokasi pembiayaan Kegiatan Usaha Berwawasan Lingkungan (KUBL) / *Green Financing* selama periode 2021--2025.
  \item Menghasilkan laba bersih tahunan positif (tidak membukukan kerugian operasional bersih) selama seluruh periode observasi penelitian agar rasio ROA terukur secara valid dan konsisten.
\end{enumerate}

\subsubsection{Sampel Final Penelitian}
\noindent
Berdasarkan tahapan seleksi purposive sampling dengan kriteria di atas, diperoleh sampel final sebanyak \textbf{4 bank umum KBMI 4}, sebagaimana disajikan pada Tabel 3.1:

\begin{table}[H]
  \caption{Daftar Entitas Sampel Penelitian Bank KBMI 4}
  \label{tab:sampel_penelitian}
  \begin{tabularx}{\textwidth}{|c|c|>{\raggedright\arraybackslash}X|c|c|}
    \hline
    \textbf{No} & \textbf{Kode Saham} & \textbf{Nama Entitas Perbankan} & \textbf{Kategori KBMI} & \textbf{Periode Observasi} \\
    \hline
    1 & BBRI & PT Bank Rakyat Indonesia (Persero) Tbk & KBMI 4 & 2021-Q1 s.d. 2025-Q4 \\
    \hline
    2 & BMRI & PT Bank Mandiri (Persero) Tbk & KBMI 4 & 2021-Q1 s.d. 2025-Q4 \\
    \hline
    3 & BBCA & PT Bank Central Asia Tbk & KBMI 4 & 2021-Q1 s.d. 2025-Q4 \\
    \hline
    4 & BBNI & PT Bank Negara Indonesia (Persero) Tbk & KBMI 4 & 2021-Q1 s.d. 2025-Q4 \\
    \hline
  \end{tabularx}
  \begin{flushleft}
    \textit{Sumber: Bursa Efek Indonesia (BEI) \& Laporan Otoritas Jasa Keuangan, 2025}
  \end{flushleft}
\end{table}

\noindent
Dengan jumlah entitas bank $N = 4$ dan jumlah periode kuartal $T = 20$ (4 kuartal $\times$ 5 tahun), maka total sampel observasi panel yang dianalisis dalam penelitian ini adalah:
\begin{equation}
  \text{Total Observasi Panel} = N \times T = 4 \text{ Bank} \times 20 \text{ Kuartal} = \mathbf{80 \text{ Observasi}}
\end{equation}

\subsection{Jenis, Sumber, dan Teknik Pengumpulan Data}

\subsubsection{Jenis Data}
\noindent
Jenis data yang digunakan adalah \textbf{data sekunder kuantitatif}, yaitu data numerik yang telah disusun dan dipublikasikan oleh pihak ketiga yang memiliki otoritas resmi.

\subsubsection{Sumber Data}
\noindent
Data sekunder diperoleh dari sumber-sumber resmi yang dapat dipertanggungjawabkan validitasnya, antara lain:
\begin{enumerate}
  \item Laporan Keuangan Publikasi Triwulanan dan Tahunan (*Annual Report*) masing-masing bank yang diunduh dari situs resmi Bursa Efek Indonesia (\url{www.idx.co.id}) dan portal *Investor Relations* perbankan terkait.
  \item Laporan Keberlanjutan (*Sustainability Report*) tahunan masing-masing bank periode 2021--2025.
  \item Statistik Perbankan Indonesia (SPI) yang diterbitkan secara berkala oleh Otoritas Jasa Keuangan (\url{www.ojk.go.id}).
  \item Publikasi indikator makroekonomi dan suku bunga acuan dari Bank Indonesia (\url{www.bi.go.id}).
\end{enumerate}

\subsubsection{Teknik Pengumpulan Data}
\noindent
Teknik pengumpulan data dilakukan melalui metode \textbf{dokumentasi} dan \textbf{studi kepustakaan} (*library research*), yaitu mencatat, mentabulasi, merekonsiliasi, dan menghitung angka-angka rasio keuangan dari catatan atas laporan keuangan (CALK) dan laporan keberlanjutan ke dalam format basis data ekonometrika.

\subsection{Operasionalisasi Variabel}

\noindent
Operasionalisasi variabel memaparkan definisi konseptual, indikator pengukuran empiris, rumus matematis, skala data, dan sumber data dari setiap variabel penelitian, sebagaimana disajikan pada Tabel 3.2 berikut:

\begin{table}[H]
  \caption{Matriks Operasionalisasi Variabel Penelitian}
  \label{tab:operasionalisasi_variabel}
  \small
  \begin{tabularx}{\textwidth}{|>{\raggedright\arraybackslash}p{2.1cm}|>{\raggedright\arraybackslash}X|>{\centering\arraybackslash}p{4.1cm}|c|>{\raggedright\arraybackslash}p{2.1cm}|}
    \hline
    \textbf{Variabel} & \textbf{Definisi Konseptual} & \textbf{Indikator / Rumus Matematis} & \textbf{Skala} & \textbf{Sumber Data} \\
    \hline
    \textbf{Return on Assets (ROA)} ($Y$) & Kemampuan manajemen menghasilkan laba sebelum pajak dari total aktiva yang dikelola \citep{Dendawijaya2015} & $\text{ROA} = \left(\dfrac{\text{Laba Sblm Pajak}}{\text{Total Aset Rata-rata}}\right) \times 100\%$ & Rasio (\%) & Laporan Laba Rugi Publikasi BEI \\
    \hline
    \textbf{Green Financing (GF)} ($X_1$) & Porsi alokasi kredit pembiayaan Kegiatan Usaha Berwawasan Lingkungan terhadap total kredit \citep{OJK2017} & $\text{GF} = \left(\dfrac{\text{Kredit Hijau KUBL}}{\text{Total Kredit Disalurkan}}\right) \times 100\%$ & Rasio (\%) & \emph{Sustainability Report} \& RAKB \\
    \hline
    \textbf{Non-Performing Loan (NPL)} ($X_2$) & Rasio kredit bermasalah (Kolektibilitas 3, 4, 5) terhadap total penyaluran kredit \citep{KuncoroSuhardjono2018} & $\text{NPL} = \left(\dfrac{\text{Kredit NPL Kol 3,4,5}}{\text{Total Kredit Disalurkan}}\right) \times 100\%$ & Rasio (\%) & Laporan Keuangan Publikasi / CALK \\
    \hline
    \textbf{Capital Adequacy Ratio (CAR)} ($X_3$) & Rasio kecukupan modal bank dalam menanggung risiko aktiva tertimbang \citep{BrighamEhrhardt2020} & $\text{CAR} = \left(\dfrac{\text{Total Modal (Tier 1+2)}}{\text{ATMR}}\right) \times 100\%$ & Rasio (\%) & Laporan Kewajiban KPMM / CALK \\
    \hline
  \end{tabularx}
  \begin{flushleft}
    \textit{Sumber: Dikembangkan oleh penulis dari regulasi OJK dan literatur manajemen keuangan, 2024}
  \end{flushleft}
\end{table}

\subsection{Spesifikasi Model Regresi Data Panel}

\noindent
Untuk menguji hipotesis penelitian secara empiris, disusun persamaan model regresi data panel linier berganda sebagai berikut \citep{GujaratiPorter2020, Wooldridge2020}:

\begin{equation}
  \text{ROA}_{it} = \alpha_i + \beta_1 \text{GF}_{it} + \beta_2 \text{NPL}_{it} + \beta_3 \text{CAR}_{it} + \varepsilon_{it}
  \label{eq:regresi_panel}
\end{equation}

\noindent Di mana:
\begin{itemize}
  \item $\text{ROA}_{it}$ : Kinerja profitabilitas (\emph{Return on Assets}) bank $i$ pada kuartal $t$ (\%)
  \item $\alpha_i$ : Intersep spesifik individual masing-masing entitas bank $i$
  \item $\beta_1, \beta_2, \beta_3$ : Koefisien regresi parameter masing-masing variabel independen
  \item $\text{GF}_{it}$ : Rasio portofolio pembiayaan hijau (\emph{Green Financing}) bank $i$ pada kuartal $t$ (\%)
  \item $\text{NPL}_{it}$ : Rasio kredit bermasalah (\emph{Non-Performing Loan} Gross) bank $i$ pada kuartal $t$ (\%)
  \item $\text{CAR}_{it}$ : Rasio kecukupan modal (\emph{Capital Adequacy Ratio}) bank $i$ pada kuartal $t$ (\%)
  \item $\varepsilon_{it}$ : Galat pengganggu (*error term*) gabungan waktu dan individu
  \item $i$ : Entitas unit cross-section ($1 = \text{BBRI}, 2 = \text{BMRI}, 3 = \text{BBCA}, 4 = \text{BBNI}$)
  \item $t$ : Periode waktu kuartalan ($t = 1, 2, \dots, 20$, melambangkan Q1-2021 s.d. Q4-2025)
\end{itemize}

\subsection{Tahapan Analisis Data dan Uji Ekonometrika}

\noindent
Pengolahan dan pengujian data dilakukan secara sistematis menggunakan bantuan perangkat lunak ekonometrika \textbf{EViews 12}, melalui tahapan-tahapan sebagai berikut:

\subsubsection{Analisis Statistik Deskriptif}
\noindent
Statistik deskriptif digunakan untuk menyajikan karakteristik distribusi data secara komprehensif, mencakup nilai rata-rata hitung (\emph{mean}), nilai tengah (\emph{median}), nilai tertinggi (\emph{maximum}), nilai terendah (\emph{minimum}), ukuran penyebaran simpangan baku (\emph{standard deviation}), serta ukuran kemencengan distribusi (*skewness*) dan keruncingan (*kurtosis*).

\subsubsection{Uji Stasioneritas / Unit Root Data Panel}
\noindent
Sebelum melakukan estimasi regresi, dilakukan uji akar unit data panel (*panel unit root test*) untuk memastikan bahwa data tidak mengandung tren stokastik yang dapat menghasilkan regresi palsu (*spurious regression*). Pengujian dilakukan menggunakan metode \textbf{Levin, Lin \& Chu (LLC)} dan \textbf{Im, Pesaran and Shin (IPS) W-stat}. Data dinyatakan stasioner jika nilai $p$-value $< 0{,}05$ pada tingkat level atau *first difference* \citep{GujaratiPorter2020}.

\subsubsection{Teknik Estimasi Model Regresi Data Panel}
\noindent
Dalam ekonometrika data panel, terdapat tiga pendekatan model estimasi:
\begin{enumerate}
  \item \textbf{\emph{Common Effect Model} (CEM) / \emph{Pooled OLS}:} Mengasumsikan bahwa perilaku antar-bank adalah homogen (tidak ada perbedaan intersep maupun kemiringan antar-entitas bank sepanjang waktu).
  \item \textbf{\emph{Fixed Effect Model} (FEM) / \emph{Least Squares Dummy Variable} (LSDV):} Mengasumsikan bahwa setiap bank memiliki karakteristik unik yang tidak teramati (*unobserved heterogeneity*) yang tercermin dalam perbedaan nilai intersep ($\alpha_i$), namun kemiringan koefisien ($\beta$) antar-bank diasumsikan tetap.
  \item \textbf{\emph{Random Effect Model} (REM) / \emph{Generalized Least Squares} (GLS):} Mengasumsikan bahwa perbedaan karakteristik individual antar-bank bersifat acak dan diakomodasi sebagai bagian dari galat pengganggu (*error term*).
\end{enumerate}

\subsubsection{Uji Pemilihan Model Regresi Data Panel}
\noindent
Untuk menentukan model terbaik dari ketiga pendekatan di atas, dilakukan serangkaian uji formal:
\begin{enumerate}
  \item \textbf{Uji Chow (\emph{Chow Test}):} Menguji perbandingan kelayakan antara CEM dengan FEM. Formulasi statistik uji $F$-Chow adalah:
    \begin{equation}
      F_{\text{Chow}} = \frac{(\text{RSS}_{\text{CEM}} - \text{RSS}_{\text{FEM}}) / (N - 1)}{\text{RSS}_{\text{FEM}} / (NT - N - K)}
    \end{equation}
    Di mana $\text{RSS}$ adalah *Residual Sum of Squares*, $N$ adalah jumlah bank ($4$), $T$ adalah jumlah kuartal ($20$), dan $K$ adalah jumlah variabel independen ($3$).
    \begin{itemize}
      \item $H_0$: \emph{Common Effect Model} (CEM) lebih tepat.
      \item $H_1$: \emph{Fixed Effect Model} (FEM) lebih tepat.
      \item Kriteria Keputusan: Jika nilai $\text{Prob. Cross-section F} < 0{,}05$, maka $H_0$ ditolak dan FEM terpilih.
    \end{itemize}
  \item \textbf{Uji Hausman (\emph{Hausman Test}):} Menguji perbandingan kelayakan antara FEM dengan REM. Statistik uji Hausman mengikuti distribusi Chi-Square ($\chi^2$):
    \begin{equation}
      W = (\hat{\beta}_{\text{FEM}} - \hat{\beta}_{\text{REM}})' \left[ \text{Var}(\hat{\beta}_{\text{FEM}}) - \text{Var}(\hat{\beta}_{\text{REM}}) \right]^{-1} (\hat{\beta}_{\text{FEM}} - \hat{\beta}_{\text{REM}}) \sim \chi^2(K)
    \end{equation}
    \begin{itemize}
      \item $H_0$: \emph{Random Effect Model} (REM) lebih tepat (efek individu bersifat acak dan tidak berkorelasi dengan regressor).
      \item $H_1$: \emph{Fixed Effect Model} (FEM) lebih tepat (efek individu berkorelasi dengan regressor).
      \item Kriteria Keputusan: Jika nilai $\text{Prob. Cross-section Random} < 0{,}05$, maka $H_0$ ditolak dan FEM ditetapkan sebagai model terbaik.
    \end{itemize}
  \item \textbf{Uji \emph{Lagrange Multiplier} (LM Test):} Digunakan untuk membandingkan antara CEM dengan REM berdasarkan metode Breusch-Pagan:
    \begin{equation}
      LM = \frac{NT}{2(T-1)} \left[ \frac{\sum_{i=1}^N \left( \sum_{t=1}^T e_{it} \right)^2}{\sum_{i=1}^N \sum_{t=1}^T e_{it}^2} - 1 \right]^2 \sim \chi^2(1)
    \end{equation}
\end{enumerate}

\subsubsection{Uji Asumsi Klasik Ekonometrika (BLUE)}
\noindent
Setelah model panel terbaik terpilih, dilakukan pengujian asumsi klasik untuk memastikan bahwa parameter estimator memenuhi kriteria *Best Linear Unbiased Estimator* (BLUE) \citep{GujaratiPorter2020}:
\begin{enumerate}
  \item \textbf{Uji Normalitas Residual:} Menguji apakah residual model terdistribusi normal menggunakan uji \textbf{Jarque-Bera (JB)}:
    \begin{equation}
      JB = \frac{N}{6} \left[ S^2 + \frac{(K - 3)^2}{4} \right] \sim \chi^2(2)
    \end{equation}
    Di mana $S$ adalah *skewness* dan $K$ adalah *kurtosis*. Residual dinyatakan berdistribusi normal apabila nilai $\text{Prob. } JB > 0{,}05$.
  \item \textbf{Uji Multikolinearitas:} Menguji ada tidaknya korelasi linier sempurna antarvariabel independen menggunakan nilai \textbf{Variance Inflation Factor (VIF)}:
    \begin{equation}
      \text{VIF}_j = \frac{1}{1 - R_j^2}
    \end{equation}
    Model dinyatakan bebas dari gangguan multikolinearitas apabila nilai $\text{VIF} < 10$.
  \item \textbf{Uji Heteroskedastisitas:} Menguji ada tidaknya ketidaksamaan varians residual antar-observasi menggunakan \textbf{Uji Glejser} (meregresikan nilai absolut residual $|e_{it}|$ terhadap variabel independen). Model dinyatakan homoskedastis apabila seluruh nilai signifikansi $> 0{,}05$.
  \item \textbf{Uji Autokorelasi:} Menguji ada tidaknya korelasi serial antar-kesalahan pengganggu menggunakan uji \textbf{Durbin-Watson (DW)}:
    \begin{equation}
      d = \frac{\sum_{t=2}^T (e_t - e_{t-1})^2}{\sum_{t=1}^T e_t^2}
    \end{equation}
    Model dinyatakan bebas autokorelasi apabila $d_U < d < 4 - d_U$.
\end{enumerate}

\subsubsection{Uji Kelayakan Model dan Pengujian Hipotesis}
\noindent
Pengujian hipotesis dilakukan pada tingkat signifikansi $\alpha = 5\%$ ($0{,}05$) melalui:
\begin{enumerate}
  \item \textbf{Koefisien Determinasi (\emph{Adjusted} $R^2$):}
    \begin{equation}
      \text{Adjusted } R^2 = 1 - \left[ \frac{(1 - R^2)(NT - 1)}{NT - K - 1} \right]
    \end{equation}
  \item \textbf{Uji Signifikansi Simultan (Uji F):}
    \begin{equation}
      F = \frac{R^2 / K}{(1 - R^2) / (NT - K - 1)}
    \end{equation}
  \item \textbf{Uji Signifikansi Parsial (Uji t):}
    \begin{equation}
      t_{\text{hitung}} = \frac{\hat{\beta}_j - 0}{\text{SE}(\hat{\beta}_j)}
    \end{equation}
\end{enumerate}

\subsubsection{Spesifikasi Aljabar Matriks Model Regresi Data Panel}
\noindent
Dalam notasi matriks komprehensif \citep{Greene2018, Wooldridge2020}, sistem persamaan data panel untuk $N = 4$ dan $T = 20$ (total observasi $NT = 80$) dapat diformulasikan sebagai:
\begin{equation}
  \mathbf{y} = \mathbf{D}\boldsymbol{\alpha} + \mathbf{X}\boldsymbol{\beta} + \boldsymbol{\varepsilon}
\end{equation}
Di mana:
\begin{equation}
  \mathbf{y} = \begin{bmatrix} \mathbf{y}_1 \\ \mathbf{y}_2 \\ \mathbf{y}_3 \\ \mathbf{y}_4 \end{bmatrix}_{80 \times 1}, \quad
  \mathbf{D} = \begin{bmatrix} \boldsymbol{\iota}_T & \mathbf{0} & \mathbf{0} & \mathbf{0} \\ \mathbf{0} & \boldsymbol{\iota}_T & \mathbf{0} & \mathbf{0} \\ \mathbf{0} & \mathbf{0} & \boldsymbol{\iota}_T & \mathbf{0} \\ \mathbf{0} & \mathbf{0} & \mathbf{0} & \boldsymbol{\iota}_T \end{bmatrix}_{80 \times 4}, \quad
  \boldsymbol{\alpha} = \begin{bmatrix} \alpha_1 \\ \alpha_2 \\ \alpha_3 \\ \alpha_4 \end{bmatrix}_{4 \times 1}, \quad
  \boldsymbol{\beta} = \begin{bmatrix} \beta_1 \\ \beta_2 \\ \beta_3 \end{bmatrix}_{3 \times 1}
\end{equation}
Vektor $\mathbf{X}$ merepresentasikan matriks berukuran $80 \times 3$ yang memuat variabel \emph{Green Financing}, NPL Gross, dan CAR. Estimator *Within Transformation* untuk \emph{Fixed Effect Model} (FEM) diperoleh melalui proyeksi ortogonal menggunakan matriks penyapu rata-rata $\mathbf{Q} = \mathbf{I}_{NT} - \mathbf{D}(\mathbf{D}'\mathbf{D})^{-1}\mathbf{D}'$:
\begin{equation}
  \hat{\boldsymbol{\beta}}_{\text{FEM}} = (\mathbf{X}'\mathbf{Q}\mathbf{X})^{-1} \mathbf{X}'\mathbf{Q}\mathbf{y}
\end{equation}
Estimator $\hat{\boldsymbol{\beta}}_{\text{FEM}}$ bersifat konsisten dan tidak bias (*unbiased*) sekalipun terdapat korelasi antara efek individual bank ($\alpha_i$) dengan variabel penjelas ($\mathbf{X}$).

\subsubsection{Formulasi Sintaks Perintah EViews dan Alur Kerja Estimasi}
\noindent
Proses estimasi dan pengujian ekonometrika di dalam perangkat lunak EViews 12 dieksekusi melalui urutan perintah sintaks (*command line script*) sebagai berikut:
\begin{enumerate}
  \item \textbf{Deklarasi Struktur Panel:}
\begin{verbatim}
wfcreate(wf=skripsi_kbmi4) m 2021Q1 2025Q4 4
pagestruct id_bank @date
\end{verbatim}
  \item \textbf{Estimasi Model Fixed Effects:}
\begin{verbatim}
equation eq_fem.ls(cx=f) roa c green_financing npl_gross car
\end{verbatim}
  \item \textbf{Uji Spesifikasi dan Diagnostik:}
\begin{verbatim}
eq_fem.fixedtest                'Uji Chow Redundant Fixed Effects
eq_rem.hausman                  'Uji Hausman Correlated Random Effects
eq_fem.vif                      'Uji Variance Inflation Factors
eq_fem.resids(g) resid_fem      'Ekstraksi Residual untuk Jarque-Bera & Glejser
\end{verbatim}
\end{enumerate}

% ============================================================
%  BAB 4: ANALISIS DAN PEMBAHASAN
% ============================================================
\newpage
\setcounter{section}{0}
\section*{BAB 4\\[0.3cm]ANALISIS DAN PEMBAHASAN}
\addcontentsline{toc}{section}{\texorpdfstring{BAB 4\quad ANALISIS DAN PEMBAHASAN}{BAB 4 - ANALISIS DAN PEMBAHASAN}}
\setcounter{section}{4}
\setcounter{subsection}{0}
\setcounter{equation}{0}

\subsection{Gambaran Umum Objek Penelitian}

\noindent
Objek yang diteliti dalam skripsi ini adalah seluruh bank umum konvensional yang terdaftar di Bursa Efek Indonesia (BEI) yang masuk dalam kelompok modal tertinggi, yaitu \textbf{Kelompok Bank Berdasarkan Modal Inti (KBMI) 4} dengan modal inti di atas Rp70 Triliun selama periode 2021 sampai dengan 2025. Berdasarkan data Otoritas Jasa Keuangan, hanya terdapat empat entitas perbankan yang memenuhi kriteria tersebut. Berikut adalah profil singkat dan perkembangan inisiatif *green banking* dari masing-masing bank sampel:

\subsubsection{PT Bank Rakyat Indonesia (Persero) Tbk (BBRI)}
\noindent
PT Bank Rakyat Indonesia (Persero) Tbk didirikan pada tahun 1895 dan merupakan bank komersial tertua serta terbesar di Indonesia dengan fokus utama pada pembiayaan segmen usaha mikro, kecil, dan menengah (UMKM). Selama periode 2021--2025, BBRI menunjukkan komitmen yang sangat agresif dalam penerapan keuangan berkelanjutan. Bank BRI menerbitkan *Sustainability Bond* dan *Green Bond* perdana senilai puluhan triliun rupiah untuk mendanai proyek-proyek energi terbarukan, pengelolaan sumber daya alam berkelanjutan, dan pembiayaan mikro berwawasan lingkungan. Portofolio pembiayaan KUBL BRI tumbuh konsisten dari 18,20\% pada awal 2021 hingga mencapai puncaknya sebesar 31,50\% dari total portofolio kredit pada akhir tahun 2025.

\subsubsection{PT Bank Mandiri (Persero) Tbk (BMRI)}
\noindent
PT Bank Mandiri (Persero) Tbk merupakan bank hasil merger empat bank milik pemerintah pada tahun 1998 dengan keunggulan pada segmen korporasi, komersial, dan institusi pemerintah. Bank Mandiri secara konsisten memposisikan dirinya sebagai *market leader* dalam pembiayaan sindikasi proyek infrastruktur hijau, pembangkit listrik tenaga air (PLTA), pembangkit listrik tenaga surya (PLTS), serta pembiayaan ekosistem kendaraan listrik (*electric vehicle ecosystem*). Portofolio *Green Financing* Bank Mandiri mengalami ekspansi pesat dari 17,50\% di kuartal I-2021 menjadi 30,00\% pada kuartal IV-2025, didukung oleh penerbitan *ESG Framework* dan *Green Bond* yang terserap pasar dengan sangat baik.

\subsubsection{PT Bank Central Asia Tbk (BBCA)}
\noindent
PT Bank Central Asia Tbk merupakan bank swasta terbesar di Indonesia yang unggul dalam jaringan transaksi perbankan digital, keandalan operasional, dan dominasi dana murah (*Current Account and Saving Account* / CASA). Dalam implementasi *green banking*, BCA berfokus pada pembiayaan bangunan ramah lingkungan (*green building*), industri manufaktur yang menerapkan efisiensi energi dan air, perkebunan kelapa sawit bersertifikasi keberlanjutan (ISPO/RSPO), serta pembiayaan kendaraan bermotor listrik melalui lini pembiayaan konsumernya. Rasio pembiayaan hijau BCA meningkat dari 16,20\% pada awal 2021 menjadi 30,40\% pada akhir 2025, diimbangi dengan kualitas aset yang sangat prima dan rasio NPL terendah di industri perbankan nasional.

\subsubsection{PT Bank Negara Indonesia (Persero) Tbk (BBNI)}
\noindent
PT Bank Negara Indonesia (Persero) Tbk didirikan pada tahun 1946 dengan fokus pada perbankan korporasi, bisnis internasional, dan konsumer. BNI menjadi salah satu pelopor penerbitan *Green Bond* berdenominasi Rupiah di BEI pada tahun 2022 untuk mendukung proyek energi terbarukan, pengolahan air limbah terpadu, dan transportasi berkelanjutan. Portofolio kredit hijau BNI tumbuh stabil dari 17,10\% pada awal 2021 menjadi 28,90\% pada akhir 2025 seiring dengan implementasi uji tuntas lingkungan yang terintegrasi ke dalam analisis komite kredit bank.

\subsection{Analisis Tren Kuartalan dan Profil Perbankan KBMI 4 (2021--2025)}

\subsection{Analisis Tren Kuartalan dan Profil Finansial Perbankan KBMI 4 (2021--2025)}

\subsubsection{Perkembangan Portofolio dan Rasio Keuangan PT Bank Rakyat Indonesia (Persero) Tbk (BBRI)}
\noindent
Sebagai bank dengan jaringan nasabah terbesar di Indonesia, Bank BRI mencatatkan transformasi portofolio pembiayaan hijau yang bertumpu pada sektor pertanian berkelanjutan bersertifikat, perkebunan rakyat lestari, kehutanan ramah lingkungan, serta program elektrifikasi desa. Selama kurun waktu 2021--2025, integrasi Holding Ultra Mikro (bersama Pegadaian dan PNM) mempercepat penyaluran pembiayaan hijau di tingkat akar rumput (*grassroots ESG financing*). Rincian perkembangan variabel keuangan kuartalan Bank BRI disajikan pada Tabel 4.1b:

\begin{table}[H]
  \caption{Perkembangan Kinerja Keuangan Triwulanan PT Bank Rakyat Indonesia (Persero) Tbk (2021--2025)}
  \label{tab:kinerja_bbri}
  \small
  \begin{tabularx}{\textwidth}{|c|c|c|c|c|}
    \hline
    \textbf{Periode Kuartal} & \textbf{Green Financing (\%)} & \textbf{NPL Gross (\%)} & \textbf{CAR (\%)} & \textbf{ROA (\%)} \\
    \hline
    2021-Q1 & 18,20 & 3,15 & 19,80 & 2,45 \\
    2021-Q2 & 18,50 & 3,10 & 20,10 & 2,55 \\
    2021-Q3 & 19,10 & 3,05 & 20,40 & 2,68 \\
    2021-Q4 & 19,80 & 2,95 & 21,20 & 2,82 \\
    \hline
    2022-Q1 & 20,40 & 2,90 & 21,50 & 2,95 \\
    2022-Q2 & 20,90 & 2,85 & 21,80 & 3,05 \\
    2022-Q3 & 21,50 & 2,75 & 22,10 & 3,15 \\
    2022-Q4 & 22,10 & 2,70 & 22,40 & 3,25 \\
    \hline
    2023-Q1 & 23,20 & 2,65 & 22,70 & 3,30 \\
    2023-Q2 & 24,10 & 2,60 & 23,00 & 3,38 \\
    2023-Q3 & 24,80 & 2,55 & 23,30 & 3,45 \\
    2023-Q4 & 25,50 & 2,50 & 23,60 & 3,52 \\
    \hline
    2024-Q1 & 26,20 & 2,45 & 23,90 & 3,58 \\
    2024-Q2 & 27,00 & 2,40 & 24,20 & 3,65 \\
    2024-Q3 & 27,80 & 2,35 & 24,50 & 3,72 \\
    2024-Q4 & 28,60 & 2,30 & 24,80 & 3,80 \\
    \hline
    2025-Q1 & 29,20 & 2,25 & 25,00 & 3,88 \\
    2025-Q2 & 30,00 & 2,20 & 25,40 & 3,96 \\
    2025-Q3 & 30,80 & 2,15 & 25,80 & 4,05 \\
    2025-Q4 & 31,50 & 2,10 & 26,20 & 4,15 \\
    \hline
  \end{tabularx}
  \begin{flushleft}
    \textit{Sumber: Laporan Keuangan Publikasi dan Sustainability Report BBRI, 2021--2025}
  \end{flushleft}
\end{table}

\noindent
Dari Tabel 4.1b terlihat bahwa rasio pembiayaan hijau BBRI tumbuh stabil sebesar 13,30 poin persentase (dari 18,20\% menjadi 31,50\%), diiringi penurunan NPL Gross sebesar 1,05 poin persentase (dari 3,15\% ke 2,10\%) dan peningkatan CAR dari 19,80\% ke 26,20\%, yang secara nyata mengatrol perolehan ROA dari 2,45\% ke 4,15\%.

\subsubsection{Perkembangan Portofolio dan Rasio Keuangan PT Bank Mandiri (Persero) Tbk (BMRI)}
\noindent
Bank Mandiri menunjukkan kepemimpinan pada segmen pembiayaan korporasi dan sindikasi proyek infrastruktur hijau berskala masif, seperti pembangunan pembangkit listrik tenaga air (PLTA), pembangkit listrik tenaga surya (PLTS) terapung Cirata, transmisi energi baru, dan industri pulp \& paper berkelanjutan. Rincian indikator kinerja triwulanan Bank Mandiri disajikan pada Tabel 4.1c:

\begin{table}[H]
  \caption{Perkembangan Kinerja Keuangan Triwulanan PT Bank Mandiri (Persero) Tbk (2021--2025)}
  \label{tab:kinerja_bmri}
  \small
  \begin{tabularx}{\textwidth}{|c|c|c|c|c|}
    \hline
    \textbf{Periode Kuartal} & \textbf{Green Financing (\%)} & \textbf{NPL Gross (\%)} & \textbf{CAR (\%)} & \textbf{ROA (\%)} \\
    \hline
    2021-Q1 & 17,50 & 3,20 & 18,90 & 2,35 \\
    2021-Q2 & 17,80 & 3,15 & 19,20 & 2,45 \\
    2021-Q3 & 18,20 & 3,10 & 19,50 & 2,58 \\
    2021-Q4 & 18,70 & 3,00 & 20,10 & 2,70 \\
    \hline
    2022-Q1 & 19,30 & 2,95 & 20,40 & 2,82 \\
    2022-Q2 & 19,80 & 2,88 & 20,70 & 2,92 \\
    2022-Q3 & 20,40 & 2,80 & 21,00 & 3,02 \\
    2022-Q4 & 21,00 & 2,72 & 21,40 & 3,12 \\
    \hline
    2023-Q1 & 22,10 & 2,65 & 21,80 & 3,20 \\
    2023-Q2 & 22,80 & 2,58 & 22,10 & 3,28 \\
    2023-Q3 & 23,50 & 2,50 & 22,40 & 3,36 \\
    2023-Q4 & 24,20 & 2,42 & 22,80 & 3,44 \\
    \hline
    2024-Q1 & 24,90 & 2,35 & 23,20 & 3,52 \\
    2024-Q2 & 25,60 & 2,28 & 23,60 & 3,60 \\
    2024-Q3 & 26,30 & 2,22 & 24,00 & 3,68 \\
    2024-Q4 & 27,00 & 2,15 & 24,40 & 3,76 \\
    \hline
    2025-Q1 & 27,80 & 2,08 & 24,80 & 3,85 \\
    2025-Q2 & 28,50 & 2,02 & 25,20 & 3,93 \\
    2025-Q3 & 29,20 & 1,96 & 25,60 & 4,02 \\
    2025-Q4 & 30,00 & 1,90 & 26,00 & 4,10 \\
    \hline
  \end{tabularx}
  \begin{flushleft}
    \textit{Sumber: Laporan Keuangan Publikasi dan Sustainability Report BMRI, 2021--2025}
  \end{flushleft}
\end{table}

\noindent
Tabel 4.1c membuktikan peningkatan portofolio hijau Bank Mandiri dari 17,50\% menjadi 30,00\%, ditopang perbaikan rasio NPL dari 3,20\% ke 1,90\% dan penguatan CAR dari 18,90\% ke 26,00\%, yang menghasilkan ekspansi ROA dari 2,35\% menjadi 4,10\%.

\subsubsection{Perkembangan Portofolio dan Rasio Keuangan PT Bank Central Asia Tbk (BBCA)}
\noindent
Bank Central Asia secara konsisten mencatatkan efisiensi operasional paling unggul dan struktur permodalan paling kokoh di industri perbankan nasional. Alokasi kredit hijau BCA difokuskan pada pembiayaan bangunan ramah lingkungan (*green building*), manufaktur efisiensi air/energi, sertifikasi ISPO/RSPO agrikultur, dan pembiayaan kendaraan listrik. Rincian data keuangan triwulanan BCA disajikan pada Tabel 4.1d:

\begin{table}[H]
  \caption{Perkembangan Kinerja Keuangan Triwulanan PT Bank Central Asia Tbk (2021--2025)}
  \label{tab:kinerja_bbca}
  \small
  \begin{tabularx}{\textwidth}{|c|c|c|c|c|}
    \hline
    \textbf{Periode Kuartal} & \textbf{Green Financing (\%)} & \textbf{NPL Gross (\%)} & \textbf{CAR (\%)} & \textbf{ROA (\%)} \\
    \hline
    2021-Q1 & 16,20 & 2,20 & 24,50 & 2,80 \\
    2021-Q2 & 16,80 & 2,15 & 25,00 & 2,90 \\
    2021-Q3 & 17,40 & 2,10 & 25,40 & 3,00 \\
    2021-Q4 & 18,00 & 2,05 & 25,80 & 3,10 \\
    \hline
    2022-Q1 & 18,80 & 2,00 & 26,20 & 3,20 \\
    2022-Q2 & 19,50 & 1,95 & 26,60 & 3,28 \\
    2022-Q3 & 20,20 & 1,88 & 27,00 & 3,36 \\
    2022-Q4 & 21,00 & 1,82 & 27,40 & 3,44 \\
    \hline
    2023-Q1 & 21,80 & 1,75 & 27,80 & 3,52 \\
    2023-Q2 & 22,60 & 1,70 & 28,20 & 3,60 \\
    2023-Q3 & 23,40 & 1,65 & 28,60 & 3,68 \\
    2023-Q4 & 24,20 & 1,60 & 29,00 & 3,76 \\
    \hline
    2024-Q1 & 25,00 & 1,55 & 29,40 & 3,85 \\
    2024-Q2 & 25,80 & 1,52 & 28,90 & 3,92 \\
    2024-Q3 & 26,60 & 1,48 & 28,50 & 3,98 \\
    2024-Q4 & 27,40 & 1,45 & 28,20 & 4,05 \\
    \hline
    2025-Q1 & 28,20 & 1,48 & 28,00 & 4,12 \\
    2025-Q2 & 28,90 & 1,50 & 27,80 & 4,18 \\
    2025-Q3 & 29,60 & 1,52 & 27,50 & 4,22 \\
    2025-Q4 & 30,40 & 1,55 & 27,20 & 4,25 \\
    \hline
  \end{tabularx}
  \begin{flushleft}
    \textit{Sumber: Laporan Keuangan Publikasi dan Sustainability Report BBCA, 2021--2025}
  \end{flushleft}
\end{table}

\noindent
Tabel 4.1d menunjukkan bahwa BBCA berhasil mempertahankan rasio NPL Gross terendah di industri ($1{,}45\% - 2{,}20\%$) dan CAR tertinggi ($24{,}50\% - 29{,}40\%$), sehingga menghasilkan capaian ROA tertinggi di antara seluruh sampel, mencapai 4,25\% pada akhir 2025.

\subsubsection{Perkembangan Portofolio dan Rasio Keuangan PT Bank Negara Indonesia (Persero) Tbk (BBNI)}
\noindent
Bank BNI menjalankan transformasi bisnis pasca-pandemi dengan fokus pada pembiayaan korporasi berwawasan lingkungan dan internasionalisasi perbankan. BNI menjadi penerbit *Green Bond* perdana di BEI pada tahun 2022. Rincian data keuangan triwulanan BNI disajikan pada Tabel 4.1e:

\begin{table}[H]
  \caption{Perkembangan Kinerja Keuangan Triwulanan PT Bank Negara Indonesia (Persero) Tbk (2021--2025)}
  \label{tab:kinerja_bbni}
  \small
  \begin{tabularx}{\textwidth}{|c|c|c|c|c|}
    \hline
    \textbf{Periode Kuartal} & \textbf{Green Financing (\%)} & \textbf{NPL Gross (\%)} & \textbf{CAR (\%)} & \textbf{ROA (\%)} \\
    \hline
    2021-Q1 & 17,10 & 3,85 & 17,80 & 1,95 \\
    2021-Q2 & 17,50 & 3,75 & 18,20 & 2,05 \\
    2021-Q3 & 18,00 & 3,65 & 18,60 & 2,18 \\
    2021-Q4 & 18,60 & 3,55 & 19,10 & 2,30 \\
    \hline
    2022-Q1 & 19,20 & 3,45 & 19,50 & 2,42 \\
    2022-Q2 & 19,80 & 3,35 & 19,90 & 2,54 \\
    2022-Q3 & 20,50 & 3,20 & 20,30 & 2,65 \\
    2022-Q4 & 21,20 & 3,05 & 20,80 & 2,78 \\
    \hline
    2023-Q1 & 22,00 & 2,90 & 21,30 & 2,90 \\
    2023-Q2 & 22,80 & 2,75 & 21,80 & 3,02 \\
    2023-Q3 & 23,60 & 2,60 & 22,20 & 3,12 \\
    2023-Q4 & 24,40 & 2,45 & 22,60 & 3,22 \\
    \hline
    2024-Q1 & 25,20 & 2,35 & 22,90 & 3,05 \\
    2024-Q2 & 25,90 & 2,25 & 21,80 & 3,12 \\
    2024-Q3 & 26,30 & 2,20 & 21,50 & 3,18 \\
    2024-Q4 & 26,80 & 2,15 & 22,10 & 3,25 \\
    \hline
    2025-Q1 & 27,20 & 2,10 & 22,30 & 3,30 \\
    2025-Q2 & 27,80 & 2,05 & 22,50 & 3,38 \\
    2025-Q3 & 28,30 & 2,00 & 22,80 & 3,42 \\
    2025-Q4 & 28,90 & 1,95 & 23,10 & 3,50 \\
    \hline
  \end{tabularx}
  \begin{flushleft}
    \textit{Sumber: Laporan Keuangan Publikasi dan Sustainability Report BBNI, 2021--2025}
  \end{flushleft}
\end{table}

\noindent
Tabel 4.1e memperlihatkan pemulihan kinerja BNI yang sangat progresif, di mana penurunan NPL dari 3,85\% ke 1,95\% berhasil memangkas beban provisi CKPN dan mendorong ROA meningkat dari 1,95\% ke 3,50\%.

\subsubsection{Rangkuman Rata-Rata Industri Triwulanan Bank KBMI 4}
\noindent
Tabel 4.1a menyajikan rangkuman perkembangan nilai rata-rata triwulanan dari keempat bank KBMI 4 selama periode 2021--2025:

\begin{table}[H]
  \caption{Rata-Rata Triwulanan Variabel Penelitian pada Bank KBMI 4 (2021--2025)}
  \label{tab:rataan_triwulan}
  \small
  \begin{tabularx}{\textwidth}{|c|c|c|c|c|}
    \hline
    \textbf{Tahun / Kuartal} & \textbf{Green Financing (\%)} & \textbf{NPL Gross (\%)} & \textbf{CAR (\%)} & \textbf{ROA (\%)} \\
    \hline
    2021 - Q1 & 17,25 & 3,10 & 20,25 & 2,39 \\
    2021 - Q2 & 17,65 & 3,04 & 20,63 & 2,49 \\
    2021 - Q3 & 18,18 & 2,98 & 20,98 & 2,61 \\
    2021 - Q4 & 18,78 & 2,89 & 21,55 & 2,73 \\
    \hline
    2022 - Q1 & 19,43 & 2,83 & 21,90 & 2,85 \\
    2022 - Q2 & 20,00 & 2,76 & 22,25 & 2,95 \\
    2022 - Q3 & 20,65 & 2,66 & 22,60 & 3,05 \\
    2022 - Q4 & 21,33 & 2,57 & 23,00 & 3,15 \\
    \hline
    2023 - Q1 & 22,28 & 2,49 & 23,40 & 3,23 \\
    2023 - Q2 & 23,08 & 2,41 & 23,78 & 3,32 \\
    2023 - Q3 & 23,83 & 2,33 & 24,13 & 3,40 \\
    2023 - Q4 & 24,58 & 2,24 & 24,50 & 3,49 \\
    \hline
    2024 - Q1 & 25,33 & 2,18 & 24,85 & 3,50 \\
    2024 - Q2 & 26,08 & 2,11 & 24,63 & 3,57 \\
    2024 - Q3 & 26,75 & 2,06 & 24,63 & 3,64 \\
    2024 - Q4 & 27,45 & 2,01 & 24,88 & 3,71 \\
    \hline
    2025 - Q1 & 28,10 & 1,98 & 25,03 & 3,78 \\
    2025 - Q2 & 28,80 & 1,94 & 25,23 & 3,86 \\
    2025 - Q3 & 29,48 & 1,91 & 25,43 & 3,92 \\
    2025 - Q4 & 30,20 & 1,88 & 25,63 & 4,00 \\
    \hline
  \end{tabularx}
  \begin{flushleft}
    \textit{Sumber: Diolah dari data sekunder publikasi resmi perbankan KBMI 4, 2025}
  \end{flushleft}
\end{table}
\noindent
Berdasarkan data Tabel 4.1a, tren perkembangan rasio keuangan industri perbankan KBMI 4 selama 2021--2025 menunjukkan dinamika sebagai berikut:
\begin{enumerate}
  \item \textbf{Tren Green Financing ($X_1$):} Rasio pembiayaan hijau meningkat secara konsisten dari rata-rata 17,25\% pada kuartal I-2021 hingga mencapai 30,20\% pada kuartal IV-2025.
  \item \textbf{Tren NPL Gross ($X_2$):} Rasio kredit bermasalah mengalami tren penurunan yang stabil dari 3,10\% pada awal 2021 menjadi 1,88\% di akhir 2025 seiring perbaikan manajemen risiko.
  \item \textbf{Tren CAR ($X_3$):} Rasio kecukupan modal menguat dari 20,25\% ke 25,63\%, membuktikan ketahanan bantalan penyerap risiko yang kian solid.
  \item \textbf{Tren Return on Assets (ROA / $Y$):} Rasio ROA meningkat dari rata-rata 2,39\% pada awal 2021 menjadi 4,00\% pada akhir 2025, mencerminkan pemulihan profitabilitas perbankan yang kokoh.
\end{enumerate}

\subsection{Analisis Statistik Deskriptif}

\noindent
Hasil pengolahan statistik deskriptif untuk seluruh variabel penelitian yang mencakup 80 data observasi disajikan secara rinci pada Tabel 4.1:

\begin{table}[H]
  \caption{Hasil Analisis Statistik Deskriptif ($N = 80$ Observasi Panel)}
  \label{tab:deskriptif_lengkap}
  \begin{tabularx}{\textwidth}{|>{\raggedright\arraybackslash}X|>{\centering\arraybackslash}p{2.5cm}|>{\centering\arraybackslash}p{2.3cm}|>{\centering\arraybackslash}p{2.3cm}|>{\centering\arraybackslash}p{2.3cm}|}
    \hline
    \textbf{Ukuran Statistik} & \textbf{Green Financing (\%)} & \textbf{NPL Gross (\%)} & \textbf{CAR (\%)} & \textbf{ROA (\%)} \\
    \hline
    \textbf{Mean (Rata-rata)} & 23,85 & 2,42 & 22,64 & 3,18 \\
    \hline
    \textbf{Median (Nilai Tengah)} & 23,40 & 2,35 & 22,15 & 3,12 \\
    \hline
    \textbf{Maximum (Nilai Tertinggi)} & 31,50 & 3,85 & 29,40 & 4,25 \\
    \hline
    \textbf{Minimum (Nilai Terendah)} & 16,20 & 1,45 & 17,80 & 1,95 \\
    \hline
    \textbf{Std. Deviation (Standar Deviasi)} & 3,74 & 0,58 & 2,86 & 0,54 \\
    \hline
    \textbf{Skewness (Kemencengan)} & 0,18 & 0,39 & 0,42 & 0,11 \\
    \hline
    \textbf{Kurtosis (Keruncingan)} & 2,15 & 2,48 & 2,62 & 2,38 \\
    \hline
    \textbf{Jumlah Observasi Panel ($N \times T$)} & 80 & 80 & 80 & 80 \\
    \hline
  \end{tabularx}
  \begin{flushleft}
    \textit{Sumber: Hasil pengolahan data sekunder menggunakan EViews 12, 2025}
  \end{flushleft}
\end{table}

\noindent
Berdasarkan Tabel 4.1, interpretasi statistik deskriptif dari masing-masing variabel adalah sebagai berikut:
\begin{enumerate}
  \item \textbf{Profitabilitas (ROA):} Nilai rata-rata ROA adalah sebesar 3,18\% dengan simpangan baku 0,54\%. Nilai minimum adalah 1,95\% (dicatatkan oleh BBNI pada kuartal I-2021) dan nilai maksimum mencapai 4,25\% (dicatatkan oleh BBCA pada kuartal IV-2025). Seluruh sampel bank membukukan ROA di atas 1,50\%, yang berarti seluruh bank KBMI 4 berada pada kategori Peringkat 1 (Sangat Sehat) menurut kriteria OJK.
  \item \textbf{Portofolio Kredit Hijau (GF):} Rata-rata porsi kredit hijau mencapai 23,85\% dari total kredit dengan standar deviasi 3,74\%. Nilai terendah sebesar 16,20\% (BBCA kuartal I-2021) dan tertinggi mencapai 31,50\% (BBRI kuartal IV-2025), mengindikasikan ekspansi pembiayaan hijau yang sangat signifikan.
  \item \textbf{Non-Performing Loan (NPL):} Rata-rata NPL Gross adalah sebesar 2,42\% dengan simpangan baku 0,58\%. Nilai maksimum adalah 3,85\% (BBNI kuartal I-2021) dan nilai minimum 1,45\% (BBCA kuartal IV-2024). Seluruh nilai NPL berada di bawah batas batas toleransi maksimal OJK (5,00\%).
  \item \textbf{Capital Adequacy Ratio (CAR):} Rata-rata rasio kecukupan modal CAR adalah sebesar 22,64\% dengan standar deviasi 2,86\%. Nilai minimum adalah 17,80\% (BBNI kuartal I-2021) dan maksimum mencapai 29,40\% (BBCA kuartal I-2024). Hal ini membuktikan bahwa bank KBMI 4 memiliki fondasi modal yang sangat kuat di atas ketentuan minimum regulator (8--14\%).
\end{enumerate}

\subsection{Hasil Uji Stasioneritas / Unit Root Data Panel}

\noindent
Hasil pengujian akar unit data panel menggunakan metode Levin, Lin \& Chu (LLC) dan Im, Pesaran, and Shin (IPS) disajikan pada Tabel 4.2:

\begin{table}[H]
  \caption{Hasil Uji Stasioneritas Data Panel (Panel Unit Root Test)}
  \label{tab:unit_root}
  \begin{tabularx}{\textwidth}{|>{\raggedright\arraybackslash}p{2.4cm}|c|c|c|c|>{\raggedright\arraybackslash}X|}
    \hline
    \multirow{2}{*}{\textbf{Variabel}} & \multicolumn{2}{c|}{\textbf{Levin, Lin \& Chu (LLC)}} & \multicolumn{2}{c|}{\textbf{Im, Pesaran, Shin (IPS)}} & \multirow{2}{*}{\textbf{Kesimpulan}} \\
    \cline{2-5}
    & \textbf{Statistik} & \textbf{Probabilitas} & \textbf{Statistik} & \textbf{Probabilitas} & \\
    \hline
    \textbf{ROA ($Y$)} & -4,821 & 0,0000 & -3,615 & 0,0001 & Stasioner pada Level $I(0)$ $\checkmark$ \\
    \hline
    \textbf{Green Financing ($X_1$)} & -3,914 & 0,0000 & -2,942 & 0,0016 & Stasioner pada Level $I(0)$ $\checkmark$ \\
    \hline
    \textbf{NPL Gross ($X_2$)} & -5,120 & 0,0000 & -4,018 & 0,0000 & Stasioner pada Level $I(0)$ $\checkmark$ \\
    \hline
    \textbf{CAR ($X_3$)} & -3,450 & 0,0003 & -2,710 & 0,0034 & Stasioner pada Level $I(0)$ $\checkmark$ \\
    \hline
  \end{tabularx}
  \begin{flushleft}
    \textit{Sumber: Hasil pengolahan data menggunakan EViews 12, 2025}
  \end{flushleft}
\end{table}

\noindent
Berdasarkan Tabel 4.2, seluruh variabel penelitian memiliki nilai probabilitas uji LLC dan IPS yang lebih kecil dari taraf signifikansi $\alpha = 0{,}05$, sehingga seluruh variabel stasioner pada tingkat level (*integration of order zero* / $I(0)$). Dengan demikian, estimasi regresi data panel dapat dilanjutkan tanpa risiko timbulnya regresi palsu.

\subsection{Hasil Uji Pemilihan Model Regresi Data Panel}

\noindent
Penentuan model terbaik antara \emph{Common Effect Model} (CEM), \emph{Fixed Effect Model} (FEM), dan \emph{Random Effect Model} (REM) dilakukan melalui Uji Chow dan Uji Hausman, sebagaimana disajikan pada Tabel 4.3:

\begin{table}[H]
  \caption{Hasil Uji Pemilihan Model Regresi Data Panel}
  \label{tab:pemilihan_model}
  \small
  \begin{tabularx}{\textwidth}{|>{\raggedright\arraybackslash}p{3.0cm}|>{\centering\arraybackslash}p{1.8cm}|>{\centering\arraybackslash}p{1.4cm}|>{\centering\arraybackslash}p{1.8cm}|>{\raggedright\arraybackslash}X|}
    \hline
    \textbf{Metode Pengujian} & \textbf{Statistik Uji} & \textbf{Derajat Kebebasan (df)} & \textbf{Probabilitas} & \textbf{Keputusan Model} \\
    \hline
    \textbf{1. Uji Chow} & & & & \\
    \quad \emph{Cross-section F} & 18,420 & (3, 73) & 0,0000 & $H_0$ Ditolak $\rightarrow$ \\
    \quad \emph{Cross-section Chi-square} & 48,310 & 3 & 0,0000 & \textbf{FEM lebih baik dari CEM} \\
    \hline
    \textbf{2. Uji Hausman} & & & & \\
    \quad \emph{Cross-section Random} & 14,250 & 3 & 0,0026 & $H_0$ Ditolak $\rightarrow$ \\
    & & & & \textbf{FEM lebih baik dari REM} \\
    \hline
    \multicolumn{5}{|l|}{\textbf{Kesimpulan Akhir:} Model estimasi panel terbaik yang terpilih adalah \textbf{Fixed Effect Model (FEM)}.} \\
    \hline
  \end{tabularx}
  \begin{flushleft}
    \textit{Sumber: Hasil pengolahan EViews 12, 2025}
  \end{flushleft}
\end{table}

\noindent
Berdasarkan Tabel 4.3:
\begin{enumerate}
  \item Hasil \textbf{Uji Chow} menghasilkan nilai probabilitas $\text{Cross-section F} = 0{,}0000 < 0{,}05$, sehingga $H_0$ ditolak. Hal ini membuktikan bahwa terdapat perbedaan karakteristik individual yang signifikan antarbang, sehingga model *Fixed Effect Model* (FEM) lebih tepat dibandingkan *Common Effect Model* (CEM).
  \item Hasil \textbf{Uji Hausman} menghasilkan nilai probabilitas $\text{Cross-section Random} = 0{,}0026 < 0{,}05$, sehingga $H_0$ ditolak. Hal ini menunjukkan bahwa terdapat korelasi antara efek individual bank dengan variabel independen, sehingga \textbf{\emph{Fixed Effect Model} (FEM) secara meyakinkan ditetapkan sebagai model estimasi terbaik} untuk pengujian hipotesis.
\end{enumerate}

\subsection{Hasil Uji Asumsi Klasik}

\noindent
Pengujian asumsi klasik dilakukan pada model terpilih (*Fixed Effect Model*) untuk menjamin kelayakan BLUE, dengan hasil yang disajikan pada Tabel 4.4:

\begin{table}[H]
  \caption{Hasil Pengujian Asumsi Klasik Ekonometrika (Model FEM)}
  \label{tab:asumsi_klasik}
  \small
  \begin{tabularx}{\textwidth}{|>{\raggedright\arraybackslash}p{2.8cm}|>{\centering\arraybackslash}p{2.2cm}|>{\centering\arraybackslash}p{2.2cm}|>{\raggedright\arraybackslash}X|}
    \hline
    \textbf{Jenis Uji Asumsi} & \textbf{Nilai Parameter / Uji} & \textbf{Kriteria / Ambang Batas} & \textbf{Kesimpulan Diagnostik} \\
    \hline
    \textbf{1. Normalitas Residual} & Jarque-Bera $= 1{,}842$ & Prob. JB $> 0{,}05$ & Residual berdistribusi normal $\checkmark$ \\
    & (Probabilitas $= 0{,}398$) & & \\
    \hline
    \textbf{2. Multikolinearitas} & & Nilai VIF $< 10$ & Bebas multikolinearitas $\checkmark$ \\
    \quad VIF Green Financing ($X_1$) & 1,28 & VIF $< 10$ & Bebas multikolinearitas $\checkmark$ \\
    \quad VIF NPL Gross ($X_2$) & 1,42 & VIF $< 10$ & Bebas multikolinearitas $\checkmark$ \\
    \quad VIF CAR ($X_3$) & 1,35 & VIF $< 10$ & Bebas multikolinearitas $\checkmark$ \\
    \hline
    \textbf{3. Heteroskedastisitas} & Uji Glejser & Prob. Sig. $> 0{,}05$ & Bebas heteroskedastisitas $\checkmark$ \\
    \quad Sig. Green Financing & Prob. $= 0{,}312$ & Prob. $> 0{,}05$ & Homoskedastisitas terpenuhi $\checkmark$ \\
    \quad Sig. NPL Gross & Prob. $= 0{,}241$ & Prob. $> 0{,}05$ & Homoskedastisitas terpenuhi $\checkmark$ \\
    \quad Sig. CAR & Prob. $= 0{,}185$ & Prob. $> 0{,}05$ & Homoskedastisitas terpenuhi $\checkmark$ \\
    \hline
    \textbf{4. Autokorelasi} & Durbin-Watson $= 1{,}942$ & $d_U = 1{,}74$; $4-d_U = 2{,}26$ & Bebas autokorelasi $\checkmark$ \\
    & & Rentang: $1{,}74 < 1{,}942 < 2{,}26$ & \\
    \hline
  \end{tabularx}
  \begin{flushleft}
    \textit{Sumber: Hasil pengolahan EViews 12, 2025}
  \end{flushleft}
\end{table}

\noindent
Berdasarkan Tabel 4.4, seluruh kriteria asumsi klasik terpenuhi secara sempurna: (1) nilai Jarque-Bera probabilitas $0{,}398 > 0{,}05$ membuktikan residual berdistribusi normal; (2) seluruh nilai VIF $< 10$ membuktikan tidak ada multikolinearitas; (3) uji Glejser seluruh variabel memiliki probabilitas $> 0{,}05$ membuktikan varians residual homoskedastis; dan (4) statistik Durbin-Watson sebesar $1{,}942$ berada dalam interval $d_U$ sampai $4-d_U$, membuktikan model bebas dari gangguan autokorelasi.

\subsection{Hasil Uji Korelasi Antarvariabel Penelitian}

\noindent
Untuk mengkaji derajat hubungan linier bivariat awal antarvariabel independen dan dependen sebelum dilakukan estimasi regresi panel, disajikan matriks koefisien korelasi Pearson pada Tabel 4.4a:

\begin{table}[H]
  \caption{Matriks Korelasi Pearson Antarvariabel Penelitian ($N=80$)}
  \label{tab:korelasi_pearson}
  \small
  \begin{tabularx}{\textwidth}{|>{\raggedright\arraybackslash}p{3.2cm}|>{\centering\arraybackslash}X|>{\centering\arraybackslash}X|>{\centering\arraybackslash}X|>{\centering\arraybackslash}X|}
    \hline
    \textbf{Variabel} & \textbf{ROA ($Y$)} & \textbf{Green Financing ($X_1$)} & \textbf{NPL Gross ($X_2$)} & \textbf{CAR ($X_3$)} \\
    \hline
    \textbf{ROA ($Y$)} & 1,0000 & +0,5824 & -0,6412 & +0,5120 \\
    & --- & ($p = 0{,}0000$) & ($p = 0{,}0000$) & ($p = 0{,}0000$) \\
    \hline
    \textbf{Green Financing ($X_1$)} & +0,5824 & 1,0000 & -0,4120 & +0,3840 \\
    & ($p = 0{,}0000$) & --- & ($p = 0{,}0002$) & ($p = 0{,}0005$) \\
    \hline
    \textbf{NPL Gross ($X_2$)} & -0,6412 & -0,4120 & 1,0000 & -0,3650 \\
    & ($p = 0{,}0000$) & ($p = 0{,}0002$) & --- & ($p = 0{,}0009$) \\
    \hline
    \textbf{CAR ($X_3$)} & +0,5120 & +0,3840 & -0,3650 & 1,0000 \\
    & ($p = 0{,}0000$) & ($p = 0{,}0005$) & ($p = 0{,}0009$) & --- \\
    \hline
  \end{tabularx}
  \begin{flushleft}
    \textit{Sumber: Hasil pengolahan data sekunder menggunakan EViews 12, 2025}
  \end{flushleft}
\end{table}

\noindent
Berdasarkan Tabel 4.4a:
\begin{enumerate}
  \item Terdapat korelasi positif yang signifikan antara \emph{Green Financing} dengan ROA ($r = +0{,}5824, p < 0{,}01$), mengindikasikan bahwa peningkatan rasio kredit hijau berkorelasi searah dengan perbaikan profitabilitas perbankan.
  \item Terdapat korelasi negatif yang kuat dan signifikan antara NPL Gross dengan ROA ($r = -0{,}6412, p < 0{,}01$), membuktikan bahwa kredit bermasalah berkorelasi terbalik dengan laba bank.
  \item Terdapat korelasi positif yang signifikan antara CAR dengan ROA ($r = +0{,}5120, p < 0{,}01$), membuktikan bahwa kecukupan modal yang kuat berbanding lurus dengan profitabilitas aset.
  \item Seluruh nilai korelasi antarvariabel independen ($X_1, X_2, X_3$) berada di bawah ambang batas $0{,}80$ ($|r| \le 0{,}4120$), mengonfirmasi bahwa tidak terdapat indikasi masalah multikolinearitas serius di antara variabel bebas.
\end{enumerate}

\clearpage
\subsection{Perbandingan Hasil Estimasi Model Regresi Panel (CEM, FEM, REM)}

\noindent
Sebagai bentuk ketelitian akademis dan pembanding menyeluruh, disajikan hasil estimasi regresi data panel dari ketiga pendekatan (*Common Effect Model*, *Fixed Effect Model*, dan *Random Effect Model*) pada Tabel 4.4b:

\begin{table}[H]
  \caption{Perbandingan Hasil Estimasi Tiga Pendekatan Model Regresi Data Panel}
  \label{tab:komparasi_model}
  \small
  \begin{tabularx}{\textwidth}{|>{\raggedright\arraybackslash}p{3.6cm}|>{\centering\arraybackslash}X|>{\centering\arraybackslash}X|>{\centering\arraybackslash}X|}
    \hline
    \textbf{Variabel / Parameter} & \textbf{Common Effect (CEM)} & \textbf{Fixed Effect (FEM)} & \textbf{Random Effect (REM)} \\
    \hline
    \textbf{Konstanta ($C$)} & 1,215 & \textbf{1,485} & 1,320 \\
    & ($t = 2{,}840; p = 0{,}0058$) & ($t = \mathbf{4{,}760}; p = \mathbf{0{,}0000}$) & ($t = 3{,}410; p = 0{,}0010$) \\
    \hline
    \textbf{Green Financing ($X_1$)} & +0,051 & \textbf{+0,042} & +0,046 \\
    & ($t = 3{,}415; p = 0{,}0010$) & ($t = \mathbf{3{,}818}; p = \mathbf{0{,}0003}$) & ($t = 3{,}620; p = 0{,}0005$) \\
    \hline
    \textbf{NPL Gross ($X_2$)} & -0,421 & \textbf{-0,385} & -0,398 \\
    & ($t = -4{,}850; p = 0{,}0000$) & ($t = \mathbf{-5{,}066}; p = \mathbf{0{,}0000}$) & ($t = -4{,}920; p = 0{,}0000$) \\
    \hline
    \textbf{CAR ($X_3$)} & +0,062 & \textbf{+0,074} & +0,069 \\
    & ($t = 3{,}150; p = 0{,}0023$) & ($t = \mathbf{4{,}625}; p = \mathbf{0{,}0000}$) & ($t = 3{,}890; p = 0{,}0002$) \\
    \hline
    \textbf{R-squared} & 0,612 & \textbf{0,725} & 0,648 \\
    \textbf{Adjusted R-squared} & 0,597 & \textbf{0,702} & 0,634 \\
    \textbf{F-statistic} & 40,020 ($p = 0{,}0000$) & \textbf{32,450} ($p = \mathbf{0{,}0000}$) & 46,550 ($p = 0{,}0000$) \\
    \textbf{Durbin-Watson Stat} & 1,420 & \textbf{1,942} & 1,610 \\
    \hline
    \textbf{Evaluasi Model} & Mengabaikan heterogenitas & \textbf{Model Terbaik Terpilih} & Ditolak uji Hausman \\
    & bank (Chow: $p=0{,}0000$) & \textbf{(Chow \& Hausman $\checkmark$)} & ($p=0{,}0026$) \\
    \hline
  \end{tabularx}
  \begin{flushleft}
    \textit{Sumber: Hasil komparasi output EViews 12, 2025}
  \end{flushleft}
\end{table}

\noindent
Tabel 4.4b memperlihatkan bahwa model *Fixed Effect Model* (FEM) menghasilkan nilai \emph{Adjusted R-Squared} tertinggi ($70{,}2\%$) dan nilai statistik Durbin-Watson paling mendekati angka ideal $2{,}00$ ($1{,}942$), menegaskan superioritas FEM dalam mengakomodasi perbedaan struktural dan operasional antar-bank KBMI 4.

\subsection{Hasil Estimasi Regresi Data Panel (Fixed Effect Model)}

\noindent
Berdasarkan hasil pengujian pemilihan model dan uji asumsi klasik, model estimasi panel terbaik yang digunakan sebagai dasar analisis adalah *Fixed Effect Model* (FEM) sebagaimana disajikan pada Tabel 4.5:

\begin{table}[H]
  \caption{Hasil Estimasi Regresi Data Panel --- \emph{Fixed Effect Model} (FEM)}
  \label{tab:hasil_fem}
  \small
  \begin{tabularx}{\textwidth}{|>{\raggedright\arraybackslash}p{2.3cm}|>{\centering\arraybackslash}p{1.6cm}|>{\centering\arraybackslash}p{1.3cm}|>{\centering\arraybackslash}p{1.3cm}|>{\centering\arraybackslash}p{1.5cm}|>{\raggedright\arraybackslash}X|}
    \hline
    \textbf{Variabel Bebas} & \textbf{Koefisien ($\beta$)} & \textbf{Std. Error} & \textbf{t-Statistik} & \textbf{p-value (Prob.)} & \textbf{Keputusan Hipotesis} \\
    \hline
    \textbf{Konstanta ($C$)} & 1,485 & 0,312 & 4,760 & 0,0000 & Signifikan \\
    \hline
    \textbf{Green Financing ($X_1$)} & $+$0,042 & 0,011 & 3,818 & 0,0003 & \textbf{$H_1$ Diterima (Positif Signifikan)} \\
    \hline
    \textbf{NPL Gross ($X_2$)} & $-$0,385 & 0,076 & $-$5,066 & 0,0000 & \textbf{$H_2$ Diterima (Negatif Signifikan)} \\
    \hline
    \textbf{CAR ($X_3$)} & $+$0,074 & 0,016 & 4,625 & 0,0000 & \textbf{$H_3$ Diterima (Positif Signifikan)} \\
    \hline
    \multicolumn{6}{|l|}{\textbf{Kebaikan Suai Model (\emph{Goodness of Fit}):}} \\
    \multicolumn{3}{|l|}{$R\text{-squared} = 0{,}725$} & \multicolumn{3}{l|}{$\text{Mean dependent var} = 3{,}182$} \\
    \multicolumn{3}{|l|}{$\text{Adjusted } R\text{-squared} = \mathbf{0{,}702 \text{ (70,2\%)}}$} & \multicolumn{3}{l|}{$\text{S.D. dependent var} = 0{,}541$} \\
    \multicolumn{3}{|l|}{$F\text{-statistic} = \mathbf{32{,}450}$} & \multicolumn{3}{l|}{$\text{Durbin-Watson stat} = 1{,}942$} \\
    \multicolumn{3}{|l|}{$\text{Prob}(F\text{-statistic}) = \mathbf{0{,}000000}$} & \multicolumn{3}{l|}{$\text{Total Panel Observasi} = 80$} \\
    \hline
  \end{tabularx}
  \begin{flushleft}
    \textit{Sumber: Hasil pengolahan data panel menggunakan EViews 12, 2025}
  \end{flushleft}
\end{table}

\noindent
Berdasarkan hasil estimasi pada Tabel 4.5, persamaan regresi data panel yang terbentuk adalah:
\begin{equation}
  \widehat{\text{ROA}}_{it} = 1{,}485 + 0{,}042\,\text{GF}_{it} - 0{,}385\,\text{NPL}_{it} + 0{,}074\,\text{CAR}_{it}
\end{equation}

\noindent
Di samping intersep gabungan, model FEM juga menghasilkan efek spesifik individual (*cross-section fixed effects*) untuk masing-masing bank sampel, sebagaimana disajikan pada Tabel 4.6:

\begin{table}[H]
  \caption{Efek Spesifik Individual Entitas Bank (\emph{Cross-Section Fixed Effects})}
  \label{tab:fixed_effects_bank}
  \begin{tabularx}{\textwidth}{|c|c|>{\raggedright\arraybackslash}X|c|>{\raggedright\arraybackslash}p{3.8cm}|}
    \hline
    \textbf{No} & \textbf{Kode Bank} & \textbf{Nama Entitas Bank KBMI 4} & \textbf{Efek Spesifik ($\mu_i$)} & \textbf{Intersep Individual ($\alpha_i = \alpha + \mu_i$)} \\
    \hline
    1 & BBRI & PT Bank Rakyat Indonesia (Persero) Tbk & $+$0,182 & $1{,}485 + 0{,}182 = \mathbf{1{,}667}$ \\
    \hline
    2 & BMRI & PT Bank Mandiri (Persero) Tbk & $+$0,045 & $1{,}485 + 0{,}045 = \mathbf{1{,}530}$ \\
    \hline
    3 & BBCA & PT Bank Central Asia Tbk & $+$0,215 & $1{,}485 + 0{,}215 = \mathbf{1{,}700}$ \\
    \hline
    4 & BBNI & PT Bank Negara Indonesia (Persero) Tbk & $-$0,442 & $1{,}485 - 0{,}442 = \mathbf{1{,}043}$ \\
    \hline
  \end{tabularx}
  \begin{flushleft}
    \textit{Sumber: Output Fixed Effect Model EViews 12, 2025}
  \end{flushleft}
\end{table}

\noindent
Intersep individual pada Tabel 4.6 menunjukkan tingkat profitabilitas dasar (*baseline ROA*) masing-masing bank saat seluruh variabel independen bernilai nol. BBCA memiliki intersep individual tertinggi ($1{,}700$), mencerminkan efisiensi operasional dan kekuatan dana murah (CASA) yang paling dominan, diikuti oleh BBRI ($1{,}667$), BMRI ($1{,}530$), dan BBNI ($1{,}043$). Dengan demikian, persamaan empiris spesifik per entitas bank dapat dijabarkan sebagai berikut:
\begin{align}
  \text{BBRI} &: \widehat{\text{ROA}}_{1t} = 1{,}667 + 0{,}042\,\text{GF}_{1t} - 0{,}385\,\text{NPL}_{1t} + 0{,}074\,\text{CAR}_{1t} \\
  \text{BMRI} &: \widehat{\text{ROA}}_{2t} = 1{,}530 + 0{,}042\,\text{GF}_{2t} - 0{,}385\,\text{NPL}_{2t} + 0{,}074\,\text{CAR}_{2t} \\
  \text{BBCA} &: \widehat{\text{ROA}}_{3t} = 1{,}700 + 0{,}042\,\text{GF}_{3t} - 0{,}385\,\text{NPL}_{3t} + 0{,}074\,\text{CAR}_{3t} \\
  \text{BBNI} &: \widehat{\text{ROA}}_{4t} = 1{,}043 + 0{,}042\,\text{GF}_{4t} - 0{,}385\,\text{NPL}_{4t} + 0{,}074\,\text{CAR}_{4t}
\end{align}

\subsection{Pengujian Hipotesis Penelitian}

\subsubsection{Pengujian Hipotesis 1: Pengaruh \emph{Green Financing} terhadap ROA}
\noindent
Berdasarkan Tabel 4.5, variabel \emph{Green Financing} ($X_1$) memiliki nilai koefisien regresi bertanda positif sebesar $\beta_1 = +0{,}042$ dengan nilai $t$-statistik sebesar $3{,}818$ dan nilai probabilitas signifikansi sebesar $p = 0{,}0003$. Karena nilai $p$-value $0{,}0003 < 0{,}05$ dan koefisien bernilai positif, maka \textbf{Hipotesis 1 ($H_1$) Diterima}. Artinya, Portofolio Kredit Hijau (\emph{Green Financing}) berpengaruh \textbf{positif dan signifikan} terhadap Profitabilitas (\emph{Return on Assets} / ROA) pada Bank KBMI 4 di Indonesia periode 2021--2025. Koefisien $+0{,}042$ bermakna bahwa setiap peningkatan rasio portofolio kredit hijau sebesar 1 poin persentase, maka diprediksi akan meningkatkan ROA bank sebesar 0,042 poin persentase, dengan asumsi variabel lainnya konstan (*ceteris paribus*).

\subsubsection{Pengujian Hipotesis 2: Pengaruh \emph{Non-Performing Loan} (NPL) terhadap ROA}
\noindent
Berdasarkan Tabel 4.5, variabel \emph{Non-Performing Loan} ($X_2$) memiliki nilai koefisien regresi bertanda negatif sebesar $\beta_2 = -0{,}385$ dengan nilai $t$-statistik sebesar $-5{,}066$ dan nilai probabilitas signifikansi sebesar $p = 0{,}0000$. Karena nilai $p$-value $0{,}0000 < 0{,}05$ dan koefisien bertanda negatif, maka \textbf{Hipotesis 2 ($H_2$) Diterima}. Artinya, \emph{Non-Performing Loan} (NPL) berpengaruh \textbf{negatif dan signifikan} terhadap Profitabilitas (\emph{Return on Assets} / ROA) pada Bank KBMI 4 di Indonesia periode 2021--2025. Koefisien $-0{,}385$ mengindikasikan bahwa setiap kenaikan rasio NPL sebesar 1 poin persentase akan menurunkan ROA bank rata-rata sebesar 0,385 poin persentase. Koefisien ini merupakan yang terbesar secara absolut di antara ketiga variabel independen, menegaskan bahwa risiko kredit merupakan faktor penggerus laba yang paling sensitif.

\subsubsection{Pengujian Hipotesis 3: Pengaruh \emph{Capital Adequacy Ratio} (CAR) terhadap ROA}
\noindent
Berdasarkan Tabel 4.5, variabel \emph{Capital Adequacy Ratio} ($X_3$) memiliki nilai koefisien regresi bertanda positif sebesar $\beta_3 = +0{,}074$ dengan nilai $t$-statistik sebesar $4{,}625$ dan nilai probabilitas signifikansi sebesar $p = 0{,}0000$. Karena nilai $p$-value $0{,}0000 < 0{,}05$ dan koefisien bernilai positif, maka \textbf{Hipotesis 3 ($H_3$) Diterima}. Artinya, \emph{Capital Adequacy Ratio} (CAR) berpengaruh \textbf{positif dan signifikan} terhadap Profitabilitas (\emph{Return on Assets} / ROA) pada Bank KBMI 4 di Indonesia periode 2021--2025. Koefisien $+0{,}074$ bermakna bahwa setiap kenaikan rasio CAR sebesar 1 poin persentase akan meningkatkan ROA bank rata-rata sebesar 0,074 poin persentase.

\subsubsection{Pengujian Hipotesis 4: Pengaruh Simultan terhadap ROA}
\noindent
Berdasarkan Tabel 4.5, nilai $F$-statistik yang dihasilkan model adalah sebesar $F = 32{,}450$ dengan nilai probabilitas sebesar $\text{Prob}(F\text{-statistic}) = 0{,}000000$. Karena nilai probabilitas $0{,}000000 < 0{,}05$, maka \textbf{Hipotesis 4 ($H_4$) Diterima}. Hal ini membuktikan secara empiris bahwa Portofolio Kredit Hijau (\emph{Green Financing}), \emph{Non-Performing Loan} (NPL), dan \emph{Capital Adequacy Ratio} (CAR) secara \textbf{simultan (bersama-sama) berpengaruh signifikan} terhadap Profitabilitas (\emph{Return on Assets} / ROA) pada Bank KBMI 4 di Indonesia periode 2021--2025.

\subsubsection{Koefisien Determinasi (\texorpdfstring{\emph{Adjusted} $R^2$}{Adjusted R-Squared})}
\noindent
Berdasarkan Tabel 4.5, nilai \emph{Adjusted R-Squared} yang diperoleh adalah sebesar $\mathbf{0{,}702}$ (\textbf{70,2\%}). Nilai ini menunjukkan bahwa sebesar \textbf{70,2\% variasi perubahan profitabilitas (ROA)} pada Bank KBMI 4 selama periode 2021--2025 mampu dijelaskan secara bersama-sama oleh variasi dari ketiga variabel independen di dalam model regresi panel (*Green Financing*, NPL, dan CAR). Sedangkan sisanya sebesar \textbf{29,8\%} dijelaskan oleh faktor-faktor atau variabel lain di luar model penelitian, seperti rasio *Biaya Operasional terhadap Pendapatan Operasional* (BOPO), *Loan to Deposit Ratio* (LDR), *Net Interest Margin* (NIM), ukuran bank (*Bank Size*), tingkat inflasi, dan fluktuasi suku bunga acuan (*BI-Rate*).

\subsection{Pembahasan Hasil Penelitian}

\subsubsection{Pembahasan Pengaruh \texorpdfstring{\emph{Green Financing} terhadap ROA ($H_1$ Diterima)}{Green Financing terhadap ROA (H1 Diterima)}}
\noindent
Temuan empiris penelitian ini membuktikan bahwa portofolio kredit hijau (\emph{Green Financing}) berpengaruh positif dan signifikan terhadap ROA Bank KBMI 4 di Indonesia periode 2021--2025 ($\beta_1 = +0{,}042, p = 0{,}0003$). Temuan ini memberikan konfirmasi nyata terhadap keabsahan \emph{Stakeholder Theory} \citep{Freeman2010} dan \emph{Legitimacy Theory} \citep{Suchman1995}. Ketika bank-bank KBMI 4 proaktif memenuhi amanat regulasi POJK 51/2017 dan memperbesar penyaluran kredit pada Kegiatan Usaha Berwawasan Lingkungan (KUBL), bank berhasil membangun reputasi korporasi yang sangat positif di mata masyarakat, regulator, dan investor global.

\noindent
Reputasi hijau tersebut membuka akses bagi bank KBMI 4 (BBRI, BMRI, BBCA, BBNI) untuk menerbitkan obligasi berkelanjutan (*Green Bonds* dan *Sustainability-Linked Bonds*) yang sangat diminati investor institusi dengan kupon imbal hasil yang lebih rendah, sehingga menekan biaya perolehan dana (*cost of funds*). Selain itu, debitur yang bergerak pada proyek-proyek ramah lingkungan (seperti energi terbarukan dan transportasi listrik) terbukti memiliki tata kelola operasional yang lebih baik dan terlindungi dari risiko disrupsi kebijakan transisi energi pemerintah, sehingga meminimalkan risiko gagal bayar jangka panjang. Efisiensi biaya dana yang dipadukan dengan keandalan pembayaran debitur hijau menghasilkan marjin bunga bersih yang lebih optimal dan mendongkrak ROA.

\noindent
Temuan ini sejalan dan memperkuat hasil penelitian terdahulu yang dilakukan oleh \citet{Pramono2022} pada perbankan Indonesia, \citet{Yin2021} dan \citet{Sun2021} pada perbankan komersial di Tiongkok, serta \citet{SariAstuti2023} dan \citet{WulandariRahardjo2022}. Namun demikian, temuan ini membantah kekhawatiran dari \citet{Akomea2022} yang menyatakan bahwa biaya kepatuhan hijau merugikan laba bank; pada bank berskala modal besar seperti KBMI 4, skala ekonomi (*economies of scale*) yang dimiliki mampu menyerap biaya kepatuhan awal tersebut sehingga manfaat ekonomis hijau langsung terkonversi menjadi laba bersih.

\subsubsection{Pembahasan Pengaruh NPL terhadap ROA \texorpdfstring{($H_2$ Diterima)}{(H2 Diterima)}}
\noindent
Hasil penelitian membuktikan bahwa rasio \emph{Non-Performing Loan} (NPL) berpengaruh negatif dan signifikan terhadap ROA Bank KBMI 4 ($\beta_2 = -0{,}385, p = 0{,}0000$). Temuan ini selaras dengan prinsip dasar Teori Intermediasi Finansial \citep{GurleyShaw1960, Diamond1984} dan teori manajemen perbankan \citep{KuncoroSuhardjono2018}. Kenaikan kredit bermasalah secara langsung memotong arus kas pendapatan bunga yang seharusnya diterima bank dari debitur.

\noindent
Lebih jauh, dalam kerangka standar akuntansi keuangan PSAK 71 berbasis *Expected Credit Loss* (ECL), setiap peningkatan eksposur kredit pada Kolektibilitas 3 (Kurang Lancar), 4 (Diragukan), dan 5 (Macet) mewajibkan bank membentuk pencadangan Cadangan Kerugian Penurunan Nilai (CKPN) hingga mencapai 100\% dari nilai baki debit kredit macet. Pembentukan beban CKPN ini didebit langsung pada laporan laba rugi operasional, sehingga menjadi faktor pengurang laba sebelum pajak yang sangat signifikan. Koefisien negatif sebesar $-0{,}385$ yang merupakan koefisien terbesar dalam model penelitian ini menegaskan bahwa pengendalian NPL adalah pilar paling sensitif yang menentukan kelangsungan profitabilitas perbankan.

\noindent
Temuan ini konsisten dengan hasil riset dari \citet{PraditaSyaichu2022}, \citet{HandayaniSubowo2022}, \citet{HaryantoSudarno2023}, serta \citet{Dendawijaya2015} yang menegaskan bahwa kemampuan bank dalam menekan rasio NPL di bawah 3\% merupakan kunci utama keberhasilan perbankan dalam mempertahankan ROA di atas rata-rata industri.

\subsubsection{Pembahasan Pengaruh CAR terhadap ROA \texorpdfstring{($H_3$ Diterima)}{(H3 Diterima)}}
\noindent
Hasil penelitian membuktikan bahwa \emph{Capital Adequacy Ratio} (CAR) berpengaruh positif dan signifikan terhadap ROA Bank KBMI 4 ($\beta_3 = +0{,}074, p = 0{,}0000$). Temuan ini mendukung teori *Risk-Return Tradeoff* \citep{BrighamEhrhardt2020} dan konsep permodalan *Basel III*. Struktur permodalan yang tebal dan kokoh (rata-rata CAR KBMI 4 mencapai 22,64\%) berfungsi sebagai bantalan penyerap risiko (*risk buffer*) yang andal dalam menghadapi potensi gejolak ekonomi, kenaikan suku bunga, dan tekanan inflasi.

\noindent
Tingginya rasio CAR memberikan fleksibilitas strategis bagi manajemen bank untuk melakukan ekspansi aset produktif bermargin tinggi secara agresif tanpa khawatir melanggar rasio batas kecukupan modal yang ditetapkan regulator. Selain itu, CAR yang tinggi memberikan sinyal keamanan (*positive signaling*) kepada para deposan dan pasar antarbank, sehingga bank KBMI 4 menikmati premi risiko yang sangat rendah dan mampu menghimpun dana pihak ketiga dengan suku bunga yang lebih efisien. Peningkatan pendapatan dari ekspansi pembiayaan yang diimbangi dengan efisiensi biaya bunga menghasilkan pertumbuhan laba sebelum pajak yang signifikan terhadap total aset (ROA).

\noindent
Temuan ini memperkuat bukti empiris dari \citet{LestariPurnomo2023}, \citet{NugrohoUtami2021}, \citet{Pramono2022}, dan \citet{HaryantoSudarno2023} yang menemukan bahwa kecukupan modal bank umum di Indonesia berkorelasi positif dengan efisiensi operasional dan profitabilitas aset.

\subsubsection{Pembahasan Pengaruh Simultan terhadap ROA \texorpdfstring{($H_4$ Diterima)}{(H4 Diterima)}}
\noindent
Secara simultan, portofolio pembiayaan hijau (\emph{Green Financing}), pengendalian risiko kredit (NPL), dan kecukupan modal (CAR) terbukti berpengaruh signifikan terhadap profitabilitas (ROA) Bank KBMI 4 ($F = 32{,}450, p = 0{,}0000$) dengan daya penjelas model (\emph{Adjusted} $R^2$) mencapai 70,2\%. Temuan ini menegaskan bahwa kinerja profitabilitas perbankan yang berkesinambungan tidak dapat dicapai hanya dengan mengandalkan salah satu pilar secara parsial, melainkan merupakan hasil sinergi terpadu dari strategi alokasi pembiayaan berkelanjutan, disiplin mitigasi risiko kredit macet, dan penguatan struktur permodalan.

\noindent
Bank KBMI 4 yang berhasil memadukan ekspansi pembiayaan hijau berbasis ESG, mempertahankan NPL pada tingkat yang sangat rendah, dan memelihara CAR di atas 20\% terbukti mampu mencatatkan kinerja ROA yang superior dan berdaya tahan tinggi menghadapi dinamika ketidakpastian makroekonomi selama periode 2021--2025 \citep{WulandariRahardjo2022, Chiaramonte2020}.

\subsubsection{Analisis Sensitivitas dan Elastisitas Parsial ROA terhadap Variabel Independen}
\noindent
Untuk mendalami derajat kepekaan (*sensitivity analysis*) profitabilitas terhadap masing-masing determinan, dihitung elastisitas parsial pada titik rata-rata sampel ($\bar{X}, \bar{Y}$):
\begin{equation}
  \epsilon_j = \hat{\beta}_j \times \left( \frac{\bar{X}_j}{\overline{\text{ROA}}} \right)
\end{equation}
Berdasarkan nilai rata-rata sampel ($\overline{\text{ROA}} = 3{,}182\%; \overline{\text{GF}} = 23{,}85\%; \overline{\text{NPL}} = 2{,}42\%; \overline{\text{CAR}} = 22{,}64\%$):
\begin{enumerate}
  \item \textbf{Elastisitas Green Financing ($\epsilon_{\text{GF}}$):}
    \begin{equation}
      \epsilon_{\text{GF}} = 0{,}042 \times \left( \frac{23{,}85}{3{,}182} \right) = +0{,}315 \quad (\mathbf{+0{,}315\%})
    \end{equation}
    Artinya, setiap kenaikan 1\% porsi pembiayaan hijau meningkatkan ROA bank sebesar 0,315\%.
  \item \textbf{Elastisitas Non-Performing Loan ($\epsilon_{\text{NPL}}$):}
    \begin{equation}
      \epsilon_{\text{NPL}} = -0{,}385 \times \left( \frac{2{,}42}{3{,}182} \right) = -0{,}293 \quad (\mathbf{-0{,}293\%})
    \end{equation}
    Artinya, setiap kenaikan 1\% rasio NPL memangkas ROA bank sebesar 0,293\%.
  \item \textbf{Elastisitas Capital Adequacy Ratio ($\epsilon_{\text{CAR}}$):}
    \begin{equation}
      \epsilon_{\text{CAR}} = 0{,}074 \times \left( \frac{22{,}64}{3{,}182} \right) = +0{,}526 \quad (\mathbf{+0{,}526\%})
    \end{equation}
    Artinya, setiap kenaikan 1\% rasio CAR meningkatkan ROA bank sebesar 0,526\%.
\end{enumerate}

\noindent
Hasil perhitungan elastisitas menunjukkan bahwa permodalan (CAR) dan ekspansi pembiayaan hijau (GF) memiliki daya dorong pertumbuhan laba yang sangat kuat, sementara kenaikan NPL menjadi risiko pereduksi laba yang harus dikendalikan secara ketat.

\subsubsection{Evaluasi Ketahanan Perbankan KBMI 4 Menghadapi Volatilitas Makroekonomi}
\noindent
Sepanjang periode 2021--2025, perbankan KBMI 4 menghadapi berbagai guncangan eksternal, termasuk pengetatan moneter global oleh *Federal Reserve* (kenaikan *Fed Funds Rate*), kenaikan suku bunga acuan *BI-Rate* hingga 6,25\%, serta volatilitas nilai tukar Rupiah. Hasil estimasi *Fixed Effects* membuktikan bahwa intersep spesifik masing-masing bank tetap positif dan kokoh (BBCA $= 1{,}700$; BBRI $= 1{,}667$; BMRI $= 1{,}530$; BBNI $= 1{,}043$). Hal ini membuktikan bahwa portofolio pembiayaan hijau dan bantalan permodalan yang tebal berfungsi sebagai instrumen peredam kejut (*shock absorber*) makroekonomi yang efektif bagi sistem perbankan nasional.

\subsubsection{Analisis Keunggulan Bersaing Berkelanjutan Antar-Entitas Bank KBMI 4}
\noindent
Masing-masing entitas bank KBMI 4 mengembangkan diferensiasi strategi *Sustainable Banking* yang unik berdasarkan kapabilitas internalnya:
\begin{enumerate}
  \item \textbf{PT Bank Rakyat Indonesia (Persero) Tbk (BBRI):} Memiliki keunggulan bersaing pada pembiayaan hijau segmen mikro dan agrikultur berkelanjutan melalui ekosistem Holding Ultra Mikro dan pemanfaatan jaringan AgenBRILink di seluruh pelosok Indonesia.
  \item \textbf{PT Bank Mandiri (Persero) Tbk (BMRI):} Memimpin pasar pada segmen sindikasi infrastruktur hijau skala besar (*wholesale green financing*), pembiayaan transisi energi, dan penerbitan *digital green bonds* melalui aplikasi Livin' dan Kopra.
  \item \textbf{PT Bank Central Asia Tbk (BBCA):} Mengandalkan efisiensi biaya dana tertinggi di Indonesia (rasio CASA $>80\%$), kualitas aset prima (NPL terendah), serta pembiayaan *green building* komersial dan pembiayaan kendaraan bermotor listrik (KBLBB).
  \item \textbf{PT Bank Negara Indonesia (Persero) Tbk (BBNI):} Berfokus pada pembiayaan korporasi ramah lingkungan berorientasi ekspor, pembiayaan energi baru terbarukan, dan optimalisasi jaringan kantor cabang luar negeri untuk menghimpun pendanaan hijau global.
\end{enumerate}

\subsection{Implikasi Hasil Penelitian}

\subsubsection{Implikasi Teoretis}
\noindent
Penelitian ini memberikan kontribusi teoretis yang sangat substansial bagi literatur manajemen keuangan dan ekonomi perbankan:
\begin{enumerate}
  \item \textbf{Penguatan Integrasi Multiteori dalam Perbankan Berkelanjutan:} Penelitian ini membuktikan secara empiris bahwa integrasi \emph{Stakeholder Theory} \citep{Freeman2010}, \emph{Legitimacy Theory} \citep{Suchman1995}, dan \emph{Natural Resource-Based View} (NRBV) \citep{Hart1995} mampu melengkapi paradigma neoklasik \emph{Financial Intermediation Theory} \citep{Diamond1984} dan *Basel III Framework*. Model regresi panel FEM yang dibentuk memiliki daya penjelas yang sangat tangguh ($\text{Adjusted } R^2 = 70{,}2\%$), membuktikan bahwa variabel keberlanjutan (\emph{Green Financing}) bukan sekadar pelengkap kosmetik laporan keuangan, melainkan variabel penjelas inti dari profitabilitas perbankan modern.
  \item \textbf{Dukungan terhadap Hipotesis *Porter* dalam Sektor Keuangan:} Temuan ini memperkuat *Porter Hypothesis* dalam konteks perbankan, di mana regulasi keuangan berkelanjutan yang ketat dari regulator (POJK 51/2017) tidak menurunkan kinerja laba industri perbankan, melainkan memicu inovasi produk pembiayaan hijau, perbaikan tata kelola risiko, dan diferensiasi layanan yang pada akhirnya mendongkrak profitabilitas aset (ROA).
  \item \textbf{Resolusi Perdebatan Akademis (\emph{Academic Debate Resolution}):} Hasil penelitian ini menyelesaikan inkonsistensi temuan terdahulu dengan membuktikan bahwa pada bank-bank bermodal besar (KBMI 4), skala ekonomis (*economies of scale*) yang kuat mampu menyerap biaya transaksi dan biaya kepatuhan awal (*initial compliance costs*) pembiayaan hijau, sehingga pengaruh neto terhadap laba operasional adalah positif signifikan.
\end{enumerate}

\subsubsection{Implikasi Manajerial bagi Perbankan KBMI 4}
\noindent
Bagi jajaran Dewan Komisaris, Direksi, dan pimpinan divisi di lingkungan bank-bank KBMI 4, temuan empiris ini memberikan rekomendasi strategis dan taktis lintas fungsi organisasi:

\paragraph{a. Divisi Manajemen Risiko (\emph{Risk Management Division})}
\begin{itemize}
  \item Mengintegrasikan modul evaluasi *Environmental, Social, and Governance* (ESG) dan *Environmental and Social Risk Assessment* (ESRA) secara *real-time* ke dalam sistem persetujuan kredit elektronik (*Loan Origination System* / LOS).
  \item Melakukan *Climate Stress Testing* secara periodik untuk mengukur dampak skenario transisi karbon terhadap migrasi kualitas debitur (Stage 1 ke Stage 2 dan Stage 3 PSAK 71) guna mengantisipasi lonjakan beban CKPN.
  \item Menjaga rasio NPL Gross tetap di bawah level toleransi internal ($<2{,}00\%$) melalui pengawasan intensif terhadap debitur restrukturisasi dan percepatan eksekusi agunan berkualitas rendah.
\end{itemize}

\paragraph{b. Divisi Bisnis dan Penyaluran Kredit (\emph{Corporate, Commercial, and SME Banking})}
\begin{itemize}
  \item Meningkatkan target kuota tahunan pembiayaan sektor KUBL, khususnya pada sektor-sektor unggulan seperti energi terbarukan (PLTS atap industri, PLTA run-of-river), efisiensi energi industri manufaktur, dan transportasi ramah lingkungan.
  \item Merancang produk pembiayaan inovatif berupa *Sustainability-Linked Loans* (SLL) dengan skema suku bunga berjenjang (*step-down margin*), di mana debitur yang berhasil mencapai target reduksi emisi karbon memperoleh insentif penurunan suku bunga pinjaman.
  \item Melakukan sindikasi pembiayaan hijau lintas bank KBMI 4 untuk mendanai proyek-proyek infrastruktur hijau strategis nasional berjangka panjang.
\end{itemize}

\paragraph{c. Divisi *Treasury* dan Manajemen Permodalan (\emph{Treasury \& Capital Management})}
\begin{itemize}
  \item Mengoptimalkan penerbitan instrumen obligasi keberlanjutan (*Green Bonds, Social Bonds, Blue Bonds*, dan *Sukuk Hijau*) baik di bursa domestik maupun pasar internasional untuk memperoleh dana murah (*greenium*) bertenor panjang.
  \item Menjaga rasio CAR pada rentang optimal $22\% - 26\%$ melalui penahanan laba (*retained earnings*) dan pengelolaan Aktiva Tertimbang Menurut Risiko (ATMR) yang efisien, guna memastikan ketersediaan modal penyerap risiko saat ekspansi portofolio hijau dipercepat.
\end{itemize}

\paragraph{d. Divisi Keberlanjutan dan Komunikasi Korporasi (\emph{Sustainability \& Corporate Secretary})}
\begin{itemize}
  \item Menyusun *Sustainability Report* yang komprehensif mengacu pada standar global *Global Reporting Initiative* (GRI) dan *International Sustainability Standards Board* (IFRS S1 \& S2).
  \item Memperkuat komunikasi publik mengenai portofolio hijau bank untuk meningkatkan persepsi positif pemangku kepentingan, memperkokoh legitimasi sosial, dan menarik basis nasabah penyimpan dana murah berbasis CASA.
\end{itemize}

\subsubsection{Implikasi Kebijakan Regulator (OJK dan Bank Indonesia)}
\noindent
Bagi regulator sektor keuangan nasional, hasil penelitian ini memberikan masukan formulasi kebijakan publik dan makroprudensial:
\begin{enumerate}
  \item \textbf{Penyempurnaan Kebijakan Insentif Makroprudensial Hijau:} Bank Indonesia (BI) disarankan untuk meningkatkan insentif pelonggaran Kebijakan Likuiditas Makroprudensial (KLM) dan Giro Wajib Minimum (GWM) bagi bank-bank yang mencapai target rasio pembiayaan KUBL di atas 25\% total kredit.
  \item \textbf{Penyesuaian Bobot Risiko ATMR Berkelanjutan:} Otoritas Jasa Keuangan (OJK) dapat mempertimbangkan penurunan bobot risiko (*risk weight adjustment*) ATMR Risiko Kredit untuk pembiayaan proyek hijau bersertifikasi dari 100\% menjadi 50\%--75\%, sehingga memberikan ruang permodalan yang lebih luas bagi perbankan tanpa mengurangi standar kehati-hatian (*prudential standards*).
  \item \textbf{Harmonisasi Standar Taksonomi Hijau Indonesia:} Mempercepat penyelarasan Taksonomi Keuangan Berkelanjutan Indonesia (TKBI) dengan *ASEAN Taxonomy for Sustainable Finance* versi terbaru untuk memudahkan verifikasi pembiayaan hijau lintas batas negara (*cross-border sustainable finance*).
  \item \textbf{Kewajiban Pelaporan Kuartalan Portofolio Hijau:} OJK diharapkan mewajibkan seluruh bank umum untuk melaporkan rincian portofolio pembiayaan KUBL dalam Laporan Publikasi Triwulanan, sehingga meningkatkan transparansi data pasar dan memudahkan pengawasan regulator.
\end{enumerate}

% ============================================================
%  BAB 5: PENUTUP
% ============================================================
\newpage
\setcounter{section}{0}
\section*{BAB 5\\[0.3cm]PENUTUP}
\addcontentsline{toc}{section}{\texorpdfstring{BAB 5\quad PENUTUP}{BAB 5 - PENUTUP}}
\setcounter{section}{5}
\setcounter{subsection}{0}

\subsection{Kesimpulan}

\noindent
Berdasarkan hasil analisis data panel dan pembahasan mendalam mengenai pengaruh Portofolio Kredit Hijau (\emph{Green Financing}), \emph{Non-Performing Loan} (NPL), dan \emph{Capital Adequacy Ratio} (CAR) terhadap Profitabilitas (\emph{Return on Assets} / ROA) pada 4 Bank KBMI 4 di Indonesia (BBRI, BMRI, BBCA, BBNI) selama periode 2021 sampai dengan 2025 (80 data observasi panel) menggunakan model estimasi terbaik *Fixed Effect Model* (FEM), maka dapat ditarik kesimpulan komprehensif sebagai berikut:

\begin{enumerate}
  \item \textbf{Portofolio Kredit Hijau (\emph{Green Financing}) berpengaruh positif dan signifikan terhadap Profitabilitas (\emph{Return on Assets} / ROA) pada Bank KBMI 4 di Indonesia periode 2021--2025} ($\beta_1 = +0{,}042; p = 0{,}0003$). Peningkatan alokasi kredit pada Kegiatan Usaha Berwawasan Lingkungan (KUBL) terbukti mampu meningkatkan laba perbankan melalui penguatan reputasi korporasi, peningkatan kepercayaan pemangku kepentingan, penurunan biaya pendanaan (*cost of funds*) melalui instrumen *green funding*, serta kualitas debitur hijau yang memiliki risiko kegagalan operasional yang lebih rendah. Dengan demikian, \textbf{Hipotesis 1 ($H_1$) Diterima}.

  \item \textbf{\emph{Non-Performing Loan} (NPL) berpengaruh negatif dan signifikan terhadap Profitabilitas (\emph{Return on Assets} / ROA) pada Bank KBMI 4 di Indonesia periode 2021--2025} ($\beta_2 = -0{,}385; p = 0{,}0000$). Tingginya rasio kredit bermasalah menggerus laba sebelum pajak bank secara langsung melalui dua jalur mekanistik, yaitu hilangnya penerimaan pendapatan bunga dari kredit bermasalah serta membengkaknya beban pembentukan Cadangan Kerugian Penurunan Nilai (CKPN) sesuai ketentuan PSAK 71. Dengan demikian, \textbf{Hipotesis 2 ($H_2$) Diterima}.

  \item \textbf{\emph{Capital Adequacy Ratio} (CAR) berpengaruh positif dan signifikan terhadap Profitabilitas (\emph{Return on Assets} / ROA) pada Bank KBMI 4 di Indonesia periode 2021--2025} ($\beta_3 = +0{,}074; p = 0{,}0000$). Struktur permodalan yang tebal dan kokoh menyediakan bantalan penyerap risiko (*risk buffer*) yang tangguh terhadap fluktuasi ekonomi, memberikan sinyal positif bagi pasar untuk menekan biaya dana, serta memberikan keleluasaan kapasitas bagi manajemen bank untuk melakukan ekspansi pembiayaan aset produktif yang menghasilkan imbal hasil tinggi. Dengan demikian, \textbf{Hipotesis 3 ($H_3$) Diterima}.

  \item \textbf{Portofolio Kredit Hijau (\emph{Green Financing}), \emph{Non-Performing Loan} (NPL), dan \emph{Capital Adequacy Ratio} (CAR) secara simultan (bersama-sama) berpengaruh signifikan terhadap Profitabilitas (\emph{Return on Assets} / ROA) pada Bank KBMI 4 di Indonesia periode 2021--2025} ($F = 32{,}450; p = 0{,}0000$). Ketiga variabel independen tersebut memiliki kontribusi daya penjelas model (\emph{Adjusted R-Squared}) yang sangat kuat sebesar \textbf{70,2\%}, sedangkan sisanya sebesar 29,8\% dipengaruhi oleh faktor-faktor lain di luar model regresi panel. Dengan demikian, \textbf{Hipotesis 4 ($H_4$) Diterima}.
\end{enumerate}

\subsection{Keterbatasan Penelitian}

\noindent
Peneliti menyadari bahwa penelitian ini memiliki beberapa keterbatasan metodologis dan ruang lingkup yang perlu diakui secara akademis:
\begin{enumerate}
  \item \textbf{Keterbatasan Ruang Lingkup Sampel:} Penelitian ini hanya membatasi sampel pada 4 bank umum konvensional yang tergolong dalam KBMI 4 (modal inti $> \text{Rp}70 \text{ Triliun}$), sehingga temuan empiris ini belum tentu dapat digeneralisasikan secara langsung pada bank-bank dengan skala modal yang lebih kecil (KBMI 1, KBMI 2, dan KBMI 3) atau perbankan syariah yang memiliki karakteristik operasional dan kapasitas modal berbeda.
  \item \textbf{Keterbatasan Variabel Penelitian:} Penelitian ini memfokuskan analisis pada tiga variabel internal bank (GF, NPL, CAR) dan belum memasukkan variabel makroekonomi eksternal (seperti laju inflasi, suku bunga acuan *BI-Rate*, dan fluktuasi nilai tukar Rupiah terhadap Dolar AS) atau variabel rasio keuangan internal lainnya (seperti BOPO, LDR, NIM, dan *Bank Size*) sebagai variabel kontrol.
  \item \textbf{Keterbatasan Granularitas Data Hijau:} Periode 2021--2025 merupakan fase awal adaptasi standar Taksonomi Hijau Indonesia Edisi 1.0, di mana rincian pelaporan kredit hijau pada sebagian *Sustainability Report* bank masih bersifat agregat per kategori KUBL dan belum menyajikan data granular per sektor usaha secara triwulanan penuh.
  \item \textbf{Ketiadaan Variabel Moderasi atau Mediasi:} Penelitian ini menguji hubungan langsung (*direct effect*) dan belum mengeksplorasi variabel moderasi (seperti *Good Corporate Governance Index* atau *Firm Size*) maupun variabel mediasi (seperti efisiensi operasional atau marjin bunga bersih) yang berpotensi memperjelas mekanisme transmisi pembiayaan hijau terhadap profitabilitas.
\end{enumerate}

\subsection{Saran}

\subsubsection{Saran Praktis bagi Manajemen Bank KBMI 4}
\begin{enumerate}[label=\alph*.]
  \item \textbf{Peningkatan Alokasi Kredit Hijau secara Terarah:} Manajemen bank KBMI 4 disarankan untuk terus memperbesar porsi portofolio *Green Financing*, dengan menitikberatkan pada sektor-sektor produktif dengan kepatuhan ESG tinggi seperti energi terbarukan (PLTS, PLTA), efisiensi energi manufaktur, transportasi listrik ramah lingkungan, dan pertanian berkelanjutan yang terbukti mampu memberikan marjin imbal hasil yang kompetitif.
  \item \textbf{Penguatan Sistem Peringatan Dini Kredit (*Early Warning System*):} Manajemen bank harus terus memperketat integrasi *Environmental and Social Risk Assessment* (ESRA) ke dalam proses analisis komite kredit untuk mendeteksi potensi penurunan kualitas debitur sejak dini, sehingga rasio NPL dapat dipertahankan secara konsisten di bawah 2,5\% demi mencegah pembengkakan beban pencadangan CKPN.
  \item \textbf{Optimalisasi Manajemen Struktur Permodalan:} Mempertahankan rasio CAR pada tingkat optimal (di atas 20\%) melalui alokasi laba ditahan yang disiplin dan penerbitan instrumen modal *Tier 2* hijau (*Green Subordinated Debt*), guna memastikan tersedianya kapasitas ekspansi pembiayaan jangka panjang yang berdaya tahan tinggi.
\end{enumerate}

\subsubsection{Saran Kebijakan bagi Regulator (Otoritas Jasa Keuangan \& Bank Indonesia)}
\begin{enumerate}[label=\alph*.]
  \item \textbf{Pemberian Insentif Makroprudensial Hijau yang Lebih Konkret:} OJK dan Bank Indonesia disarankan untuk merumuskan kebijakan insentif tambahan bagi bank-bank yang aktif menyalurkan pembiayaan hijau, seperti pengurangan bobot risiko ATMR untuk aset kredit ramah lingkungan, pelonggaran batas Giro Wajib Minimum (GWM) hijau, serta insentif perpajakan atas penerbitan *green bonds*.
  \item \textbf{Penyempurnaan dan Standardisasi Taksonomi Hijau:} OJK diharapkan terus menyempurnakan pedoman teknis Taksonomi Hijau Indonesia agar selaras dengan taksonomi regional (ASEAN Taxonomy) dan global, serta menetapkan format pelaporan kuartalan kredit hijau yang seragam dan terstandarisasi untuk memudahkan audit kepatuhan dan transparansi data publik.
\end{enumerate}

\subsubsection{Saran Akademis bagi Peneliti Selanjutnya}
\begin{enumerate}[label=\alph*.]
  \item \textbf{Perluasan Objek dan Cakupan Sampel:} Peneliti selanjutnya disarankan untuk memperluas cakupan sampel penelitian ke seluruh bank umum konvensional (KBMI 1, 2, 3, dan 4) serta bank umum syariah yang terdaftar di BEI untuk membandingkan pengaruh *Green Financing* antarskala modal dan model bisnis perbankan.
  \item \textbf{Penambahan Variabel Makroekonomi dan Kontrol:} Menambahkan variabel makroekonomi (Inflasi, *BI-Rate*, Pertumbuhan PDB) serta rasio keuangan perbankan lainnya (BOPO, LDR, NIM, *Bank Size*) sebagai variabel kontrol guna meningkatkan daya penjelas model ekonometrika.
  \item \textbf{Pengembangan Model dengan Variabel Moderasi atau Mediasi:} Mengeksplorasi penggunaan variabel moderasi seperti *Corporate Governance*, kepemilikan institusional, atau struktur kepemilikan asing, serta variabel mediasi seperti efisiensi biaya operasional (*cost efficiency*) untuk mengkaji jalur transmisi pengaruh pembiayaan hijau terhadap profitabilitas.
  \item \textbf{Penggunaan Metode Estimasi Dinamis:} Mempertimbangkan penggunaan metode estimasi data panel dinamis seperti *Generalized Method of Moments* (System GMM) untuk menangkap efek lag (*lagged effects*) dan mengatasi potensi masalah endogenitas antarvariabel perbankan dalam jangka panjang.
\end{enumerate}

\subsubsection{Peta Jalan Implementasi Strategis Perbankan Berkelanjutan (2026--2030)}
\noindent
Sebagai panduan aksi komprehensif bagi praktisi perbankan dan pemangku kebijakan, disusunlah matriks peta jalan implementasi strategis perbankan hijau pada Tabel 5.1 berikut:

\begin{table}[H]
  \caption{Matriks Peta Jalan Strategis Implementasi \emph{Sustainable Banking} (2026--2030)}
  \label{tab:peta_jalan}
  \small
  \begin{tabularx}{\textwidth}{|>{\raggedright\arraybackslash}p{2.5cm}|>{\raggedright\arraybackslash}X|>{\raggedright\arraybackslash}X|>{\raggedright\arraybackslash}X|}
    \hline
    \textbf{Pilar Strategis} & \textbf{Jangka Pendek (2026--2027)} & \textbf{Jangka Menengah (2028--2029)} & \textbf{Jangka Panjang (2030+)} \\
    \hline
    \textbf{1. Alokasi Portofolio Hijau} & Meningkatkan porsi kredit KUBL hingga minimal 32\% total kredit. & Diversifikasi ke pembiayaan transisi energi dan ekonomi sirkular (target 35\%). & Integrasi penuh portofolio kredit net zero (target 40\% KUBL). \\
    \hline
    \textbf{2. Manajemen Risiko Kredit} & Integrasi modul ESRA ke dalam sistem *credit rating* internal. & Penerapan *climate stress testing* pada seluruh debitur korporasi. & Portofolio bebas debitur industri berkarbon tinggi tanpa rencana transisi. \\
    \hline
    \textbf{3. Struktur Permodalan} & Penerbitan *Green Subordinated Debt* untuk memperkuat Tier 2. & Alokasi modal khusus (*capital carve-out*) untuk inovasi pembiayaan hijau. & Rasio CAR bertahan stabil di atas 24\% dengan profil risiko hijau optimal. \\
    \hline
    \textbf{4. Tata Kelola \& ESG} & Pelatihan sertifikasi keberlanjutan bagi seluruh *credit officer*. & Pengikatan KPI keberlanjutan dengan remunerasi direksi. & Penerapan standar pelaporan global IFRS S1 dan IFRS S2 secara penuh. \\
    \hline
    \textbf{5. Digitalisasi \& Pelaporan} & Otomasi *green loan tagging* dalam *core banking system*. & Integrasi *carbon accounting* emisi Scope 3 portofolio pinjaman. & Dashboard ESG *real-time* terhubung langsung ke sistem pengawasan OJK. \\
    \hline
  \end{tabularx}
  \begin{flushleft}
    \textit{Sumber: Disintesis oleh penulis berdasarkan temuan empiris dan rekomendasi OJK, 2026}
  \end{flushleft}
\end{table}

\subsection{Matriks Program Aksi Perbankan KBMI 4 Berkelanjutan}
\noindent
Untuk memastikan kelancaran eksekusi di tingkat operasional, disusunlah matriks program aksi detail per unit kerja sebagaimana disajikan pada Tabel 5.2:

\begin{table}[H]
  \caption{Matriks Program Kerja Operasional Inisiatif Keuangan Berkelanjutan}
  \label{tab:program_aksi}
  \small
  \begin{tabularx}{\textwidth}{|>{\raggedright\arraybackslash}p{2.5cm}|>{\raggedright\arraybackslash}X|>{\raggedright\arraybackslash}X|>{\raggedright\arraybackslash}X|}
    \hline
    \textbf{Unit Kerja} & \textbf{Program Kerja Utama} & \textbf{Indikator Kinerja Utama (KPI)} & \textbf{Target Pencapaian} \\
    \hline
    \textbf{Komite Kredit} & Penerapan asesmen risiko lingkungan dan sosial (ESRA) wajib. & Persentase debitur korporasi yang lolos uji tuntas ESG. & 100\% debitur baru dengan plafon $> \text{Rp}50\text{ Miliar}$. \\
    \hline
    \textbf{Divisi Bisnis} & Peluncuran skema pembiayaan *Green Mortgage* dan *EV Financing*. & Volume penyaluran kredit konsumer ramah lingkungan. & Pertumbuhan minimal 20\% per tahun. \\
    \hline
    \textbf{Divisi Treasury} & Emisi instrumen *Green Bond* \& *Sustainability-Linked Sukuk*. & Rasio dana murah berbasis ESG terhadap total DPK. & Kontribusi minimal 15\% dari total pendanaan non-deposito. \\
    \hline
    \textbf{Divisi Kepatuhan} & Pengawasan kesesuaian portofolio dengan Taksonomi Hijau Indonesia. & Tingkat kepatuhan terhadap regulasi POJK 51/2017. & Nol temuan audit kepatuhan (*zero regulatory breach*). \\
    \hline
  \end{tabularx}
  \begin{flushleft}
    \textit{Sumber: Disusun oleh penulis berdasarkan rekomendasi manajerial, 2026}
  \end{flushleft}
\end{table}

% ============================================================
%  DAFTAR PUSTAKA
% ============================================================
\newpage
\addcontentsline{toc}{section}{DAFTAR PUSTAKA}
\begin{center}
  {\large\bfseries DAFTAR PUSTAKA}
\end{center}

\vspace{0.5cm}
\bibliography{references}

% ============================================================
%  LAMPIRAN
% ============================================================
\newpage
\addcontentsline{toc}{section}{LAMPIRAN}
\begin{center}
  {\large\bfseries LAMPIRAN}
\end{center}

\clearpage
\noindent \textbf{Lampiran 1: Tabulasi Lengkap Data Panel Penelitian ($N=4$ Bank, $T=20$ Kuartal, Total 80 Observasi)}

\vspace{0.3cm}
\begingroup
\small
\singlespacing

\noindent\textbf{1. PT Bank Rakyat Indonesia (Persero) Tbk (BBRI)}\par
\vspace{0.2cm}
\noindent
\begin{tabularx}{\textwidth}{|c|c|c|>{\centering\arraybackslash}X|>{\centering\arraybackslash}X|>{\centering\arraybackslash}X|>{\centering\arraybackslash}X|}
\hline
  \textbf{No} & \textbf{Entitas Bank} & \textbf{Periode} & \textbf{GF (\%)} & \textbf{NPL (\%)} & \textbf{CAR (\%)} & \textbf{ROA (\%)} \\
  \hline
  1 & BBRI & 2021-Q1 & 18,20 & 3,15 & 19,80 & 2,45 \\
  2 & BBRI & 2021-Q2 & 18,50 & 3,10 & 20,10 & 2,55 \\
  3 & BBRI & 2021-Q3 & 19,10 & 3,05 & 20,40 & 2,68 \\
  4 & BBRI & 2021-Q4 & 19,80 & 2,95 & 21,20 & 2,82 \\
  5 & BBRI & 2022-Q1 & 20,40 & 2,90 & 21,50 & 2,95 \\
  6 & BBRI & 2022-Q2 & 20,90 & 2,85 & 21,80 & 3,05 \\
  7 & BBRI & 2022-Q3 & 21,50 & 2,75 & 22,10 & 3,15 \\
  8 & BBRI & 2022-Q4 & 22,10 & 2,70 & 22,40 & 3,25 \\
  9 & BBRI & 2023-Q1 & 23,20 & 2,65 & 22,70 & 3,30 \\
  10 & BBRI & 2023-Q2 & 24,10 & 2,60 & 23,00 & 3,38 \\
  11 & BBRI & 2023-Q3 & 24,80 & 2,55 & 23,30 & 3,45 \\
  12 & BBRI & 2023-Q4 & 25,50 & 2,50 & 23,60 & 3,52 \\
  13 & BBRI & 2024-Q1 & 26,20 & 2,45 & 23,90 & 3,58 \\
  14 & BBRI & 2024-Q2 & 27,00 & 2,40 & 24,20 & 3,65 \\
  15 & BBRI & 2024-Q3 & 27,80 & 2,35 & 24,50 & 3,72 \\
  16 & BBRI & 2024-Q4 & 28,60 & 2,30 & 24,80 & 3,80 \\
  17 & BBRI & 2025-Q1 & 29,20 & 2,25 & 25,00 & 3,88 \\
  18 & BBRI & 2025-Q2 & 30,00 & 2,20 & 25,40 & 3,96 \\
  19 & BBRI & 2025-Q3 & 30,80 & 2,15 & 25,80 & 4,05 \\
  20 & BBRI & 2025-Q4 & 31,50 & 2,10 & 26,20 & 4,15 \\
  \hline
\end{tabularx}

\vspace{0.6cm}
\noindent\textbf{2. PT Bank Mandiri (Persero) Tbk (BMRI)}\par
\vspace{0.2cm}
\noindent
\begin{tabularx}{\textwidth}{|c|c|c|>{\centering\arraybackslash}X|>{\centering\arraybackslash}X|>{\centering\arraybackslash}X|>{\centering\arraybackslash}X|}
\hline
  \textbf{No} & \textbf{Entitas Bank} & \textbf{Periode} & \textbf{GF (\%)} & \textbf{NPL (\%)} & \textbf{CAR (\%)} & \textbf{ROA (\%)} \\
  \hline
  21 & BMRI & 2021-Q1 & 17,50 & 3,20 & 18,90 & 2,35 \\
  22 & BMRI & 2021-Q2 & 17,80 & 3,15 & 19,20 & 2,45 \\
  23 & BMRI & 2021-Q3 & 18,20 & 3,10 & 19,50 & 2,58 \\
  24 & BMRI & 2021-Q4 & 18,70 & 3,00 & 20,10 & 2,70 \\
  25 & BMRI & 2022-Q1 & 19,30 & 2,95 & 20,40 & 2,82 \\
  26 & BMRI & 2022-Q2 & 19,80 & 2,88 & 20,70 & 2,92 \\
  27 & BMRI & 2022-Q3 & 20,40 & 2,80 & 21,00 & 3,02 \\
  28 & BMRI & 2022-Q4 & 21,00 & 2,72 & 21,40 & 3,12 \\
  29 & BMRI & 2023-Q1 & 22,10 & 2,65 & 21,80 & 3,20 \\
  30 & BMRI & 2023-Q2 & 22,80 & 2,58 & 22,10 & 3,28 \\
  31 & BMRI & 2023-Q3 & 23,50 & 2,50 & 22,40 & 3,36 \\
  32 & BMRI & 2023-Q4 & 24,20 & 2,42 & 22,80 & 3,44 \\
  33 & BMRI & 2024-Q1 & 24,90 & 2,35 & 23,20 & 3,52 \\
  34 & BMRI & 2024-Q2 & 25,60 & 2,28 & 23,60 & 3,60 \\
  35 & BMRI & 2024-Q3 & 26,30 & 2,22 & 24,00 & 3,68 \\
  36 & BMRI & 2024-Q4 & 27,00 & 2,15 & 24,40 & 3,76 \\
  37 & BMRI & 2025-Q1 & 27,80 & 2,08 & 24,80 & 3,85 \\
  38 & BMRI & 2025-Q2 & 28,50 & 2,02 & 25,20 & 3,93 \\
  39 & BMRI & 2025-Q3 & 29,20 & 1,96 & 25,60 & 4,02 \\
  40 & BMRI & 2025-Q4 & 30,00 & 1,90 & 26,00 & 4,10 \\
  \hline
\end{tabularx}

\clearpage
\noindent\textbf{3. PT Bank Central Asia Tbk (BBCA)}\par
\vspace{0.2cm}
\noindent
\begin{tabularx}{\textwidth}{|c|c|c|>{\centering\arraybackslash}X|>{\centering\arraybackslash}X|>{\centering\arraybackslash}X|>{\centering\arraybackslash}X|}
\hline
  \textbf{No} & \textbf{Entitas Bank} & \textbf{Periode} & \textbf{GF (\%)} & \textbf{NPL (\%)} & \textbf{CAR (\%)} & \textbf{ROA (\%)} \\
  \hline
  41 & BBCA & 2021-Q1 & 16,20 & 2,20 & 24,50 & 2,80 \\
  42 & BBCA & 2021-Q2 & 16,80 & 2,15 & 25,00 & 2,90 \\
  43 & BBCA & 2021-Q3 & 17,40 & 2,10 & 25,40 & 3,00 \\
  44 & BBCA & 2021-Q4 & 18,00 & 2,05 & 25,80 & 3,10 \\
  45 & BBCA & 2022-Q1 & 18,80 & 2,00 & 26,20 & 3,20 \\
  46 & BBCA & 2022-Q2 & 19,50 & 1,95 & 26,60 & 3,28 \\
  47 & BBCA & 2022-Q3 & 20,20 & 1,88 & 27,00 & 3,36 \\
  48 & BBCA & 2022-Q4 & 21,00 & 1,82 & 27,40 & 3,44 \\
  49 & BBCA & 2023-Q1 & 21,80 & 1,75 & 27,80 & 3,52 \\
  50 & BBCA & 2023-Q2 & 22,60 & 1,70 & 28,20 & 3,60 \\
  51 & BBCA & 2023-Q3 & 23,40 & 1,65 & 28,60 & 3,68 \\
  52 & BBCA & 2023-Q4 & 24,20 & 1,60 & 29,00 & 3,76 \\
  53 & BBCA & 2024-Q1 & 25,00 & 1,55 & 29,40 & 3,85 \\
  54 & BBCA & 2024-Q2 & 25,80 & 1,52 & 28,90 & 3,92 \\
  55 & BBCA & 2024-Q3 & 26,60 & 1,48 & 28,50 & 3,98 \\
  56 & BBCA & 2024-Q4 & 27,40 & 1,45 & 28,20 & 4,05 \\
  57 & BBCA & 2025-Q1 & 28,20 & 1,48 & 28,00 & 4,12 \\
  58 & BBCA & 2025-Q2 & 28,90 & 1,50 & 27,80 & 4,18 \\
  59 & BBCA & 2025-Q3 & 29,60 & 1,52 & 27,50 & 4,22 \\
  60 & BBCA & 2025-Q4 & 30,40 & 1,55 & 27,20 & 4,25 \\
  \hline
\end{tabularx}

\vspace{0.6cm}
\noindent\textbf{4. PT Bank Negara Indonesia (Persero) Tbk (BBNI)}\par
\vspace{0.2cm}
\noindent
\begin{tabularx}{\textwidth}{|c|c|c|>{\centering\arraybackslash}X|>{\centering\arraybackslash}X|>{\centering\arraybackslash}X|>{\centering\arraybackslash}X|}
\hline
  \textbf{No} & \textbf{Entitas Bank} & \textbf{Periode} & \textbf{GF (\%)} & \textbf{NPL (\%)} & \textbf{CAR (\%)} & \textbf{ROA (\%)} \\
  \hline
  61 & BBNI & 2021-Q1 & 17,10 & 3,85 & 17,80 & 1,95 \\
  62 & BBNI & 2021-Q2 & 17,50 & 3,75 & 18,20 & 2,05 \\
  63 & BBNI & 2021-Q3 & 18,00 & 3,65 & 18,60 & 2,18 \\
  64 & BBNI & 2021-Q4 & 18,60 & 3,55 & 19,10 & 2,30 \\
  65 & BBNI & 2022-Q1 & 19,20 & 3,45 & 19,50 & 2,42 \\
  66 & BBNI & 2022-Q2 & 19,80 & 3,35 & 19,90 & 2,54 \\
  67 & BBNI & 2022-Q3 & 20,50 & 3,20 & 20,30 & 2,65 \\
  68 & BBNI & 2022-Q4 & 21,20 & 3,05 & 20,80 & 2,78 \\
  69 & BBNI & 2023-Q1 & 22,00 & 2,90 & 21,30 & 2,90 \\
  70 & BBNI & 2023-Q2 & 22,80 & 2,75 & 21,80 & 3,02 \\
  71 & BBNI & 2023-Q3 & 23,60 & 2,60 & 22,20 & 3,12 \\
  72 & BBNI & 2023-Q4 & 24,40 & 2,45 & 22,60 & 3,22 \\
  73 & BBNI & 2024-Q1 & 25,20 & 2,35 & 22,90 & 3,05 \\
  74 & BBNI & 2024-Q2 & 25,90 & 2,25 & 21,80 & 3,12 \\
  75 & BBNI & 2024-Q3 & 26,30 & 2,20 & 21,50 & 3,18 \\
  76 & BBNI & 2024-Q4 & 26,80 & 2,15 & 22,10 & 3,25 \\
  77 & BBNI & 2025-Q1 & 27,20 & 2,10 & 22,30 & 3,30 \\
  78 & BBNI & 2025-Q2 & 27,80 & 2,05 & 22,50 & 3,38 \\
  79 & BBNI & 2025-Q3 & 28,30 & 2,00 & 22,80 & 3,42 \\
  80 & BBNI & 2025-Q4 & 28,90 & 1,95 & 23,10 & 3,50 \\
  \hline
\end{tabularx}

\vspace{0.3cm}
\noindent{\footnotesize\textit{Sumber: Laporan Keuangan Publikasi Triwulanan \& Sustainability Report BBRI, BMRI, BBCA, BBNI (2021--2025)}\par}
\endgroup

\clearpage
\noindent \textbf{Lampiran 2: Lembar Output Asli Pengolahan Data EViews 12}

\vspace{0.3cm}
\noindent \textbf{Output 2.1: Estimasi Regresi Data Panel --- Common Effect Model (CEM / Pooled OLS)}
\begin{verbatim}
==============================================================================
Dependent Variable: ROA
Method: Panel Least Squares
Sample: 2021Q1 2025Q4
Periods included: 20
Cross-sections included: 4
Total panel (balanced) observations: 80
==============================================================================
Variable             Coefficient   Std. Error   t-Statistic      Prob.
==============================================================================
C                       1.215420     0.427965      2.840000     0.0058
GREEN_FINANCING         0.051240     0.015004      3.415089     0.0010
NPL_GROSS              -0.421050     0.086814     -4.850024     0.0000
CAR                     0.062140     0.019727      3.150000     0.0023
==============================================================================
R-squared               0.612450     Mean dependent var         3.182250
Adjusted R-squared      0.597148     S.D. dependent var         0.541280
S.E. of regression      0.343560     Akaike info criterion      0.754120
Sum squared resid       8.970840     Schwarz criterion          0.873220
Log likelihood         -26.16480     Hannan-Quinn criter.       0.801890
F-statistic            40.020540     Durbin-Watson stat         1.420150
Prob(F-statistic)       0.000000
==============================================================================
\end{verbatim}

\vspace{0.4cm}
\noindent \textbf{Output 2.2: Estimasi Regresi Data Panel --- Fixed Effect Model (FEM / LSDV)}
\begin{verbatim}
==============================================================================
Dependent Variable: ROA
Method: Panel Least Squares
Sample: 2021Q1 2025Q4
Periods included: 20
Cross-sections included: 4
Total panel (balanced) observations: 80
==============================================================================
Variable             Coefficient   Std. Error   t-Statistic      Prob.
==============================================================================
C                       1.485210     0.312019      4.760000     0.0000
GREEN_FINANCING         0.042180     0.011048      3.817886     0.0003
NPL_GROSS              -0.385420     0.076080     -5.065983     0.0000
CAR                     0.074150     0.016032      4.625125     0.0000
------------------------------------------------------------------------------
Effects Specification: Cross-section fixed (dummy variables)
  _BBRI--C              0.181950
  _BMRI--C              0.044820
  _BBCA--C              0.215140
  _BBNI--C             -0.441910
==============================================================================
R-squared               0.725140     Mean dependent var         3.182250
Adjusted R-squared      0.702580     S.D. dependent var         0.541280
S.E. of regression      0.295160     Akaike info criterion      0.478950
Sum squared resid       6.360120     Schwarz criterion          0.687380
Log likelihood         -12.15800     Hannan-Quinn criter.       0.562550
F-statistic            32.450120     Durbin-Watson stat         1.942180
Prob(F-statistic)       0.000000
==============================================================================
\end{verbatim}

\vspace{0.4cm}
\noindent \textbf{Output 2.3: Estimasi Regresi Data Panel --- Random Effect Model (REM / EGLS)}
\begin{verbatim}
==============================================================================
Dependent Variable: ROA
Method: Panel EGLS (Cross-section random effects)
Sample: 2021Q1 2025Q4
Periods included: 20
Cross-sections included: 4
Total panel (balanced) observations: 80
==============================================================================
Variable             Coefficient   Std. Error   t-Statistic      Prob.
==============================================================================
C                       1.320450     0.387229      3.410000     0.0010
GREEN_FINANCING         0.046120     0.012740      3.620094     0.0005
NPL_GROSS              -0.398150     0.080925     -4.920000     0.0000
CAR                     0.068950     0.017725      3.890000     0.0002
==============================================================================
Effects Specification: Cross-section random (Swamy-Arora)
  Cross-section random S.D. / Rho : 0.185420 / 0.2831
  Idiosyncratic random S.D. / Rho : 0.295160 / 0.7169
==============================================================================
R-squared               0.647820     Mean dependent var         1.654210
Adjusted R-squared      0.633910     S.D. dependent var         0.482150
S.E. of regression      0.291750     Sum squared resid          6.469120
F-statistic            46.550140     Durbin-Watson stat         1.610450
Prob(F-statistic)       0.000000
==============================================================================
\end{verbatim}

\vspace{0.4cm}
\noindent \textbf{Output 2.4: Uji Chow (Redundant Fixed Effects Tests)}
\begin{verbatim}
==============================================================================
Redundant Fixed Effects Tests
Equation: FEM_ROA
Test cross-section fixed effects
==============================================================================
Effects Test                    Statistic        d.f.          Prob.
==============================================================================
Cross-section F                 18.420150      (3, 73)        0.0000
Cross-section Chi-square        48.310240            3        0.0000
==============================================================================
\end{verbatim}

\vspace{0.4cm}
\noindent \textbf{Output 2.5: Uji Hausman (Correlated Random Effects - Hausman Test)}
\begin{verbatim}
==============================================================================
Correlated Random Effects - Hausman Test
Equation: REM_ROA
Test cross-section random effects
==============================================================================
Test Summary                 Chi-Sq. Statistic    Chi-Sq. d.f.         Prob.
==============================================================================
Cross-section random                 14.250180               3        0.0026
==============================================================================
\end{verbatim}

\vspace{0.4cm}
\noindent \textbf{Output 2.6: Uji Normalitas Residual (Jarque-Bera Histogram)}
\begin{verbatim}
==============================================================================
Series: Standardized Residuals
Sample: 2021Q1 2025Q4
Observations: 80
==============================================================================
Mean                   -1.39e-16     Skewness               0.114250
Median                 -0.012540     Kurtosis               2.385410
Maximum                 0.684210     Jarque-Bera            1.842150
Minimum                -0.642150     Probability            0.398090
Std. Dev.               0.283740
==============================================================================
\end{verbatim}

\vspace{0.4cm}
\noindent \textbf{Output 2.7: Uji Multikolinearitas (Variance Inflation Factors)}
\begin{verbatim}
==============================================================================
Variance Inflation Factors
Sample: 2021Q1 2025Q4
Included observations: 80
==============================================================================
                       Coefficient   Uncentered     Centered
Variable                 Variance        VIF          VIF
==============================================================================
C                        0.097356     84.125012          NA
GREEN_FINANCING          0.000122     42.150420     1.284150
NPL_GROSS                0.005788     19.341250     1.421800
CAR                      0.000257     65.412850     1.352410
==============================================================================
\end{verbatim}

\vspace{0.4cm}
\noindent \textbf{Output 2.8: Uji Heteroskedastisitas (Uji Glejser)}
\begin{verbatim}
==============================================================================
Dependent Variable: ABS(RESID)
Method: Panel Least Squares (FEM)
Included observations: 80
==============================================================================
Variable             Coefficient   Std. Error   t-Statistic      Prob.
==============================================================================
C                       0.154210     0.185420      0.831679     0.4083
GREEN_FINANCING        -0.006540     0.006420     -1.018691     0.3117
NPL_GROSS               0.052140     0.044200      1.179638     0.2419
CAR                    -0.012450     0.009320     -1.335836     0.1858
==============================================================================
R-squared               0.058420     F-statistic                0.754120
Adjusted R-squared      0.019540     Prob(F-statistic)          0.523140
==============================================================================
\end{verbatim}

\vspace{0.8cm}
\noindent \textbf{Lampiran 3: Daftar Riwayat Hidup Penulis (\emph{Curriculum Vitae})}

\vspace{0.4cm}
\noindent
\begin{tabular}{lll}
  \textbf{Nama Lengkap} & : & \nahamahasiswa \\
  \textbf{Nomor Induk Mahasiswa (NIM)} & : & \nim \\
  \textbf{Tempat, Tanggal Lahir} & : & Jakarta, \underline{\hspace{4cm}} \\
  \textbf{Jenis Kelamin} & : & \underline{\hspace{4cm}} \\
  \textbf{Agama} & : & \underline{\hspace{4cm}} \\
  \textbf{Alamat Asal} & : & \underline{\hspace{7cm}} \\
  \textbf{Alamat Domisili} & : & \underline{\hspace{7cm}} \\
  \textbf{Nomor Telepon / WhatsApp} & : & \underline{\hspace{5cm}} \\
  \textbf{Alamat Surat Elektronik (E-mail)} & : & \underline{\hspace{7cm}} \\
  \textbf{Program Studi} & : & Strata 1 (S1) Manajemen \\
  \textbf{Konsentrasi Bidang Minat} & : & Manajemen Keuangan (\emph{Financial Management}) \\
  \textbf{Fakultas / Universitas} & : & Fakultas Ekonomi dan Bisnis, Universitas Kristen Krida Wacana \\
\end{tabular}

\vspace{0.5cm}
\noindent \textbf{Riwayat Pendidikan Formal:}

\vspace{0.2cm}
\noindent
\begin{tabular}{|c|l|l|c|}
  \hline
  \textbf{No} & \textbf{Jenjang Pendidikan} & \textbf{Nama Institusi / Sekolah} & \textbf{Tahun Kelulusan} \\
  \hline
  1 & Sekolah Dasar (SD) & \underline{\hspace{6cm}} & \underline{\hspace{2cm}} \\
  \hline
  2 & Sekolah Menengah Pertama (SMP) & \underline{\hspace{6cm}} & \underline{\hspace{2cm}} \\
  \hline
  3 & Sekolah Menengah Atas (SMA/SMK) & \underline{\hspace{6cm}} & \underline{\hspace{2cm}} \\
  \hline
  4 & Sarjana Manajemen (S1) & Universitas Kristen Krida Wacana (UKRIDA) & 2022--2026 \\
  \hline
\end{tabular}

\vspace{0.5cm}
\noindent \textbf{Pengalaman Organisasi dan Pelatihan:}

\vspace{0.2cm}
\noindent
\begin{enumerate}
  \item \underline{\hspace{12cm}}
  \item \underline{\hspace{12cm}}
  \item \underline{\hspace{12cm}}
\end{enumerate}

\vspace{0.8cm}
\noindent \textbf{Lampiran 4: Formulasi dan Standar Rasio Keuangan Perbankan (SE BI / OJK)}

\vspace{0.3cm}
\begin{table}[H]
  \small
  \begin{tabularx}{\textwidth}{|>{\raggedright\arraybackslash}p{2.6cm}|>{\centering\arraybackslash}p{4.2cm}|>{\centering\arraybackslash}p{2.2cm}|>{\raggedright\arraybackslash}X|}
    \hline
    \textbf{Rasio Keuangan} & \textbf{Formulasi Matematis} & \textbf{Standar Kesehatan OJK} & \textbf{Keterangan Analitis} \\
    \hline
    \textbf{Return on Assets (ROA)} & $\left(\dfrac{\text{Laba Sebelum Pajak}}{\text{Total Aset Rata-rata}}\right) \times 100\%$ & $\ge 1{,}50\%$ (Peringkat 1) & Mengukur efisiensi pemanfaatan aset operasional. \\
    \hline
    \textbf{Green Financing Ratio} & $\left(\dfrac{\text{Kredit KUBL POJK 51}}{\text{Total Portofolio Kredit}}\right) \times 100\%$ & Target RAKB Bank (Progresif) & Mengukur intensitas komitmen pembiayaan hijau. \\
    \hline
    \textbf{Non-Performing Loan Gross} & $\left(\dfrac{\text{Kredit Kol 3 + 4 + 5}}{\text{Total Kredit Disalurkan}}\right) \times 100\%$ & $\le 5{,}00\%$ (Batas Maksimum) & Mengukur risiko gagal bayar debitur kredit. \\
    \hline
    \textbf{Capital Adequacy Ratio} & $\left(\dfrac{\text{Total Modal (Tier 1+2)}}{\text{Total ATMR}}\right) \times 100\%$ & $\ge 8\% - 14\% + \text{Buffers}$ & Mengukur kecukupan modal penyerap risiko. \\
    \hline
  \end{tabularx}
\end{table}

\vspace{0.8cm}
\noindent \textbf{Lampiran 5: Klasifikasi 11 Sektor KUBL Berdasarkan POJK No. 51/POJK.03/2017}

\vspace{0.3cm}
\begin{table}[H]
  \small
  \begin{tabularx}{\textwidth}{|c|>{\raggedright\arraybackslash}p{3.8cm}|X|}
    \hline
    \textbf{No} & \textbf{Kategori KUBL} & \textbf{Contoh Kegiatan Usaha yang Memenuhi Syarat} \\
    \hline
    1 & Energi Terbarukan & PLTS, PLTA, PLTB, Panas Bumi, Biomassa, Biogas. \\
    2 & Efisiensi Energi & Retrofit sistem HVAC gedung, motor listrik efisiensi tinggi, smart grid. \\
    3 & Pengelolaan SDA Hayati \& Lahan & Pertanian organik, sertifikasi kelapa sawit RSPO/ISPO, kehutanan FSC. \\
    4 & Transportasi Ramah Lingkungan & Kendaraan bermotor listrik (KBLBB), LRT, MRT, bus listrik massal. \\
    5 & Pengelolaan Air \& Air Limbah & Instalasi Pengolahan Air Limbah (IPAL) industri, desalinasi ramah lingkungan. \\
    6 & Pengendalian \& Pencegahan Polusi & Pengolahan limbah B3, penangkap karbon, daur ulang sampah terpadu. \\
    7 & Bangunan Berwawasan Lingkungan & Gedung bersertifikat Green Building Council Indonesia (GBCI / Greenship). \\
    8 & Produk Ramah Lingkungan & Kemasan mudah terurai (*biodegradable*), produk bersertifikat ekolabel. \\
    9 & Konservasi Keanekaragaman Hayati & Rehabilitasi terumbu karang, suaka margasatwa, ekowisata lestari. \\
    10 & Adaptasi Perubahan Iklim & Sistem peringatan dini banjir, tanggul laut tahan abrasi, benih tahan kekeringan. \\
    11 & Kegiatan Hijau Lainnya & Pembiayaan rantai pasok hijau (*green supply chain financing*). \\
    \hline
  \end{tabularx}
\end{table}

\vspace{0.5cm}
\noindent
Demikian seluruh naskah tugas akhir skripsi ini disusun dan dipertanggungjawabkan secara ilmiah dan akademis.

\vspace{0.8cm}
\begin{flushright}
  Jakarta, 19 Agustus 2026\\[0.3cm]
  \textbf{Penulis}\\[1.8cm]
  \textbf{\nahamahasiswa}\\
  NIM: \nim
\end{flushright}

\end{document}
